% ----------------------------------------------------------------------
%  1.6. Кристаллография пластической деформации
% ----------------------------------------------------------------------
\begin{orientir}
Как кристалл меняет форму, не разрушаясь? В~1.6.1 разберём два способа
пластической деформации~— \vazhno{скольжение и двойникование}, закон
Шмида и то, какие системы скольжения работают в ГЦК, ОЦК и ГПУ-металлах.
Скольжение поворачивает решётку, поэтому деформированный поликристалл
приобретает \vazhno{текстуру}; в~1.6.2~— как её описывают (полюсные
фигуры и функция распределения ориентаций).
\end{orientir}

Под действием механического напряжения кристалл деформируется. При малых
напряжениях деформация \term{упругая}: после разгрузки форма полностью
восстанавливается, а деформация и напряжение связаны законом Гука
(п.~\ref{ssec:elastic-crystals}). Выше \term{предела текучести} появляется
\term{остаточная} (пластическая) деформация: после разгрузки кристалл
сохраняет изменённую форму. Практически предел упругости (последнее
напряжение, при котором выполняется закон Гука) неотличим от предела
текучести.

\subsection{Геометрия пластической деформации монокристаллов}\label{ssec:slip-geometry}

Монокристалл деформируется пластически двумя способами
(рис.~\ref{fig:slip-twin}).

\begin{opredelenie}
\term{Скольжение}~— сдвиг одной части кристалла относительно другой по
определённой кристаллографической плоскости (\term{плоскости скольжения})
в определённом направлении (\term{направлении скольжения}) на величину,
\emph{кратную} периоду решётки в этом направлении. Пару «плоскость +
направление» называют \term{системой скольжения} и записывают как
$\plfam{hkl}\nfam{uvw}$.

\term{Механическое двойникование}~— сдвиг, в результате которого часть
кристалла переходит в положение, зеркально-симметричное исходному (или
повёрнутое относительно него на $180^\circ$). Атомы при этом смещаются на
\emph{доли} периода, пропорционально расстоянию от плоскости
двойникования.
\end{opredelenie}

\begin{figure}[ht]
\centering
\begin{tikzpicture}[scale=0.38]
  % (а) исходный кристалл
  \begin{scope}
    \foreach \j in {0,...,7}{\foreach \i in {0,...,6}{
      \fill[atomA!80!black] ({\i+0.35*\j},{0.8*\j}) circle (0.17);}}
    \draw[atomLine, line width=0.6pt] (-0.575,-0.4) -- (6.225,-0.4) -- (9.025,6.0) -- (2.225,6.0) -- cycle;
    \node[font=\small] at (4.2,-1.5) {\textit{а}};
  \end{scope}
  % (б) скольжение: верхняя половина сдвинута на два периода
  \begin{scope}[shift={(11,0)}]
    \foreach \j in {0,...,3}{\foreach \i in {0,...,6}{
      \fill[atomA!80!black] ({\i+0.35*\j},{0.8*\j}) circle (0.17);}}
    \foreach \j in {4,...,7}{\foreach \i in {0,...,6}{
      \fill[atomA!80!black] ({\i+0.35*\j+2},{0.8*\j}) circle (0.17);}}
    \draw[atomLine, line width=0.6pt] (-0.575,-0.4) -- (6.225,-0.4) -- (7.625,2.8) -- (9.625,2.8)
      -- (11.025,6.0) -- (4.225,6.0) -- (2.825,2.8) -- (0.825,2.8) -- cycle;
    \draw[fmRed, dashed, line width=0.7pt] (-0.6,2.8) -- (11.4,2.8);
    \draw[arrB] (6.2,6.7) -- (8.2,6.7);
    \node[font=\small] at (5.2,-1.5) {\textit{б}};
  \end{scope}
  % (в) двойникование: слои 1 и 5 — границы двойниковой прослойки
  \begin{scope}[shift={(25,0)}]
    \fill[fmAmber!20] (-1.6,0.8) rectangle (7.0,4.0);
    \foreach \j in {0,...,7}{\foreach \i in {0,...,6}{
      \pgfmathsetmacro{\sh}{\j<=1 ? 0.35*\j : (\j<=5 ? 0.35-0.35*(\j-1) : 0.35-1.4+0.35*(\j-5))}
      \pgfmathsetmacro{\col}{(\j>=2 && \j<=4) ? 1 : 0}
      \ifdim\col pt>0.5pt
        \fill[atomB!90!black] ({\i+\sh},{0.8*\j}) circle (0.17);
      \else
        \fill[atomA!80!black] ({\i+\sh},{0.8*\j}) circle (0.17);
      \fi}}
    \draw[fmAmber!80!black, dashed, line width=0.7pt] (-1.6,0.8) -- (7.0,0.8) (-1.6,4.0) -- (7.0,4.0);
    \node[font=\small] at (2.7,-1.5) {\textit{в}};
  \end{scope}
\end{tikzpicture}
\caption{Два механизма пластической деформации: \textit{а}~— исходный
кристалл; \textit{б}~— скольжение по плоскости, показанной штриховой линией:
верхняя часть сдвинута на целое число периодов, решётка не изменилась,
на боковых гранях появились ступеньки; \textit{в}~— двойникование: в
прослойке (оранжевые атомы) каждый слой сдвинут относительно
предыдущего на долю периода, и решётка зеркально отражена относительно
границ прослойки (штриховые линии)}
\label{fig:slip-twin}
\end{figure}

Главный механизм пластической деформации при обычных и низких
температурах~— скольжение. Объём при скольжении не меняется, а решётка
после сдвига остаётся той же (сдвиг кратен трансляции), поэтому скольжение
называют также \term{трансляционным}. На атомном уровне скольжение~— это
движение дислокаций (гл.~2, п.~2.2.2).

\subsubsection{Закон Шмида}

Растянем цилиндрический монокристалл напряжением $\sigma$ вдоль оси $F$
(рис.~\ref{fig:schmid}). Пусть нормаль $\vb n$ к плоскости скольжения
образует с осью угол $\chi$, а направление скольжения $\vb m$~— угол
$\lambda$. Сила $\sigma A$ действует на площадку $A/\cos\chi$, её проекция
на направление $\vb m$ равна $\sigma A\cos\lambda$. Поэтому касательное
(«скалывающее») напряжение в системе скольжения
\begin{keyeq}\label{eq:schmid}
\tau=\sigma\cos\chi\cos\lambda .
\end{keyeq}
Множитель $\cos\chi\cos\lambda$ называют \term{фактором Шмида}
(ориентационным множителем).

\begin{glavnoe}[Закон Шмида (закон критического скалывающего напряжения)]
Скольжение начинается, когда касательное напряжение хотя бы в одной из
систем скольжения достигает критического значения $\tau_\text{кр}$~—
постоянной материала, \emph{не зависящей} от ориентации кристалла.
Первой включается система с наибольшим фактором Шмида.
\end{glavnoe}

\begin{figure}[ht]
\centering
\begin{tikzpicture}[scale=1.15]
  \fill[fmBlueBg] (-0.9,-2) -- (-0.9,2) arc[start angle=180, end angle=360, x radius=0.9, y radius=0.25]
    -- (0.9,-2) arc[start angle=0, end angle=-180, x radius=0.9, y radius=0.25];
  \draw[atomLine] (-0.9,-2) -- (-0.9,2) (0.9,-2) -- (0.9,2);
  \draw[atomLine] (0,2) ellipse (0.9 and 0.25);
  \draw[atomLine] (-0.9,-2) arc[start angle=180, end angle=360, x radius=0.9, y radius=0.25];
  \draw[atomLine, dashed] (0.9,-2) arc[start angle=0, end angle=180, x radius=0.9, y radius=0.25];
  \fill[fmRed!25, rotate=30] (0,0) ellipse (1.04 and 0.25);
  \draw[fmRed!80!black, rotate=30] (0,0) ellipse (1.04 and 0.25);
  \node[font=\footnotesize, fmRed!80!black] at (-1.35,-0.75) {$S$};
  \draw[arrB, fmBlue] (0,2.3) -- (0,3.1) node[above, font=\small]{$F$};
  \draw[arrB, fmBlue] (0,-2.3) -- (0,-3.1);
  \draw[fmGray, dashed] (0,-2.0) -- (0,2.0);
  \draw[bvec, fmGreen!70!black] (0,0) -- (-0.75,1.30) node[above left=-2pt, font=\small]{$\vb n$};
  \draw[bvec] (0,0) -- (0.62,0.82) node[right, font=\small]{$\vb m$};
  \draw[fmGreen!70!black] (0,0.75) arc[start angle=90, end angle=120, radius=0.75];
  \node[font=\footnotesize, fmGreen!70!black] at (-0.22,0.95) {$\chi$};
  \draw[fmRed] (0,0.45) arc[start angle=90, end angle=53, radius=0.45];
  \node[font=\footnotesize, fmRed] at (0.17,0.62) {$\lambda$};
  \node[font=\small, fmBlue] at (0.35,-2.75) {$\sigma$};
\end{tikzpicture}
\caption{К закону Шмида: ось растяжения $F$, плоскость скольжения $S$
с нормалью $\vb n$ и направление скольжения $\vb m$}
\label{fig:schmid}
\end{figure}

Фактор Шмида не превышает $\frac12$: максимум достигается, когда $\vb n$,
$\vb m$ и $F$ лежат в одной плоскости и $\chi=\lambda=45^\circ$. Если ось
лежит в плоскости скольжения ($\chi=90^\circ$) или перпендикулярна
направлению скольжения ($\lambda=90^\circ$), скольжение в этой системе
невозможно при любом напряжении~— кристалл в такой ориентации кажется
очень прочным.

\begin{masshtab}[где происходит скольжение]{-9.75/-9.35/{$b$ и $d_{111}$~— доли нм}, -8/-6.6/{ступеньки 10--300 нм}, -6.3/-5/{полосы через 1--10 мкм}, -5/-4/{зерно}, -3.3/-1.7/{образец}}
Цилиндр на рис.~\ref{fig:schmid}~— это \emph{образец} (миллиметры--сантиметры),
а плоскость $S$~— не одна атомная плоскость, а \emph{ориентация}: все
параллельные плоскости $\plfam{111}$ кристалла. В меди
($d_{111}=0{,}21$~нм) на 1~мм вдоль нормали их около $5\cdot10^6$, и
фактор Шмида у всех одинаков. Поэтому он характеризует ориентацию
кристалла (зерна) как целого, а не какую-то одну плоскость.

Сдвиг же идёт не по всем этим плоскостям, а по немногим, где работают
источники дислокаций (рис.~\ref{fig:slip-scales}). Каждая дислокация
сдвигает кристалл на $b=0{,}26$~нм; сотни дислокаций, вышедших по одной
плоскости, дают на поверхности ступеньку высотой десятки--сотни
нанометров~— \term{линию скольжения}. Линии собираются в \term{полосы
скольжения}, отстоящие друг от друга на микрометры; их видно в
оптический микроскоп. В поликристалле роль «цилиндра» играет отдельное
зерно (10--100~мкм) со своим фактором Шмида, а предел текучести
поликристалла $\sigma_\text{т}\approx M\tau_\text{кр}$, где $M\approx3{,}1$
для ГЦК (фактор Тейлора, усреднение по ориентациям зёрен) вместо
$1/0{,}5=2$ у наилучшим образом ориентированного монокристалла.
\end{masshtab}

\begin{figure}[ht]
\centering
\begin{tikzpicture}[font=\footnotesize]
  % ---------- (а) образец: миллиметры
  \begin{scope}
    \fill[fmBlueBg] (-0.7,-1.5) rectangle (0.7,1.5);
    \begin{scope}
      \clip (-0.7,-1.5) rectangle (0.7,1.5);
      \foreach \c in {-2.4,-2.25,...,2.4}{\draw[fmBlue!45, line width=0.3pt] (-0.7,{\c-0.42}) -- (0.7,{\c+0.42});}
    \end{scope}
    \draw[atomLine] (-0.7,-1.5) -- (-0.7,1.5) (0.7,-1.5) -- (0.7,1.5);
    \draw[atomLine, fill=fmBlueBg] (0,1.5) ellipse (0.7 and 0.16);
    \draw[atomLine] (-0.7,-1.5) arc[start angle=180, end angle=360, x radius=0.7, y radius=0.16];
    \draw[arrB, fmBlue] (0,1.75) -- (0,2.2) node[above, inner sep=1pt]{$F$};
    \draw[arrB, fmBlue] (0,-1.75) -- (0,-2.2);
    \draw[fmAmber!80!black, line width=0.8pt] (-0.12,-0.12) rectangle (0.12,0.12);
    \draw[atomLine, line width=0.8pt] (-0.7,-2.45) -- (-0.23,-2.45) node[right, inner sep=2pt]{1~мм};
    \node[font=\small] at (0,-2.85) {\textit{а}\ образец};
  \end{scope}
  % ---------- (б) поверхность: микрометры
  \begin{scope}[shift={(3.0,0)}]
    \fill[fmBlueBg] (-1.4,-1.4) rectangle (1.4,1.4);
    \begin{scope}
      \clip (-1.4,-1.4) rectangle (1.4,1.4);
      \foreach \c in {-2.1,-1.35,...,2.2}{
        \foreach \d in {0,0.06,0.12,0.18}{\draw[fmBlue!80!black, line width=0.35pt] (-1.4,{\c+\d-0.84}) -- (1.4,{\c+\d+0.84});}}
    \end{scope}
    \draw[atomLine] (-1.4,-1.4) rectangle (1.4,1.4);
    \draw[fmAmber!80!black, line width=0.8pt] (0.05,0.25) rectangle (0.45,0.55);
    \draw[atomLine, line width=0.8pt] (-1.4,-1.75) -- (-0.65,-1.75) node[right, inner sep=2pt]{5~мкм};
    \node[font=\small] at (0,-2.85) {\textit{б}\ полосы};
  \end{scope}
  \draw[atomLine!50, densely dotted] (0.12,0.12) -- (1.6,1.4) (0.12,-0.12) -- (1.6,-1.4);
  % ---------- (в) профиль ступенек: сотни нанометров
  \begin{scope}[shift={(6.3,0)}]
    \fill[fmBlueBg] (-1.5,-1.4) -- (-1.5,-0.55) -- (-0.75,-0.55) -- (-0.45,-0.05) -- (0.05,-0.05) -- (0.2,0.2)
      -- (0.65,0.2) -- (1.0,0.8) -- (1.5,0.8) -- (1.5,-1.4) -- cycle;
    \draw[atomLine, line width=0.9pt] (-1.5,-0.55) -- (-0.75,-0.55) -- (-0.45,-0.05) -- (0.05,-0.05) -- (0.2,0.2)
      -- (0.65,0.2) -- (1.0,0.8) -- (1.5,0.8);
    \foreach \x/\y in {-0.75/-0.55, 0.05/-0.05, 0.65/0.2}{
      \draw[fmRed!70, dashed, line width=0.5pt] (\x,\y) -- ++(-0.9,-1.5);}
    \draw[dimline, fmRed] (1.15,0.2) -- node[right, inner sep=1pt, fmRed]{$h=nb$} (1.15,0.8);
    \draw[fmAmber!80!black, line width=0.8pt] (-0.86,-0.75) rectangle (-0.56,-0.45);
    \draw[atomLine, line width=0.8pt] (-1.5,-1.75) -- (-0.9,-1.75) node[right, inner sep=2pt]{100~нм};
    \node[font=\small] at (0,-2.85) {\textit{в}\ ступеньки};
  \end{scope}
  \draw[atomLine!50, densely dotted] (3.45,0.55) -- (4.8,1.4) (3.45,0.25) -- (4.8,-1.4);
  % ---------- (г) атомы: нанометры
  \begin{scope}[shift={(9.75,0)}]
    \begin{scope}
      \clip (-1.45,-1.4) rectangle (1.45,1.4);
      \begin{scope}[scale=0.27]\EdgeDislLattice{6}{5}{5}\end{scope}
    \end{scope}
    \draw[atomLine] (-1.45,-1.4) rectangle (1.45,1.4);
    \draw[fmRed, dashed, line width=0.7pt] (-1.45,0) -- (1.45,0);
    \dislT[fmRed]{(0,0)}{0}
    \draw[atomLine, line width=0.8pt] (-1.45,-1.75) -- (-0.4,-1.75) node[right, inner sep=2pt]{1~нм};
    \node[font=\small] at (0,-2.85) {\textit{г}\ дислокация};
  \end{scope}
  \draw[atomLine!50, densely dotted] (5.74,-0.45) -- (8.3,1.4) (5.74,-0.75) -- (8.3,-1.4);
\end{tikzpicture}
\caption{Масштабы скольжения: \textit{а}~— растягиваемый образец с
линиями скольжения на поверхности (миллиметры); \textit{б}~— полосы
скольжения на поверхности зерна или монокристалла: пучки близких линий,
отстоящие друг от друга на микрометры; \textit{в}~— разрез поверхности:
каждая линия скольжения~— ступенька высотой $h=nb$ (десятки--сотни
нанометров), образованная выходом $n\sim10^2$--$10^3$ дислокаций по одной
плоскости (пунктир); \textit{г}~— атомный масштаб: одна краевая
дислокация в плоскости скольжения сдвигает кристалл на $b\approx0{,}26$~нм.
Каждая следующая панель~— увеличение выделенного участка предыдущей}
\label{fig:slip-scales}
\end{figure}

Если кристалл находится под действием произвольного тензора напряжений
$\sigma$, касательное напряжение в системе с единичной нормалью $\vb h$ и
единичным направлением $\vb m$ вычисляют как
\begin{equation}\label{eq:schmid-tensor}
\tau_{\vb h\vb m}=\tilde h\,\sigma\, m=\sum_{i,j}h_i\sigma_{ij}m_j ,
\end{equation}
где $\tilde h$~— строка направляющих косинусов нормали, $m$~— столбец
направляющих косинусов направления скольжения. Для одноосного растяжения
$\sigma_{ij}=\sigma F_iF_j$ формула~\eqref{eq:schmid-tensor} переходит
в~\eqref{eq:schmid}.

\subsubsection{Тензор пластической дисторсии}

При скольжении на величину $\alpha$ (сдвиг на единицу расстояния от
плоскости скольжения) точка $P$ с радиус-вектором $\vb r$ смещается
вдоль $\vb m$ на величину, пропорциональную её расстоянию $\vb r\cdot\vb h$
от плоскости скольжения, проходящей через начало координат
(рис.~\ref{fig:shear-vector}):
\begin{equation}\label{eq:slip-displacement}
\vb u=\vb r'-\vb r=\alpha\,(\vb r\cdot\vb h)\,\vb m .
\end{equation}

\begin{figure}[ht]
\centering
\begin{tikzpicture}[scale=1.1]
  \fill[fmBlueBg] (-0.3,0) rectangle (4.3,2.2);
  \fill[fmRed!10] (-0.3,0) -- (4.3,0) -- (5.2,2.2) -- (0.6,2.2) -- cycle;
  \draw[atomLine, dashed] (-0.3,0) rectangle (4.3,2.2);
  \draw[fmRed!80!black, line width=0.8pt] (-0.3,0) -- (4.3,0) -- (5.2,2.2) -- (0.6,2.2) -- cycle;
  \draw[fmGray, line width=1pt] (-0.6,0) -- (5.4,0) node[right, font=\footnotesize]{плоскость скольжения};
  \fill (0,0) circle (1.5pt) node[below left, font=\small]{$O$};
  \fill (2.2,1.5) circle (1.5pt) node[above left, font=\small]{$P$};
  \fill (2.81,1.5) circle (1.5pt) node[above right, font=\small]{$P'$};
  \draw[arr] (0,0) -- (2.2,1.5) node[midway, above left=-2pt, font=\small]{$\vb r$};
  \draw[arr, fmRed] (2.2,1.5) -- (2.81,1.5);
  \draw[fmGray, densely dotted] (2.2,1.5) -- (2.2,0);
  \node[font=\footnotesize, anchor=west] at (2.25,0.6) {$\vb r\cdot\vb h$};
  \draw[bvec, fmGreen!70!black] (-0.3,0) -- (-0.3,0.9) node[left, font=\small]{$\vb h$};
  \draw[bvec] (0.6,-0.35) -- (1.6,-0.35) node[right, font=\small]{$\vb m$};
\end{tikzpicture}
\caption{Смещение точки при скольжении: $\overrightarrow{PP'}=\alpha(\vb r\cdot\vb h)\vb m$
(пунктир~— форма кристалла до сдвига)}
\label{fig:shear-vector}
\end{figure}

Компоненты \term{тензора пластической дисторсии}
$u^0_{ij}=\partial u_i/\partial x_j$ (сравните с~\eqref{eq:strain}): так как
$u_i=\alpha\,m_i\sum_kx_kh_k$, то $u^0_{ij}=\alpha\,m_ih_j$, или в матричной форме
\begin{keyeq}\label{eq:plastic-distortion}
u^0=\alpha\, m\,\tilde h=\alpha
\begin{pmatrix}
m_1h_1 & m_1h_2 & m_1h_3\\
m_2h_1 & m_2h_2 & m_2h_3\\
m_3h_1 & m_3h_2 & m_3h_3
\end{pmatrix}\!.
\end{keyeq}
Например, $u^0_{11}=\alpha m_1h_1$, $u^0_{23}=\partial u_2/\partial x_3=\alpha m_2h_3$.

Свойства этого тензора:
\begin{itemize}
\item След $u^0_{11}+u^0_{22}+u^0_{33}=\alpha(\vb m\cdot\vb h)=0$, так как
$\vb m\perp\vb h$: \vazhno{скольжение не меняет объём}.
\item Тензор несимметричен. Его симметричная часть
$\varepsilon_{ij}=\frac\alpha2(m_ih_j+m_jh_i)$~— \emph{чистая пластическая
деформация}, антисимметричная
$\omega_{ij}=\frac\alpha2(m_ih_j-m_jh_i)$~— \emph{поворот решётки}.
Именно этот поворот приводит к образованию текстуры (п.~\ref{ssec:texture}).
\item Отрезок материала вдоль оси $F$ после скольжения переходит, по
\eqref{eq:slip-displacement}, в $F'=F+\alpha(\vb h\cdot F)\,\vb m$~— он
\emph{поворачивается к направлению скольжения} $\vb m$, оставаясь в
плоскости, содержащей $F$ и $\vb m$.
\end{itemize}

\begin{smysl}
Тензор деформаций с нулевым следом имеет $6-1=5$ независимых компонент.
Значит, чтобы зерно поликристалла могло принять \emph{любую} форму,
навязанную соседями (без нарушения сплошности), в нём должны работать
не менее \vazhno{пяти независимых систем скольжения} (критерий Мизеса).
ГЦК- и ОЦК-металлы этому условию удовлетворяют и пластичны в
поликристаллическом состоянии; ГПУ-металлы с одним базисным скольжением~—
нет.
\end{smysl}

\subsubsection{Скольжение в ГЦК-металлах}

Плоскостями и направлениями скольжения, как правило, служат плоскости и
направления \emph{плотнейшей упаковки}: у них наибольшее межплоскостное
расстояние и наименьший вектор трансляции (а значит, и вектор Бюргерса
дислокаций). В ГЦК-металлах (Al, Cu, Ni, Ag, Au, $\gamma$-Fe) система
скольжения~— $\plfam{111}\nfam{110}$: четыре плоскости $\plfam{111}$, в
каждой три направления $\nfam{110}$, итого \vazhno{12 систем}
скольжения\ispr{В оригинале: «8 плоскостей скольжения с шестью
направлениями в каждой плоскости»~— это подсчёт с учётом знаков
нормалей и направлений; различных систем 12.}.

Какая из 12 систем включится первой, зависит от ориентации оси
растяжения. Благодаря кубической симметрии достаточно рассматривать оси,
лежащие в одном \term{стандартном стереографическом треугольнике}~—
здесь, как и в оригинале курса, это треугольник
$[100]$--$[110]$--$[111]$ (рис.~\ref{fig:slip-stereo}): любая другая
ориентация переводится в него операцией симметрии. В табл.~\ref{tab:schmid}
приведены факторы Шмида всех 12 систем для оси $F\parallel[531]$, лежащей
внутри треугольника.

\begin{table}[ht]
\centering
\caption{Факторы Шмида систем $\plfam{111}\nfam{110}$ ГЦК-кристалла при
растяжении вдоль $[531]$ (расчёт по~\eqref{eq:schmid})}
\label{tab:schmid}
\small
\begin{tabular}{@{}lcl@{\hspace{8mm}}lc@{}}
\toprule
Система & $\cos\chi\cos\lambda$ & Название & Система & $\cos\chi\cos\lambda$\\
\midrule
$(11\br1)[101]$ & \textbf{0,490} & первичная   & $(111)[1\br10]$   & 0,210\\
$(111)[10\br1]$ & 0,420 & критическая          & $(11\br1)[1\br10]$ & 0,163\\
$(11\br1)[011]$ & 0,327 &                      & $(1\br11)[10\br1]$ & 0,140\\
$(1\br11)[110]$ & 0,280 & сопряжённая          & $(1\br11)[011]$    & 0,140\\
$(111)[01\br1]$ & 0,210 &                      & $(1\br1\br1)[110]$ & 0,093\\
$(1\br1\br1)[101]$ & 0,070 & поперечная        & $(1\br1\br1)[01\br1]$ & 0,023\\
\bottomrule
\end{tabular}
\end{table}

\begin{figure}[t]
\centering
\tikzset{zone100/.style={fmBlue!70, line width=0.5pt},
  zone110/.style={fmBlue!35, line width=0.4pt},
  pole100/.style={rectangle, fill=atomLine, inner sep=0pt, minimum size=4pt},
  pole110/.style={ellipse, fill=atomLine, inner sep=0pt, minimum width=4.5pt, minimum height=3pt},
  pole111/.style={regular polygon, regular polygon sides=3, fill=atomLine, inner sep=0pt, minimum size=6pt},
  ring/.style={circle, draw=#1, line width=1pt, minimum size=8pt, inner sep=0pt},
  poleH/.style={circle, fill=fmGreen!70!black, inner sep=0pt, minimum size=3pt}}
\begin{tikzpicture}
  % создано tools/gen_plastic.py: системы скольжения ГЦК, проекция вдоль [100]
\begin{scope}
\fill[fmAmber!25] (0.0000,0.0000) -- (0.0242,0.0000) -- (0.0483,0.0000) -- (0.0725,0.0000) -- (0.0967,0.0000) -- (0.1209,0.0000) -- (0.1452,0.0000) -- (0.1694,0.0000) -- (0.1937,0.0000) -- (0.2181,0.0000) -- (0.2425,0.0000) -- (0.2669,0.0000) -- (0.2914,0.0000) -- (0.3160,0.0000) -- (0.3406,0.0000) -- (0.3653,0.0000) -- (0.3900,0.0000) -- (0.4149,0.0000) -- (0.4398,0.0000) -- (0.4648,0.0000) -- (0.4900,0.0000) -- (0.5152,0.0000) -- (0.5405,0.0000) -- (0.5660,0.0000) -- (0.5915,0.0000) -- (0.6172,0.0000) -- (0.6431,0.0000) -- (0.6691,0.0000) -- (0.6952,0.0000) -- (0.7214,0.0000) -- (0.7479,0.0000) -- (0.7745,0.0000) -- (0.8012,0.0000) -- (0.8282,0.0000) -- (0.8553,0.0000) -- (0.8827,0.0000) -- (0.9102,0.0000) -- (0.9379,0.0000) -- (0.9659,0.0000) -- (0.9941,0.0000) -- (0.9941,0.0000) -- (0.9940,0.0222) -- (0.9938,0.0444) -- (0.9935,0.0666) -- (0.9930,0.0888) -- (0.9923,0.1110) -- (0.9915,0.1332) -- (0.9906,0.1554) -- (0.9895,0.1776) -- (0.9882,0.1998) -- (0.9868,0.2221) -- (0.9853,0.2444) -- (0.9836,0.2666) -- (0.9818,0.2889) -- (0.9798,0.3112) -- (0.9777,0.3336) -- (0.9754,0.3559) -- (0.9730,0.3783) -- (0.9704,0.4007) -- (0.9676,0.4231) -- (0.9647,0.4455) -- (0.9617,0.4680) -- (0.9585,0.4905) -- (0.9551,0.5130) -- (0.9516,0.5356) -- (0.9479,0.5582) -- (0.9441,0.5808) -- (0.9400,0.6034) -- (0.9359,0.6261) -- (0.9315,0.6489) -- (0.9270,0.6716) -- (0.9223,0.6944) -- (0.9175,0.7173) -- (0.9124,0.7402) -- (0.9072,0.7631) -- (0.9018,0.7861) -- (0.8963,0.8091) -- (0.8905,0.8322) -- (0.8846,0.8553) -- (0.8785,0.8785) -- (0.8785,0.8785) -- (0.8523,0.8523) -- (0.8264,0.8264) -- (0.8008,0.8008) -- (0.7756,0.7756) -- (0.7506,0.7506) -- (0.7259,0.7259) -- (0.7014,0.7014) -- (0.6772,0.6772) -- (0.6532,0.6532) -- (0.6295,0.6295) -- (0.6059,0.6059) -- (0.5826,0.5826) -- (0.5594,0.5594) -- (0.5365,0.5365) -- (0.5137,0.5137) -- (0.4911,0.4911) -- (0.4687,0.4687) -- (0.4464,0.4464) -- (0.4242,0.4242) -- (0.4022,0.4022) -- (0.3803,0.3803) -- (0.3585,0.3585) -- (0.3369,0.3369) -- (0.3153,0.3153) -- (0.2939,0.2939) -- (0.2725,0.2725) -- (0.2512,0.2512) -- (0.2300,0.2300) -- (0.2089,0.2089) -- (0.1878,0.1878) -- (0.1668,0.1668) -- (0.1459,0.1459) -- (0.1249,0.1249) -- (0.1041,0.1041) -- (0.0832,0.0832) -- (0.0624,0.0624) -- (0.0416,0.0416) -- (0.0208,0.0208) -- (0.0000,0.0000) -- cycle;
\draw[zone100] (-0.0000,2.3768) -- (-0.0000,2.2953) -- (-0.0000,2.2164) -- (-0.0000,2.1400) -- (-0.0000,2.0659) -- (-0.0000,1.9941) -- (-0.0000,1.9243) -- (-0.0000,1.8564) -- (-0.0000,1.7903) -- (-0.0000,1.7260) -- (-0.0000,1.6632) -- (0.0000,1.6019) -- (0.0000,1.5421) -- (0.0000,1.4835) -- (0.0000,1.4263) -- (0.0000,1.3702) -- (0.0000,1.3152) -- (0.0000,1.2612) -- (0.0000,1.2082) -- (0.0000,1.1562) -- (0.0000,1.1050) -- (0.0000,1.0546) -- (0.0000,1.0050) -- (0.0000,0.9561) -- (0.0000,0.9079) -- (0.0000,0.8604) -- (0.0000,0.8134) -- (0.0000,0.7670) -- (0.0000,0.7210) -- (0.0000,0.6756) -- (0.0000,0.6306) -- (0.0000,0.5860) -- (0.0000,0.5418) -- (0.0000,0.4980) -- (0.0000,0.4544) -- (0.0000,0.4112) -- (0.0000,0.3682) -- (0.0000,0.3254) -- (0.0000,0.2829) -- (0.0000,0.2405) -- (0.0000,0.1982) -- (0.0000,0.1561) -- (0.0000,0.1141) -- (0.0000,0.0722) -- (0.0000,0.0302) -- (0.0000,-0.0116) -- (0.0000,-0.0535) -- (0.0000,-0.0955) -- (0.0000,-0.1375) -- (0.0000,-0.1795) -- (0.0000,-0.2217) -- (0.0000,-0.2640) -- (0.0000,-0.3065) -- (0.0000,-0.3492) -- (0.0000,-0.3921) -- (0.0000,-0.4352) -- (0.0000,-0.4352) -- (0.0000,-0.4786) -- (0.0000,-0.5223) -- (0.0000,-0.5664) -- (0.0000,-0.6108) -- (0.0000,-0.6556) -- (0.0000,-0.7008) -- (0.0000,-0.7465) -- (0.0000,-0.7927) -- (0.0000,-0.8394) -- (0.0000,-0.8867) -- (0.0000,-0.9347) -- (0.0000,-0.9832) -- (0.0000,-1.0325) -- (0.0000,-1.0825) -- (0.0000,-1.1333) -- (0.0000,-1.1850) -- (0.0000,-1.2376) -- (0.0000,-1.2911) -- (0.0000,-1.3456) -- (0.0000,-1.4012) -- (0.0000,-1.4580) -- (0.0000,-1.5159) -- (0.0000,-1.5752) -- (0.0000,-1.6358) -- (0.0000,-1.6979) -- (0.0000,-1.7616) -- (0.0000,-1.8269) -- (0.0000,-1.8939) -- (0.0000,-1.9628) -- (0.0000,-2.0338) -- (0.0000,-2.1068) -- (0.0000,-2.1821) -- (0.0000,-2.2599) -- (0.0000,-2.3403);
\draw[zone100] (-2.3570,0.0000) -- (-2.2761,0.0000) -- (-2.1978,0.0000) -- (-2.1220,0.0000) -- (-2.0485,0.0000) -- (-1.9771,0.0000) -- (-1.9078,0.0000) -- (-1.8404,0.0000) -- (-1.7747,0.0000) -- (-1.7108,0.0000) -- (-1.6484,0.0000) -- (-1.5874,0.0000) -- (-1.5279,0.0000) -- (-1.4697,0.0000) -- (-1.4127,0.0000) -- (-1.3569,0.0000) -- (-1.3021,0.0000) -- (-1.2484,0.0000) -- (-1.1956,0.0000) -- (-1.1438,0.0000) -- (-1.0928,0.0000) -- (-1.0426,0.0000) -- (-0.9932,0.0000) -- (-0.9445,0.0000) -- (-0.8965,0.0000) -- (-0.8490,0.0000) -- (-0.8022,0.0000) -- (-0.7559,0.0000) -- (-0.7101,0.0000) -- (-0.6648,0.0000) -- (-0.6199,-0.0000) -- (-0.5754,-0.0000) -- (-0.5313,-0.0000) -- (-0.4875,-0.0000) -- (-0.4440,-0.0000) -- (-0.4008,-0.0000) -- (-0.3579,-0.0000) -- (-0.3152,-0.0000) -- (-0.2727,-0.0000) -- (-0.2303,-0.0000) -- (-0.1881,-0.0000) -- (-0.1460,-0.0000) -- (-0.1040,-0.0000) -- (-0.0621,-0.0000) -- (-0.0202,-0.0000) -- (0.0217,-0.0000) -- (0.0636,-0.0000) -- (0.1055,-0.0000) -- (0.1476,-0.0000) -- (0.1896,-0.0000) -- (0.2319,-0.0000) -- (0.2742,-0.0000) -- (0.3167,-0.0000) -- (0.3595,-0.0000) -- (0.4024,-0.0000) -- (0.4456,-0.0000) -- (0.4891,-0.0000) -- (0.5329,-0.0000) -- (0.5770,-0.0000) -- (0.6215,-0.0000) -- (0.6664,-0.0000) -- (0.7117,-0.0000) -- (0.7576,-0.0000) -- (0.8039,-0.0000) -- (0.8507,-0.0000) -- (0.8982,-0.0000) -- (0.9463,-0.0000) -- (0.9950,-0.0000) -- (1.0445,-0.0000) -- (1.0947,-0.0000) -- (1.1457,-0.0000) -- (1.1975,-0.0000) -- (1.2503,-0.0000) -- (1.3041,-0.0000) -- (1.3589,-0.0000) -- (1.4147,0.0000) -- (1.4718,0.0000) -- (1.5300,0.0000) -- (1.5896,0.0000) -- (1.6506,0.0000) -- (1.7130,0.0000) -- (1.7771,0.0000) -- (1.8428,0.0000) -- (1.9103,0.0000) -- (1.9797,0.0000) -- (2.0511,0.0000) -- (2.1247,0.0000) -- (2.2006,0.0000) -- (2.2790,0.0000) -- (2.3600,0.0000);
\draw[zone100] (-1.6754,1.6754) -- (-1.6179,1.6179) -- (-1.5623,1.5623) -- (-1.5084,1.5084) -- (-1.4562,1.4562) -- (-1.4055,1.4055) -- (-1.3563,1.3563) -- (-1.3084,1.3084) -- (-1.2618,1.2618) -- (-1.2164,1.2164) -- (-1.1721,1.1721) -- (-1.1289,1.1289) -- (-1.0867,1.0867) -- (-1.0454,1.0454) -- (-1.0049,1.0049) -- (-0.9653,0.9653) -- (-0.9265,0.9265) -- (-0.8884,0.8884) -- (-0.8510,0.8510) -- (-0.8143,0.8143) -- (-0.7781,0.7781) -- (-0.7426,0.7426) -- (-0.7075,0.7075) -- (-0.6730,0.6730) -- (-0.6390,0.6390) -- (-0.6054,0.6054) -- (-0.5722,0.5722) -- (-0.5394,0.5394) -- (-0.5070,0.5070) -- (-0.4749,0.4749) -- (-0.4431,0.4431) -- (-0.4116,0.4116) -- (-0.3803,0.3803) -- (-0.3493,0.3493) -- (-0.3186,0.3186) -- (-0.2880,0.2880) -- (-0.2576,0.2576) -- (-0.2274,0.2274) -- (-0.1973,0.1973) -- (-0.1674,0.1674) -- (-0.1375,0.1375) -- (-0.1077,0.1077) -- (-0.0780,0.0780) -- (-0.0484,0.0484) -- (-0.0187,0.0187) -- (0.0109,-0.0109) -- (0.0405,-0.0405) -- (0.0702,-0.0702) -- (0.0999,-0.0999) -- (0.1296,-0.1296) -- (0.1595,-0.1595) -- (0.1894,-0.1894) -- (0.2194,-0.2194) -- (0.2496,-0.2496) -- (0.2800,-0.2800) -- (0.3105,-0.3105) -- (0.3412,-0.3412) -- (0.3721,-0.3721) -- (0.4033,-0.4033) -- (0.4347,-0.4347) -- (0.4664,-0.4664) -- (0.4984,-0.4984) -- (0.5308,-0.5308) -- (0.5635,-0.5635) -- (0.5966,-0.5966) -- (0.6300,-0.6300) -- (0.6640,-0.6640) -- (0.6984,-0.6984) -- (0.7332,-0.7332) -- (0.7687,-0.7687) -- (0.8047,-0.8047) -- (0.8412,-0.8412) -- (0.8785,-0.8785) -- (0.8785,-0.8785) -- (0.9164,-0.9164) -- (0.9550,-0.9550) -- (0.9944,-0.9944) -- (1.0346,-1.0346) -- (1.0756,-1.0756) -- (1.1176,-1.1176) -- (1.1606,-1.1606) -- (1.2046,-1.2046) -- (1.2497,-1.2497) -- (1.2960,-1.2960) -- (1.3435,-1.3435) -- (1.3924,-1.3924) -- (1.4427,-1.4427) -- (1.4945,-1.4945) -- (1.5479,-1.5479) -- (1.6030,-1.6030) -- (1.6600,-1.6600);
\draw[zone100] (1.6949,1.6949) -- (1.6367,1.6367) -- (1.5805,1.5805) -- (1.5260,1.5260) -- (1.4733,1.4733) -- (1.4221,1.4221) -- (1.3724,1.3724) -- (1.3241,1.3241) -- (1.2771,1.2771) -- (1.2313,1.2313) -- (1.1866,1.1866) -- (1.1431,1.1431) -- (1.1005,1.1005) -- (1.0589,1.0589) -- (1.0182,1.0182) -- (0.9783,0.9783) -- (0.9393,0.9393) -- (0.9009,0.9009) -- (0.8633,0.8633) -- (0.8263,0.8263) -- (0.7900,0.7900) -- (0.7543,0.7543) -- (0.7191,0.7191) -- (0.6844,0.6844) -- (0.6502,0.6502) -- (0.6164,0.6164) -- (0.5831,0.5831) -- (0.5502,0.5502) -- (0.5176,0.5176) -- (0.4854,0.4854) -- (0.4535,0.4535) -- (0.4220,0.4220) -- (0.3906,0.3906) -- (0.3596,0.3596) -- (0.3287,0.3287) -- (0.2981,0.2981) -- (0.2677,0.2677) -- (0.2374,0.2374) -- (0.2073,0.2073) -- (0.1773,0.1773) -- (0.1474,0.1474) -- (0.1176,0.1176) -- (0.0878,0.0878) -- (0.0582,0.0582) -- (0.0285,0.0285) -- (-0.0011,-0.0011) -- (-0.0307,-0.0307) -- (-0.0604,-0.0604) -- (-0.0900,-0.0900) -- (-0.1198,-0.1198) -- (-0.1496,-0.1496) -- (-0.1795,-0.1795) -- (-0.2095,-0.2095) -- (-0.2396,-0.2396) -- (-0.2699,-0.2699) -- (-0.2699,-0.2699) -- (-0.3004,-0.3004) -- (-0.3310,-0.3310) -- (-0.3619,-0.3619) -- (-0.3930,-0.3930) -- (-0.4243,-0.4243) -- (-0.4559,-0.4559) -- (-0.4878,-0.4878) -- (-0.5200,-0.5200) -- (-0.5526,-0.5526) -- (-0.5856,-0.5856) -- (-0.6189,-0.6189) -- (-0.6527,-0.6527) -- (-0.6869,-0.6869) -- (-0.7217,-0.7217) -- (-0.7569,-0.7569) -- (-0.7927,-0.7927) -- (-0.8291,-0.8291) -- (-0.8661,-0.8661) -- (-0.9038,-0.9038) -- (-0.9421,-0.9421) -- (-0.9813,-0.9813) -- (-1.0212,-1.0212) -- (-1.0620,-1.0620) -- (-1.1036,-1.1036) -- (-1.1463,-1.1463) -- (-1.1899,-1.1899) -- (-1.2347,-1.2347) -- (-1.2806,-1.2806) -- (-1.3277,-1.3277) -- (-1.3761,-1.3761) -- (-1.4259,-1.4259) -- (-1.4772,-1.4772) -- (-1.5300,-1.5300) -- (-1.5846,-1.5846) -- (-1.6410,-1.6410);
\draw[zone110] (-0.0571,2.3415) -- (-0.1121,2.2824) -- (-0.1645,2.2233) -- (-0.2144,2.1644) -- (-0.2619,2.1058) -- (-0.3071,2.0474) -- (-0.3501,1.9893) -- (-0.3910,1.9314) -- (-0.4299,1.8739) -- (-0.4669,1.8168) -- (-0.5022,1.7599) -- (-0.5357,1.7035) -- (-0.5675,1.6474) -- (-0.5978,1.5916) -- (-0.6265,1.5362) -- (-0.6538,1.4812) -- (-0.6798,1.4266) -- (-0.7043,1.3723) -- (-0.7276,1.3183) -- (-0.7497,1.2647) -- (-0.7706,1.2114) -- (-0.7903,1.1584) -- (-0.8089,1.1057) -- (-0.8265,1.0534) -- (-0.8431,1.0013) -- (-0.8586,0.9495) -- (-0.8732,0.8980) -- (-0.8868,0.8467) -- (-0.8995,0.7957) -- (-0.9114,0.7449) -- (-0.9223,0.6943) -- (-0.9325,0.6439) -- (-0.9418,0.5937) -- (-0.9503,0.5437) -- (-0.9580,0.4938) -- (-0.9649,0.4441) -- (-0.9711,0.3945) -- (-0.9765,0.3450) -- (-0.9812,0.2956) -- (-0.9852,0.2463) -- (-0.9852,0.2463) -- (-0.9884,0.1971) -- (-0.9909,0.1479) -- (-0.9927,0.0988) -- (-0.9937,0.0497) -- (-0.9941,0.0006) -- (-0.9938,-0.0485) -- (-0.9927,-0.0976) -- (-0.9909,-0.1467) -- (-0.9885,-0.1958) -- (-0.9853,-0.2451) -- (-0.9813,-0.2944) -- (-0.9767,-0.3437) -- (-0.9713,-0.3932) -- (-0.9651,-0.4428) -- (-0.9582,-0.4925) -- (-0.9505,-0.5424) -- (-0.9420,-0.5924) -- (-0.9327,-0.6426) -- (-0.9226,-0.6930) -- (-0.9117,-0.7436) -- (-0.8998,-0.7944) -- (-0.8871,-0.8454) -- (-0.8735,-0.8967) -- (-0.8590,-0.9482) -- (-0.8435,-1.0000) -- (-0.8269,-1.0521) -- (-0.8094,-1.1044) -- (-0.7908,-1.1571) -- (-0.7711,-1.2100) -- (-0.7502,-1.2633) -- (-0.7282,-1.3169) -- (-0.7049,-1.3709) -- (-0.6804,-1.4252) -- (-0.6545,-1.4799) -- (-0.6272,-1.5349) -- (-0.5985,-1.5902) -- (-0.5683,-1.6460) -- (-0.5365,-1.7021) -- (-0.5030,-1.7585) -- (-0.4678,-1.8153) -- (-0.4309,-1.8725) -- (-0.3920,-1.9300) -- (-0.3511,-1.9878) -- (-0.3082,-2.0459) -- (-0.2630,-2.1043) -- (-0.2156,-2.1630) -- (-0.1658,-2.2218) -- (-0.1135,-2.2809) -- (-0.0585,-2.3400) -- (-0.0007,-2.3993);
\draw[zone110] (0.0216,-2.3782) -- (0.0784,-2.3190) -- (0.1324,-2.2599) -- (0.1838,-2.2009) -- (0.2327,-2.1421) -- (0.2793,-2.0836) -- (0.3237,-2.0253) -- (0.3659,-1.9672) -- (0.4060,-1.9095) -- (0.4442,-1.8522) -- (0.4805,-1.7951) -- (0.5151,-1.7384) -- (0.5480,-1.6821) -- (0.5792,-1.6261) -- (0.6089,-1.5705) -- (0.6371,-1.5153) -- (0.6639,-1.4604) -- (0.6893,-1.4059) -- (0.7133,-1.3517) -- (0.7362,-1.2979) -- (0.7578,-1.2444) -- (0.7782,-1.1912) -- (0.7975,-1.1383) -- (0.8158,-1.0858) -- (0.8329,-1.0335) -- (0.8491,-0.9816) -- (0.8642,-0.9299) -- (0.8785,-0.8785) -- (0.8917,-0.8273) -- (0.9041,-0.7763) -- (0.9156,-0.7256) -- (0.9263,-0.6751) -- (0.9361,-0.6248) -- (0.9451,-0.5746) -- (0.9533,-0.5247) -- (0.9607,-0.4749) -- (0.9674,-0.4252) -- (0.9733,-0.3756) -- (0.9784,-0.3262) -- (0.9828,-0.2768) -- (0.9865,-0.2276) -- (0.9894,-0.1784) -- (0.9917,-0.1292) -- (0.9932,-0.0801) -- (0.9940,-0.0310) -- (0.9941,0.0181) -- (0.9934,0.0671) -- (0.9921,0.1162) -- (0.9901,0.1654) -- (0.9873,0.2146) -- (0.9838,0.2638) -- (0.9796,0.3131) -- (0.9747,0.3625) -- (0.9690,0.4121) -- (0.9626,0.4617) -- (0.9554,0.5115) -- (0.9474,0.5614) -- (0.9386,0.6115) -- (0.9290,0.6618) -- (0.9185,0.7122) -- (0.9073,0.7629) -- (0.8951,0.8138) -- (0.8821,0.8649) -- (0.8681,0.9163) -- (0.8532,0.9679) -- (0.8373,1.0198) -- (0.8204,1.0719) -- (0.8025,1.1244) -- (0.7834,1.1772) -- (0.7633,1.2303) -- (0.7420,1.2837) -- (0.7195,1.3374) -- (0.6958,1.3915) -- (0.6707,1.4459) -- (0.6443,1.5007) -- (0.6165,1.5559) -- (0.5872,1.6114) -- (0.5564,1.6673) -- (0.5240,1.7235) -- (0.4899,1.7801) -- (0.4540,1.8370) -- (0.4163,1.8943) -- (0.3767,1.9520) -- (0.3350,2.0099) -- (0.2913,2.0681) -- (0.2453,2.1266) -- (0.1970,2.1853) -- (0.1462,2.2443) -- (0.0929,2.3034) -- (0.0369,2.3626);
\draw[zone110] (-2.3626,-0.0369) -- (-2.3034,-0.0929) -- (-2.2443,-0.1462) -- (-2.1853,-0.1970) -- (-2.1266,-0.2453) -- (-2.0681,-0.2913) -- (-2.0099,-0.3350) -- (-1.9520,-0.3767) -- (-1.8943,-0.4163) -- (-1.8370,-0.4540) -- (-1.7801,-0.4899) -- (-1.7235,-0.5240) -- (-1.6673,-0.5564) -- (-1.6114,-0.5872) -- (-1.5559,-0.6165) -- (-1.5007,-0.6443) -- (-1.4459,-0.6707) -- (-1.3915,-0.6958) -- (-1.3374,-0.7195) -- (-1.2837,-0.7420) -- (-1.2303,-0.7633) -- (-1.1772,-0.7834) -- (-1.1244,-0.8025) -- (-1.0719,-0.8204) -- (-1.0198,-0.8373) -- (-0.9679,-0.8532) -- (-0.9163,-0.8681) -- (-0.8649,-0.8821) -- (-0.8138,-0.8951) -- (-0.7629,-0.9073) -- (-0.7122,-0.9185) -- (-0.6618,-0.9290) -- (-0.6115,-0.9386) -- (-0.5614,-0.9474) -- (-0.5115,-0.9554) -- (-0.4617,-0.9626) -- (-0.4121,-0.9690) -- (-0.3625,-0.9747) -- (-0.3131,-0.9796) -- (-0.2638,-0.9838) -- (-0.2146,-0.9873) -- (-0.1654,-0.9901) -- (-0.1162,-0.9921) -- (-0.0671,-0.9934) -- (-0.0181,-0.9941) -- (0.0310,-0.9940) -- (0.0801,-0.9932) -- (0.1292,-0.9917) -- (0.1784,-0.9894) -- (0.2276,-0.9865) -- (0.2768,-0.9828) -- (0.3262,-0.9784) -- (0.3756,-0.9733) -- (0.4252,-0.9674) -- (0.4749,-0.9607) -- (0.5247,-0.9533) -- (0.5746,-0.9451) -- (0.6248,-0.9361) -- (0.6751,-0.9263) -- (0.7256,-0.9156) -- (0.7763,-0.9041) -- (0.8273,-0.8917) -- (0.8785,-0.8785) -- (0.9299,-0.8642) -- (0.9816,-0.8491) -- (1.0335,-0.8329) -- (1.0858,-0.8158) -- (1.1383,-0.7975) -- (1.1912,-0.7782) -- (1.2444,-0.7578) -- (1.2979,-0.7362) -- (1.3517,-0.7133) -- (1.4059,-0.6893) -- (1.4604,-0.6639) -- (1.5153,-0.6371) -- (1.5705,-0.6089) -- (1.6261,-0.5792) -- (1.6821,-0.5480) -- (1.7384,-0.5151) -- (1.7951,-0.4805) -- (1.8522,-0.4442) -- (1.9095,-0.4060) -- (1.9672,-0.3659) -- (2.0253,-0.3237) -- (2.0836,-0.2793) -- (2.1421,-0.2327) -- (2.2009,-0.1838) -- (2.2599,-0.1324) -- (2.3190,-0.0784) -- (2.3782,-0.0216);
\draw[zone110] (2.3625,0.0369) -- (2.3033,0.0929) -- (2.2442,0.1463) -- (2.1853,0.1970) -- (2.1266,0.2453) -- (2.0681,0.2913) -- (2.0098,0.3351) -- (1.9519,0.3767) -- (1.8943,0.4163) -- (1.8370,0.4540) -- (1.7800,0.4899) -- (1.7235,0.5240) -- (1.6672,0.5564) -- (1.6113,0.5872) -- (1.5558,0.6165) -- (1.5007,0.6443) -- (1.4459,0.6707) -- (1.3915,0.6958) -- (1.3374,0.7195) -- (1.2836,0.7420) -- (1.2302,0.7633) -- (1.1771,0.7834) -- (1.1244,0.8025) -- (1.0719,0.8204) -- (1.0197,0.8373) -- (0.9678,0.8532) -- (0.9162,0.8681) -- (0.8649,0.8821) -- (0.8137,0.8951) -- (0.7629,0.9073) -- (0.7122,0.9186) -- (0.6617,0.9290) -- (0.6115,0.9386) -- (0.5614,0.9474) -- (0.5114,0.9554) -- (0.4617,0.9626) -- (0.4120,0.9690) -- (0.3625,0.9747) -- (0.3131,0.9796) -- (0.2638,0.9838) -- (0.2145,0.9873) -- (0.1653,0.9901) -- (0.1162,0.9921) -- (0.0671,0.9934) -- (0.0180,0.9941) -- (-0.0311,0.9940) -- (-0.0802,0.9932) -- (-0.1293,0.9917) -- (-0.1784,0.9894) -- (-0.2276,0.9865) -- (-0.2769,0.9828) -- (-0.3262,0.9784) -- (-0.3757,0.9733) -- (-0.4252,0.9674) -- (-0.4749,0.9607) -- (-0.5247,0.9533) -- (-0.5747,0.9451) -- (-0.6248,0.9361) -- (-0.6751,0.9263) -- (-0.7257,0.9156) -- (-0.7764,0.9041) -- (-0.8273,0.8917) -- (-0.8785,0.8784) -- (-0.9299,0.8642) -- (-0.9816,0.8491) -- (-1.0336,0.8329) -- (-1.0858,0.8157) -- (-1.1384,0.7975) -- (-1.1912,0.7782) -- (-1.2444,0.7578) -- (-1.2979,0.7362) -- (-1.3517,0.7133) -- (-1.4059,0.6892) -- (-1.4604,0.6638) -- (-1.4604,0.6638) -- (-1.5153,0.6371) -- (-1.5706,0.6089) -- (-1.6262,0.5792) -- (-1.6822,0.5479) -- (-1.7385,0.5151) -- (-1.7952,0.4805) -- (-1.8522,0.4442) -- (-1.9096,0.4060) -- (-1.9673,0.3658) -- (-2.0253,0.3236) -- (-2.0836,0.2793) -- (-2.1422,0.2327) -- (-2.2010,0.1838) -- (-2.2599,0.1323) -- (-2.3191,0.0783) -- (-2.3783,0.0215);
\draw[fmAmber!80!black, line width=0.8pt] (0.0000,0.0000) -- (0.0242,0.0000) -- (0.0483,0.0000) -- (0.0725,0.0000) -- (0.0967,0.0000) -- (0.1209,0.0000) -- (0.1452,0.0000) -- (0.1694,0.0000) -- (0.1937,0.0000) -- (0.2181,0.0000) -- (0.2425,0.0000) -- (0.2669,0.0000) -- (0.2914,0.0000) -- (0.3160,0.0000) -- (0.3406,0.0000) -- (0.3653,0.0000) -- (0.3900,0.0000) -- (0.4149,0.0000) -- (0.4398,0.0000) -- (0.4648,0.0000) -- (0.4900,0.0000) -- (0.5152,0.0000) -- (0.5405,0.0000) -- (0.5660,0.0000) -- (0.5915,0.0000) -- (0.6172,0.0000) -- (0.6431,0.0000) -- (0.6691,0.0000) -- (0.6952,0.0000) -- (0.7214,0.0000) -- (0.7479,0.0000) -- (0.7745,0.0000) -- (0.8012,0.0000) -- (0.8282,0.0000) -- (0.8553,0.0000) -- (0.8827,0.0000) -- (0.9102,0.0000) -- (0.9379,0.0000) -- (0.9659,0.0000) -- (0.9941,0.0000) -- (0.9941,0.0000) -- (0.9940,0.0222) -- (0.9938,0.0444) -- (0.9935,0.0666) -- (0.9930,0.0888) -- (0.9923,0.1110) -- (0.9915,0.1332) -- (0.9906,0.1554) -- (0.9895,0.1776) -- (0.9882,0.1998) -- (0.9868,0.2221) -- (0.9853,0.2444) -- (0.9836,0.2666) -- (0.9818,0.2889) -- (0.9798,0.3112) -- (0.9777,0.3336) -- (0.9754,0.3559) -- (0.9730,0.3783) -- (0.9704,0.4007) -- (0.9676,0.4231) -- (0.9647,0.4455) -- (0.9617,0.4680) -- (0.9585,0.4905) -- (0.9551,0.5130) -- (0.9516,0.5356) -- (0.9479,0.5582) -- (0.9441,0.5808) -- (0.9400,0.6034) -- (0.9359,0.6261) -- (0.9315,0.6489) -- (0.9270,0.6716) -- (0.9223,0.6944) -- (0.9175,0.7173) -- (0.9124,0.7402) -- (0.9072,0.7631) -- (0.9018,0.7861) -- (0.8963,0.8091) -- (0.8905,0.8322) -- (0.8846,0.8553) -- (0.8785,0.8785) -- (0.8785,0.8785) -- (0.8523,0.8523) -- (0.8264,0.8264) -- (0.8008,0.8008) -- (0.7756,0.7756) -- (0.7506,0.7506) -- (0.7259,0.7259) -- (0.7014,0.7014) -- (0.6772,0.6772) -- (0.6532,0.6532) -- (0.6295,0.6295) -- (0.6059,0.6059) -- (0.5826,0.5826) -- (0.5594,0.5594) -- (0.5365,0.5365) -- (0.5137,0.5137) -- (0.4911,0.4911) -- (0.4687,0.4687) -- (0.4464,0.4464) -- (0.4242,0.4242) -- (0.4022,0.4022) -- (0.3803,0.3803) -- (0.3585,0.3585) -- (0.3369,0.3369) -- (0.3153,0.3153) -- (0.2939,0.2939) -- (0.2725,0.2725) -- (0.2512,0.2512) -- (0.2300,0.2300) -- (0.2089,0.2089) -- (0.1878,0.1878) -- (0.1668,0.1668) -- (0.1459,0.1459) -- (0.1249,0.1249) -- (0.1041,0.1041) -- (0.0832,0.0832) -- (0.0624,0.0624) -- (0.0416,0.0416) -- (0.0208,0.0208) -- (0.0000,0.0000) -- cycle;
\end{scope}
\draw[atomLine, line width=0.7pt] (0,0) circle (2.4);
\begin{scope}
\clip (3.2500,-2.3500) rectangle (8.0000,2.4000);
\fill[fmAmber!8] (3.5500,-2.0500) -- (3.6366,-1.9634) -- (3.7232,-1.8768) -- (3.8099,-1.7901) -- (3.8967,-1.7033) -- (3.9836,-1.6164) -- (4.0706,-1.5294) -- (4.1577,-1.4423) -- (4.2451,-1.3549) -- (4.3326,-1.2674) -- (4.4204,-1.1796) -- (4.5084,-1.0916) -- (4.5968,-1.0032) -- (4.6855,-0.9145) -- (4.7745,-0.8255) -- (4.8639,-0.7361) -- (4.9537,-0.6463) -- (5.0439,-0.5561) -- (5.1346,-0.4654) -- (5.2258,-0.3742) -- (5.3176,-0.2824) -- (5.4099,-0.1901) -- (5.5028,-0.0972) -- (5.5963,-0.0037) -- (5.6905,0.0905) -- (5.7854,0.1854) -- (5.8810,0.2810) -- (5.9774,0.3774) -- (6.0747,0.4747) -- (6.1727,0.5727) -- (6.2717,0.6717) -- (6.3716,0.7716) -- (6.4725,0.8725) -- (6.5744,0.9744) -- (6.6774,1.0774) -- (6.7815,1.1815) -- (6.8868,1.2868) -- (6.9933,1.3933) -- (7.1011,1.5011) -- (7.2103,1.6103) -- (7.2103,1.6103) -- (7.1137,1.6358) -- (7.0174,1.6605) -- (6.9213,1.6844) -- (6.8253,1.7076) -- (6.7296,1.7301) -- (6.6340,1.7518) -- (6.5387,1.7727) -- (6.4435,1.7930) -- (6.3484,1.8125) -- (6.2536,1.8313) -- (6.1589,1.8494) -- (6.0643,1.8668) -- (5.9699,1.8836) -- (5.8757,1.8996) -- (5.7815,1.9150) -- (5.6876,1.9297) -- (5.5937,1.9437) -- (5.5000,1.9571) -- (5.4064,1.9698) -- (5.3129,1.9818) -- (5.2195,1.9933) -- (5.1262,2.0040) -- (5.0329,2.0142) -- (4.9398,2.0237) -- (4.8468,2.0326) -- (4.7538,2.0408) -- (4.6610,2.0484) -- (4.5681,2.0554) -- (4.4754,2.0618) -- (4.3827,2.0676) -- (4.2900,2.0728) -- (4.1974,2.0773) -- (4.1049,2.0812) -- (4.0123,2.0846) -- (3.9198,2.0873) -- (3.8274,2.0894) -- (3.7349,2.0909) -- (3.6424,2.0918) -- (3.5500,2.0921) -- (3.5500,2.0921) -- (3.5500,1.9747) -- (3.5500,1.8581) -- (3.5500,1.7425) -- (3.5500,1.6278) -- (3.5500,1.5139) -- (3.5500,1.4008) -- (3.5500,1.2885) -- (3.5500,1.1769) -- (3.5500,1.0661) -- (3.5500,0.9560) -- (3.5500,0.8465) -- (3.5500,0.7377) -- (3.5500,0.6295) -- (3.5500,0.5219) -- (3.5500,0.4148) -- (3.5500,0.3082) -- (3.5500,0.2022) -- (3.5500,0.0966) -- (3.5500,-0.0085) -- (3.5500,-0.1132) -- (3.5500,-0.2174) -- (3.5500,-0.3213) -- (3.5500,-0.4248) -- (3.5500,-0.5280) -- (3.5500,-0.6309) -- (3.5500,-0.7335) -- (3.5500,-0.8358) -- (3.5500,-0.9378) -- (3.5500,-1.0397) -- (3.5500,-1.1413) -- (3.5500,-1.2427) -- (3.5500,-1.3440) -- (3.5500,-1.4451) -- (3.5500,-1.5461) -- (3.5500,-1.6470) -- (3.5500,-1.7478) -- (3.5500,-1.8486) -- (3.5500,-1.9493) -- (3.5500,-2.0500) -- cycle;
\fill[fmAmber!25] (3.5500,-2.0500) -- (3.6507,-2.0500) -- (3.7514,-2.0500) -- (3.8522,-2.0500) -- (3.9530,-2.0500) -- (4.0539,-2.0500) -- (4.1549,-2.0500) -- (4.2560,-2.0500) -- (4.3573,-2.0500) -- (4.4587,-2.0500) -- (4.5603,-2.0500) -- (4.6622,-2.0500) -- (4.7642,-2.0500) -- (4.8665,-2.0500) -- (4.9691,-2.0500) -- (5.0720,-2.0500) -- (5.1752,-2.0500) -- (5.2787,-2.0500) -- (5.3826,-2.0500) -- (5.4868,-2.0500) -- (5.5915,-2.0500) -- (5.6966,-2.0500) -- (5.8022,-2.0500) -- (5.9082,-2.0500) -- (6.0148,-2.0500) -- (6.1219,-2.0500) -- (6.2295,-2.0500) -- (6.3377,-2.0500) -- (6.4465,-2.0500) -- (6.5560,-2.0500) -- (6.6661,-2.0500) -- (6.7769,-2.0500) -- (6.8885,-2.0500) -- (7.0008,-2.0500) -- (7.1139,-2.0500) -- (7.2278,-2.0500) -- (7.3425,-2.0500) -- (7.4581,-2.0500) -- (7.5747,-2.0500) -- (7.6921,-2.0500) -- (7.6921,-2.0500) -- (7.6918,-1.9576) -- (7.6909,-1.8651) -- (7.6894,-1.7726) -- (7.6873,-1.6802) -- (7.6846,-1.5877) -- (7.6812,-1.4951) -- (7.6773,-1.4026) -- (7.6728,-1.3100) -- (7.6676,-1.2173) -- (7.6618,-1.1246) -- (7.6554,-1.0319) -- (7.6484,-0.9390) -- (7.6408,-0.8462) -- (7.6326,-0.7532) -- (7.6237,-0.6602) -- (7.6142,-0.5671) -- (7.6040,-0.4738) -- (7.5933,-0.3805) -- (7.5818,-0.2871) -- (7.5698,-0.1936) -- (7.5571,-0.1000) -- (7.5437,-0.0063) -- (7.5297,0.0876) -- (7.5150,0.1815) -- (7.4996,0.2757) -- (7.4836,0.3699) -- (7.4668,0.4643) -- (7.4494,0.5589) -- (7.4313,0.6536) -- (7.4125,0.7484) -- (7.3930,0.8435) -- (7.3727,0.9387) -- (7.3518,1.0340) -- (7.3301,1.1296) -- (7.3076,1.2253) -- (7.2844,1.3213) -- (7.2605,1.4174) -- (7.2358,1.5137) -- (7.2103,1.6103) -- (7.2103,1.6103) -- (7.1011,1.5011) -- (6.9933,1.3933) -- (6.8868,1.2868) -- (6.7815,1.1815) -- (6.6774,1.0774) -- (6.5744,0.9744) -- (6.4725,0.8725) -- (6.3716,0.7716) -- (6.2717,0.6717) -- (6.1727,0.5727) -- (6.0747,0.4747) -- (5.9774,0.3774) -- (5.8810,0.2810) -- (5.7854,0.1854) -- (5.6905,0.0905) -- (5.5963,-0.0037) -- (5.5028,-0.0972) -- (5.4099,-0.1901) -- (5.3176,-0.2824) -- (5.2258,-0.3742) -- (5.1346,-0.4654) -- (5.0439,-0.5561) -- (4.9537,-0.6463) -- (4.8639,-0.7361) -- (4.7745,-0.8255) -- (4.6855,-0.9145) -- (4.5968,-1.0032) -- (4.5084,-1.0916) -- (4.4204,-1.1796) -- (4.3326,-1.2674) -- (4.2451,-1.3549) -- (4.1577,-1.4423) -- (4.0706,-1.5294) -- (3.9836,-1.6164) -- (3.8967,-1.7033) -- (3.8099,-1.7901) -- (3.7232,-1.8768) -- (3.6366,-1.9634) -- (3.5500,-2.0500) -- cycle;
\draw[zone100] (3.5500,7.8534) -- (3.5500,7.5135) -- (3.5500,7.1848) -- (3.5500,6.8666) -- (3.5500,6.5580) -- (3.5500,6.2586) -- (3.5500,5.9678) -- (3.5500,5.6850) -- (3.5500,5.4097) -- (3.5500,5.1415) -- (3.5500,4.8800) -- (3.5500,4.6247) -- (3.5500,4.3753) -- (3.5500,4.1314) -- (3.5500,3.8928) -- (3.5500,3.6590) -- (3.5500,3.4298) -- (3.5500,3.2050) -- (3.5500,2.9843) -- (3.5500,2.7674) -- (3.5500,2.5541) -- (3.5500,2.3443) -- (3.5500,2.1376) -- (3.5500,1.9339) -- (3.5500,1.7331) -- (3.5500,1.5348) -- (3.5500,1.3391) -- (3.5500,1.1456) -- (3.5500,0.9543) -- (3.5500,0.7650) -- (3.5500,0.5776) -- (3.5500,0.3918) -- (3.5500,0.2076) -- (3.5500,0.0249) -- (3.5500,-0.1565) -- (3.5500,-0.3367) -- (3.5500,-0.5159) -- (3.5500,-0.6940) -- (3.5500,-0.8714) -- (3.5500,-1.0480) -- (3.5500,-1.2240) -- (3.5500,-1.3995) -- (3.5500,-1.5746) -- (3.5500,-1.7494) -- (3.5500,-1.9240) -- (3.5500,-2.0985) -- (3.5500,-2.2731) -- (3.5500,-2.4478) -- (3.5500,-2.6227) -- (3.5500,-2.7980) -- (3.5500,-2.9738) -- (3.5500,-3.1501) -- (3.5500,-3.3271) -- (3.5500,-3.5049) -- (3.5500,-3.6836) -- (3.5500,-3.8633) -- (3.5500,-3.8633) -- (3.5500,-4.0442) -- (3.5500,-4.2263) -- (3.5500,-4.4099) -- (3.5500,-4.5949) -- (3.5500,-4.7816) -- (3.5500,-4.9700) -- (3.5500,-5.1604) -- (3.5500,-5.3529) -- (3.5500,-5.5476) -- (3.5500,-5.7448) -- (3.5500,-5.9444) -- (3.5500,-6.1468) -- (3.5500,-6.3521) -- (3.5500,-6.5606) -- (3.5500,-6.7723) -- (3.5500,-6.9875) -- (3.5500,-7.2065) -- (3.5500,-7.4295) -- (3.5500,-7.6567) -- (3.5500,-7.8884) -- (3.5500,-8.1248) -- (3.5500,-8.3664) -- (3.5500,-8.6133) -- (3.5500,-8.8659) -- (3.5500,-9.1246) -- (3.5500,-9.3898) -- (3.5500,-9.6619) -- (3.5500,-9.9413) -- (3.5500,-10.2285) -- (3.5500,-10.5240) -- (3.5500,-10.8284) -- (3.5500,-11.1423) -- (3.5500,-11.4663) -- (3.5500,-11.8011);
\draw[zone100] (-6.2708,-2.0500) -- (-5.9336,-2.0500) -- (-5.6075,-2.0500) -- (-5.2916,-2.0500) -- (-4.9853,-2.0500) -- (-4.6881,-2.0500) -- (-4.3992,-2.0500) -- (-4.1182,-2.0500) -- (-3.8447,-2.0500) -- (-3.5782,-2.0500) -- (-3.3182,-2.0500) -- (-3.0643,-2.0500) -- (-2.8163,-2.0500) -- (-2.5737,-2.0500) -- (-2.3362,-2.0500) -- (-2.1036,-2.0500) -- (-1.8755,-2.0500) -- (-1.6517,-2.0500) -- (-1.4319,-2.0500) -- (-1.2159,-2.0500) -- (-1.0034,-2.0500) -- (-0.7944,-2.0500) -- (-0.5884,-2.0500) -- (-0.3855,-2.0500) -- (-0.1852,-2.0500) -- (0.0124,-2.0500) -- (0.2076,-2.0500) -- (0.4005,-2.0500) -- (0.5913,-2.0500) -- (0.7802,-2.0500) -- (0.9672,-2.0500) -- (1.1526,-2.0500) -- (1.3364,-2.0500) -- (1.5188,-2.0500) -- (1.6999,-2.0500) -- (1.8798,-2.0500) -- (2.0587,-2.0500) -- (2.2367,-2.0500) -- (2.4138,-2.0500) -- (2.5903,-2.0500) -- (2.7662,-2.0500) -- (2.9415,-2.0500) -- (3.1166,-2.0500) -- (3.2913,-2.0500) -- (3.4659,-2.0500) -- (3.6404,-2.0500) -- (3.8150,-2.0500) -- (3.9898,-2.0500) -- (4.1648,-2.0500) -- (4.3402,-2.0500) -- (4.5161,-2.0500) -- (4.6926,-2.0500) -- (4.8697,-2.0500) -- (5.0477,-2.0500) -- (5.2267,-2.0500) -- (5.4067,-2.0500) -- (5.5878,-2.0500) -- (5.7703,-2.0500) -- (5.9541,-2.0500) -- (6.1396,-2.0500) -- (6.3267,-2.0500) -- (6.5156,-2.0500) -- (6.7065,-2.0500) -- (6.8995,-2.0500) -- (7.0947,-2.0500) -- (7.2925,-2.0500) -- (7.4928,-2.0500) -- (7.6958,-2.0500) -- (7.9019,-2.0500) -- (8.1111,-2.0500) -- (8.3236,-2.0500) -- (8.5398,-2.0500) -- (8.7597,-2.0500) -- (8.9837,-2.0500) -- (9.2119,-2.0500) -- (9.4447,-2.0500) -- (9.6824,-2.0500) -- (9.9252,-2.0500) -- (10.1734,-2.0500) -- (10.4275,-2.0500) -- (10.6877,-2.0500) -- (10.9545,-2.0500) -- (11.2283,-2.0500) -- (11.5095,-2.0500) -- (11.7987,-2.0500) -- (12.0963,-2.0500) -- (12.4029,-2.0500) -- (12.7191,-2.0500) -- (13.0457,-2.0500) -- (13.3832,-2.0500);
\draw[zone100] (-3.4309,4.9309) -- (-3.1913,4.6913) -- (-2.9595,4.4595) -- (-2.7351,4.2351) -- (-2.5175,4.0175) -- (-2.3064,3.8064) -- (-2.1013,3.6013) -- (-1.9018,3.4018) -- (-1.7076,3.2076) -- (-1.5184,3.0184) -- (-1.3339,2.8339) -- (-1.1537,2.6537) -- (-0.9777,2.4777) -- (-0.8056,2.3056) -- (-0.6372,2.1372) -- (-0.4722,1.9722) -- (-0.3104,1.8104) -- (-0.1517,1.6517) -- (0.0041,1.4959) -- (0.1572,1.3428) -- (0.3078,1.1922) -- (0.4560,1.0440) -- (0.6019,0.8981) -- (0.7458,0.7542) -- (0.8876,0.6124) -- (1.0276,0.4724) -- (1.1659,0.3341) -- (1.3026,0.1974) -- (1.4377,0.0623) -- (1.5714,-0.0714) -- (1.7039,-0.2039) -- (1.8351,-0.3351) -- (1.9653,-0.4653) -- (2.0944,-0.5944) -- (2.2226,-0.7226) -- (2.3499,-0.8499) -- (2.4765,-0.9765) -- (2.6025,-1.1025) -- (2.7278,-1.2278) -- (2.8527,-1.3527) -- (2.9771,-1.4771) -- (3.1011,-1.6011) -- (3.2249,-1.7249) -- (3.3485,-1.8485) -- (3.4720,-1.9720) -- (3.5954,-2.0954) -- (3.7188,-2.2188) -- (3.8424,-2.3424) -- (3.9661,-2.4661) -- (4.0901,-2.5901) -- (4.2144,-2.7144) -- (4.3391,-2.8391) -- (4.4643,-2.9643) -- (4.5901,-3.0901) -- (4.7165,-3.2165) -- (4.8437,-3.3437) -- (4.9716,-3.4716) -- (5.1005,-3.6005) -- (5.2304,-3.7304) -- (5.3613,-3.8613) -- (5.4934,-3.9934) -- (5.6268,-4.1268) -- (5.7616,-4.2616) -- (5.8978,-4.3978) -- (6.0357,-4.5357) -- (6.1752,-4.6752) -- (6.3165,-4.8165) -- (6.4598,-4.9598) -- (6.6052,-5.1052) -- (6.7528,-5.2528) -- (6.9027,-5.4027) -- (7.0552,-5.5552) -- (7.2103,-5.7103) -- (7.2103,-5.7103) -- (7.3682,-5.8682) -- (7.5291,-6.0291) -- (7.6932,-6.1932) -- (7.8607,-6.3607) -- (8.0319,-6.5319) -- (8.2068,-6.7068) -- (8.3858,-6.8858) -- (8.5692,-7.0692) -- (8.7571,-7.2571) -- (8.9499,-7.4499) -- (9.1480,-7.6480) -- (9.3516,-7.8516) -- (9.5611,-8.0611) -- (9.7769,-8.2769) -- (9.9995,-8.4995) -- (10.2293,-8.7293) -- (10.4667,-8.9667);
\draw[zone100] (10.6119,5.0119) -- (10.3696,4.7696) -- (10.1353,4.5353) -- (9.9085,4.3085) -- (9.6887,4.0887) -- (9.4755,3.8755) -- (9.2684,3.6684) -- (9.0671,3.4671) -- (8.8712,3.2712) -- (8.6804,3.0804) -- (8.4944,2.8944) -- (8.3128,2.7128) -- (8.1355,2.5355) -- (7.9621,2.3621) -- (7.7925,2.1925) -- (7.6264,2.0264) -- (7.4636,1.8636) -- (7.3039,1.7039) -- (7.1471,1.5471) -- (6.9931,1.3931) -- (6.8417,1.2417) -- (6.6927,1.0927) -- (6.5461,0.9461) -- (6.4016,0.8016) -- (6.2591,0.6591) -- (6.1185,0.5185) -- (5.9796,0.3796) -- (5.8425,0.2425) -- (5.7068,0.1068) -- (5.5726,-0.0274) -- (5.4398,-0.1602) -- (5.3081,-0.2919) -- (5.1776,-0.4224) -- (5.0482,-0.5518) -- (4.9197,-0.6803) -- (4.7921,-0.8079) -- (4.6652,-0.9348) -- (4.5391,-1.0609) -- (4.4136,-1.1864) -- (4.2886,-1.3114) -- (4.1640,-1.4360) -- (4.0398,-1.5602) -- (3.9160,-1.6840) -- (3.7923,-1.8077) -- (3.6688,-1.9312) -- (3.5454,-2.0546) -- (3.4220,-2.1780) -- (3.2985,-2.3015) -- (3.1748,-2.4252) -- (3.0509,-2.5491) -- (2.9267,-2.6733) -- (2.8022,-2.7978) -- (2.6771,-2.9229) -- (2.5515,-3.0485) -- (2.4253,-3.1747) -- (2.4253,-3.1747) -- (2.2984,-3.3016) -- (2.1708,-3.4292) -- (2.0422,-3.5578) -- (1.9127,-3.6873) -- (1.7821,-3.8179) -- (1.6504,-3.9496) -- (1.5174,-4.0826) -- (1.3831,-4.2169) -- (1.2474,-4.3526) -- (1.1101,-4.4899) -- (0.9711,-4.6289) -- (0.8304,-4.7696) -- (0.6878,-4.9122) -- (0.5431,-5.0569) -- (0.3963,-5.2037) -- (0.2471,-5.3529) -- (0.0955,-5.5045) -- (-0.0587,-5.6587) -- (-0.2156,-5.8156) -- (-0.3756,-5.9756) -- (-0.5386,-6.1386) -- (-0.7050,-6.3050) -- (-0.8749,-6.4749) -- (-1.0485,-6.6485) -- (-1.2262,-6.8262) -- (-1.4081,-7.0081) -- (-1.5945,-7.1945) -- (-1.7856,-7.3856) -- (-1.9819,-7.5819) -- (-2.1836,-7.7836) -- (-2.3912,-7.9912) -- (-2.6049,-8.2049) -- (-2.8251,-8.4251) -- (-3.0525,-8.6525) -- (-3.2873,-8.8873);
\draw[zone110] (3.3121,7.7063) -- (3.0827,7.4598) -- (2.8644,7.2138) -- (2.6566,6.9685) -- (2.4588,6.7241) -- (2.2706,6.4808) -- (2.0914,6.2386) -- (1.9209,5.9977) -- (1.7587,5.7581) -- (1.6044,5.5199) -- (1.4576,5.2831) -- (1.3180,5.0478) -- (1.1853,4.8141) -- (1.0592,4.5818) -- (0.9394,4.3510) -- (0.8256,4.1218) -- (0.7177,3.8940) -- (0.6153,3.6677) -- (0.5182,3.4429) -- (0.4262,3.2194) -- (0.3392,2.9974) -- (0.2570,2.7767) -- (0.1794,2.5572) -- (0.1062,2.3391) -- (0.0373,2.1221) -- (-0.0275,1.9063) -- (-0.0882,1.6916) -- (-0.1450,1.4780) -- (-0.1980,1.2653) -- (-0.2474,1.0537) -- (-0.2931,0.8429) -- (-0.3353,0.6329) -- (-0.3741,0.4237) -- (-0.4095,0.2152) -- (-0.4417,0.0074) -- (-0.4706,-0.1998) -- (-0.4963,-0.4064) -- (-0.5189,-0.6126) -- (-0.5384,-0.8184) -- (-0.5549,-1.0238) -- (-0.5549,-1.0238) -- (-0.5683,-1.2289) -- (-0.5787,-1.4337) -- (-0.5861,-1.6384) -- (-0.5906,-1.8430) -- (-0.5921,-2.0474) -- (-0.5907,-2.2519) -- (-0.5863,-2.4565) -- (-0.5789,-2.6611) -- (-0.5686,-2.8660) -- (-0.5552,-3.0711) -- (-0.5389,-3.2765) -- (-0.5194,-3.4822) -- (-0.4969,-3.6884) -- (-0.4713,-3.8951) -- (-0.4424,-4.1022) -- (-0.4104,-4.3100) -- (-0.3750,-4.5185) -- (-0.3363,-4.7277) -- (-0.2942,-4.9376) -- (-0.2486,-5.1484) -- (-0.1993,-5.3600) -- (-0.1464,-5.5727) -- (-0.0897,-5.7863) -- (-0.0290,-6.0009) -- (0.0356,-6.2167) -- (0.1044,-6.4336) -- (0.1775,-6.6518) -- (0.2550,-6.8712) -- (0.3371,-7.0918) -- (0.4240,-7.3139) -- (0.5158,-7.5373) -- (0.6128,-7.7621) -- (0.7150,-7.9883) -- (0.8229,-8.2161) -- (0.9365,-8.4453) -- (1.0562,-8.6760) -- (1.1821,-8.9082) -- (1.3147,-9.1420) -- (1.4541,-9.3772) -- (1.6007,-9.6139) -- (1.7548,-9.8521) -- (1.9168,-10.0917) -- (2.0871,-10.3326) -- (2.2660,-10.5747) -- (2.4540,-10.8180) -- (2.6515,-11.0624) -- (2.8591,-11.3077) -- (3.0771,-11.5536) -- (3.3062,-11.8001) -- (3.5469,-12.0469);
\draw[zone110] (3.6400,-11.9592) -- (3.8765,-11.7125) -- (4.1016,-11.4662) -- (4.3158,-11.2204) -- (4.5198,-10.9755) -- (4.7139,-10.7315) -- (4.8986,-10.4885) -- (5.0744,-10.2468) -- (5.2417,-10.0064) -- (5.4009,-9.7673) -- (5.5523,-9.5297) -- (5.6963,-9.2935) -- (5.8332,-9.0588) -- (5.9634,-8.8256) -- (6.0871,-8.5939) -- (6.2045,-8.3637) -- (6.3161,-8.1350) -- (6.4219,-7.9078) -- (6.5223,-7.6820) -- (6.6174,-7.4577) -- (6.7074,-7.2348) -- (6.7926,-7.0133) -- (6.8730,-6.7931) -- (6.9490,-6.5741) -- (7.0205,-6.3564) -- (7.0878,-6.1399) -- (7.1510,-5.9245) -- (7.2103,-5.7103) -- (7.2656,-5.4970) -- (7.3172,-5.2847) -- (7.3652,-5.0734) -- (7.4096,-4.8629) -- (7.4505,-4.6533) -- (7.4880,-4.4443) -- (7.5221,-4.2361) -- (7.5530,-4.0286) -- (7.5807,-3.8216) -- (7.6053,-3.6151) -- (7.6267,-3.4091) -- (7.6450,-3.2035) -- (7.6603,-2.9982) -- (7.6726,-2.7932) -- (7.6819,-2.5884) -- (7.6882,-2.3838) -- (7.6915,-2.1793) -- (7.6919,-1.9748) -- (7.6894,-1.7703) -- (7.6838,-1.5657) -- (7.6753,-1.3610) -- (7.6639,-1.1560) -- (7.6494,-0.9508) -- (7.6318,-0.7453) -- (7.6112,-0.5394) -- (7.5875,-0.3331) -- (7.5607,-0.1262) -- (7.5306,0.0812) -- (7.4973,0.2892) -- (7.4607,0.4980) -- (7.4207,0.7074) -- (7.3773,0.9177) -- (7.3303,1.1288) -- (7.2796,1.3408) -- (7.2253,1.5538) -- (7.1671,1.7678) -- (7.1049,1.9829) -- (7.0387,2.1991) -- (6.9683,2.4164) -- (6.8935,2.6351) -- (6.8143,2.8549) -- (6.7304,3.0761) -- (6.6417,3.2987) -- (6.5479,3.5226) -- (6.4490,3.7480) -- (6.3446,3.9748) -- (6.2346,4.2031) -- (6.1187,4.4329) -- (5.9967,4.6641) -- (5.8683,4.8970) -- (5.7332,5.1313) -- (5.5911,5.3671) -- (5.4416,5.6043) -- (5.2846,5.8430) -- (5.1195,6.0831) -- (4.9459,6.3245) -- (4.7636,6.5671) -- (4.5720,6.8109) -- (4.3707,7.0556) -- (4.1593,7.3011) -- (3.9371,7.5473) -- (3.7036,7.7940);
\draw[zone110] (-6.2940,-2.2036) -- (-6.0473,-2.4371) -- (-5.8011,-2.6593) -- (-5.5556,-2.8707) -- (-5.3109,-3.0720) -- (-5.0671,-3.2636) -- (-4.8245,-3.4459) -- (-4.5831,-3.6195) -- (-4.3430,-3.7846) -- (-4.1043,-3.9416) -- (-3.8671,-4.0911) -- (-3.6313,-4.2332) -- (-3.3970,-4.3683) -- (-3.1641,-4.4967) -- (-2.9329,-4.6187) -- (-2.7031,-4.7346) -- (-2.4748,-4.8446) -- (-2.2480,-4.9490) -- (-2.0226,-5.0479) -- (-1.7987,-5.1417) -- (-1.5761,-5.2304) -- (-1.3549,-5.3143) -- (-1.1351,-5.3935) -- (-0.9164,-5.4683) -- (-0.6991,-5.5387) -- (-0.4829,-5.6049) -- (-0.2678,-5.6671) -- (-0.0538,-5.7253) -- (0.1592,-5.7796) -- (0.3712,-5.8303) -- (0.5823,-5.8773) -- (0.7926,-5.9207) -- (1.0020,-5.9607) -- (1.2108,-5.9973) -- (1.4188,-6.0306) -- (1.6262,-6.0607) -- (1.8331,-6.0875) -- (2.0394,-6.1112) -- (2.2453,-6.1318) -- (2.4508,-6.1494) -- (2.6560,-6.1639) -- (2.8610,-6.1753) -- (3.0657,-6.1838) -- (3.2703,-6.1894) -- (3.4748,-6.1919) -- (3.6793,-6.1915) -- (3.8838,-6.1882) -- (4.0884,-6.1819) -- (4.2932,-6.1726) -- (4.4982,-6.1603) -- (4.7035,-6.1450) -- (4.9091,-6.1267) -- (5.1151,-6.1053) -- (5.3216,-6.0807) -- (5.5286,-6.0530) -- (5.7361,-6.0221) -- (5.9443,-5.9880) -- (6.1533,-5.9505) -- (6.3629,-5.9096) -- (6.5734,-5.8652) -- (6.7847,-5.8172) -- (6.9970,-5.7656) -- (7.2103,-5.7103) -- (7.4245,-5.6510) -- (7.6399,-5.5878) -- (7.8564,-5.5205) -- (8.0741,-5.4490) -- (8.2931,-5.3730) -- (8.5133,-5.2926) -- (8.7348,-5.2074) -- (8.9577,-5.1174) -- (9.1820,-5.0223) -- (9.4078,-4.9219) -- (9.6350,-4.8161) -- (9.8637,-4.7045) -- (10.0939,-4.5871) -- (10.3256,-4.4634) -- (10.5588,-4.3332) -- (10.7935,-4.1963) -- (11.0297,-4.0523) -- (11.2673,-3.9009) -- (11.5064,-3.7417) -- (11.7468,-3.5744) -- (11.9885,-3.3986) -- (12.2315,-3.2139) -- (12.4755,-3.0198) -- (12.7204,-2.8158) -- (12.9662,-2.6016) -- (13.2125,-2.3765) -- (13.4592,-2.1400);
\draw[zone110] (13.3938,-1.8962) -- (13.1471,-1.6627) -- (12.9009,-1.4406) -- (12.6554,-1.2291) -- (12.4107,-1.0278) -- (12.1669,-0.8362) -- (11.9243,-0.6539) -- (11.6829,-0.4804) -- (11.4428,-0.3153) -- (11.2041,-0.1582) -- (10.9669,-0.0088) -- (10.7311,0.1333) -- (10.4968,0.2684) -- (10.2640,0.3968) -- (10.0327,0.5188) -- (9.8029,0.6347) -- (9.5746,0.7447) -- (9.3478,0.8491) -- (9.1224,0.9480) -- (8.8985,1.0417) -- (8.6759,1.1305) -- (8.4547,1.2144) -- (8.2349,1.2936) -- (8.0163,1.3684) -- (7.7989,1.4388) -- (7.5827,1.5050) -- (7.3676,1.5671) -- (7.1536,1.6253) -- (6.9406,1.6797) -- (6.7286,1.7303) -- (6.5175,1.7773) -- (6.3072,1.8207) -- (6.0978,1.8607) -- (5.8891,1.8974) -- (5.6810,1.9307) -- (5.4736,1.9607) -- (5.2668,1.9875) -- (5.0604,2.0112) -- (4.8545,2.0318) -- (4.6490,2.0494) -- (4.4438,2.0639) -- (4.2389,2.0753) -- (4.0341,2.0838) -- (3.8295,2.0894) -- (3.6250,2.0919) -- (3.4206,2.0915) -- (3.2160,2.0882) -- (3.0114,2.0819) -- (2.8066,2.0726) -- (2.6016,2.0603) -- (2.3964,2.0450) -- (2.1907,2.0267) -- (1.9847,2.0052) -- (1.7783,1.9807) -- (1.5713,1.9530) -- (1.3637,1.9221) -- (1.1555,1.8879) -- (0.9466,1.8504) -- (0.7369,1.8095) -- (0.5264,1.7651) -- (0.3151,1.7172) -- (0.1028,1.6656) -- (-0.1104,1.6102) -- (-0.3247,1.5510) -- (-0.5401,1.4878) -- (-0.7566,1.4205) -- (-0.9743,1.3489) -- (-1.1932,1.2730) -- (-1.4135,1.1925) -- (-1.6350,1.1073) -- (-1.8579,1.0173) -- (-2.0822,0.9222) -- (-2.3080,0.8218) -- (-2.5352,0.7160) -- (-2.5352,0.7160) -- (-2.7639,0.6044) -- (-2.9941,0.4870) -- (-3.2258,0.3633) -- (-3.4590,0.2331) -- (-3.6937,0.0962) -- (-3.9299,-0.0479) -- (-4.1675,-0.1993) -- (-4.4066,-0.3584) -- (-4.6470,-0.5257) -- (-4.8888,-0.7015) -- (-5.1317,-0.8863) -- (-5.3757,-1.0804) -- (-5.6206,-1.2843) -- (-5.8664,-1.4986) -- (-6.1127,-1.7237) -- (-6.3594,-1.9602);
\draw[fmAmber!80!black, line width=0.8pt] (3.5500,-2.0500) -- (3.6507,-2.0500) -- (3.7514,-2.0500) -- (3.8522,-2.0500) -- (3.9530,-2.0500) -- (4.0539,-2.0500) -- (4.1549,-2.0500) -- (4.2560,-2.0500) -- (4.3573,-2.0500) -- (4.4587,-2.0500) -- (4.5603,-2.0500) -- (4.6622,-2.0500) -- (4.7642,-2.0500) -- (4.8665,-2.0500) -- (4.9691,-2.0500) -- (5.0720,-2.0500) -- (5.1752,-2.0500) -- (5.2787,-2.0500) -- (5.3826,-2.0500) -- (5.4868,-2.0500) -- (5.5915,-2.0500) -- (5.6966,-2.0500) -- (5.8022,-2.0500) -- (5.9082,-2.0500) -- (6.0148,-2.0500) -- (6.1219,-2.0500) -- (6.2295,-2.0500) -- (6.3377,-2.0500) -- (6.4465,-2.0500) -- (6.5560,-2.0500) -- (6.6661,-2.0500) -- (6.7769,-2.0500) -- (6.8885,-2.0500) -- (7.0008,-2.0500) -- (7.1139,-2.0500) -- (7.2278,-2.0500) -- (7.3425,-2.0500) -- (7.4581,-2.0500) -- (7.5747,-2.0500) -- (7.6921,-2.0500) -- (7.6921,-2.0500) -- (7.6918,-1.9576) -- (7.6909,-1.8651) -- (7.6894,-1.7726) -- (7.6873,-1.6802) -- (7.6846,-1.5877) -- (7.6812,-1.4951) -- (7.6773,-1.4026) -- (7.6728,-1.3100) -- (7.6676,-1.2173) -- (7.6618,-1.1246) -- (7.6554,-1.0319) -- (7.6484,-0.9390) -- (7.6408,-0.8462) -- (7.6326,-0.7532) -- (7.6237,-0.6602) -- (7.6142,-0.5671) -- (7.6040,-0.4738) -- (7.5933,-0.3805) -- (7.5818,-0.2871) -- (7.5698,-0.1936) -- (7.5571,-0.1000) -- (7.5437,-0.0063) -- (7.5297,0.0876) -- (7.5150,0.1815) -- (7.4996,0.2757) -- (7.4836,0.3699) -- (7.4668,0.4643) -- (7.4494,0.5589) -- (7.4313,0.6536) -- (7.4125,0.7484) -- (7.3930,0.8435) -- (7.3727,0.9387) -- (7.3518,1.0340) -- (7.3301,1.1296) -- (7.3076,1.2253) -- (7.2844,1.3213) -- (7.2605,1.4174) -- (7.2358,1.5137) -- (7.2103,1.6103) -- (7.2103,1.6103) -- (7.1011,1.5011) -- (6.9933,1.3933) -- (6.8868,1.2868) -- (6.7815,1.1815) -- (6.6774,1.0774) -- (6.5744,0.9744) -- (6.4725,0.8725) -- (6.3716,0.7716) -- (6.2717,0.6717) -- (6.1727,0.5727) -- (6.0747,0.4747) -- (5.9774,0.3774) -- (5.8810,0.2810) -- (5.7854,0.1854) -- (5.6905,0.0905) -- (5.5963,-0.0037) -- (5.5028,-0.0972) -- (5.4099,-0.1901) -- (5.3176,-0.2824) -- (5.2258,-0.3742) -- (5.1346,-0.4654) -- (5.0439,-0.5561) -- (4.9537,-0.6463) -- (4.8639,-0.7361) -- (4.7745,-0.8255) -- (4.6855,-0.9145) -- (4.5968,-1.0032) -- (4.5084,-1.0916) -- (4.4204,-1.1796) -- (4.3326,-1.2674) -- (4.2451,-1.3549) -- (4.1577,-1.4423) -- (4.0706,-1.5294) -- (3.9836,-1.6164) -- (3.8967,-1.7033) -- (3.8099,-1.7901) -- (3.7232,-1.8768) -- (3.6366,-1.9634) -- (3.5500,-2.0500) -- cycle;
\end{scope}
\draw[atomLine, line width=0.5pt] (3.2500,-2.3500) rectangle (8.0000,2.4000);
\draw[atomLine!70, line width=0.4pt] (-0.0720,-0.0720) rectangle (1.0680,1.0680);
\draw[atomLine!50, densely dotted] (1.0680,1.0680) -- (3.2500,2.4000) (1.0680,-0.0720) -- (3.2500,-2.3500);
\node[pole100] at (0.0000,0.0000) {};
\node[pole100] at (0.0000,2.4000) {};
\node[pole100] at (2.4000,0.0000) {};
\node[pole110] at (0.0000,0.9941) {};
\node[pole110] at (0.0000,-0.9941) {};
\node[pole110] at (0.9941,0.0000) {};
\node[pole110] at (-0.9941,0.0000) {};
\node[pole110] at (1.6971,1.6971) {};
\node[pole110] at (1.6971,-1.6971) {};
\node[pole111] at (0.8785,0.8785) {};
\node[pole111] at (0.8785,-0.8785) {};
\node[pole111] at (-0.8785,0.8785) {};
\node[pole111] at (-0.8785,-0.8785) {};
\node[pole100] (B100) at (3.5500,-2.0500) {};
\node[pole110] (B110) at (7.6921,-2.0500) {};
\node[pole110] (B101) at (3.5500,2.0921) {};
\node[pole111] (B111) at (7.2103,1.6103) {};
\node[poleH] (B211) at (5.7974,0.1974) {};
\draw[fmRed, line width=1pt, -{Stealth[length=2mm]}] (6.2982,-1.1339) -- (6.2971,-1.1322) -- (6.2959,-1.1305) -- (6.2947,-1.1287) -- (6.2935,-1.1270) -- (6.2924,-1.1253) -- (6.2912,-1.1235) -- (6.2900,-1.1218) -- (6.2888,-1.1201) -- (6.2877,-1.1184) -- (6.2865,-1.1166) -- (6.2853,-1.1149) -- (6.2841,-1.1132) -- (6.2830,-1.1114) -- (6.2818,-1.1097) -- (6.2806,-1.1080) -- (6.2794,-1.1063) -- (6.2783,-1.1045) -- (6.2771,-1.1028) -- (6.2759,-1.1011) -- (6.2747,-1.0994) -- (6.2735,-1.0976) -- (6.2724,-1.0959) -- (6.2712,-1.0942) -- (6.2700,-1.0925) -- (6.2688,-1.0907) -- (6.2677,-1.0890) -- (6.2665,-1.0873) -- (6.2653,-1.0856) -- (6.2641,-1.0838) -- (6.2629,-1.0821) -- (6.2618,-1.0804) -- (6.2606,-1.0787) -- (6.2594,-1.0769) -- (6.2582,-1.0752) -- (6.2571,-1.0735) -- (6.2559,-1.0718) -- (6.2547,-1.0700) -- (6.2535,-1.0683) -- (6.2523,-1.0666) -- (6.2512,-1.0649) -- (6.2500,-1.0631) -- (6.2488,-1.0614) -- (6.2476,-1.0597) -- (6.2464,-1.0580) -- (6.2452,-1.0563) -- (6.2441,-1.0545) -- (6.2429,-1.0528) -- (6.2417,-1.0511) -- (6.2405,-1.0494) -- (6.2393,-1.0477) -- (6.2382,-1.0459) -- (6.2370,-1.0442) -- (6.2358,-1.0425) -- (6.2346,-1.0408) -- (6.2334,-1.0391) -- (6.2323,-1.0373) -- (6.2311,-1.0356) -- (6.2299,-1.0339) -- (6.2287,-1.0322) -- (6.2275,-1.0305) -- (6.2263,-1.0287) -- (6.2252,-1.0270) -- (6.2240,-1.0253) -- (6.2228,-1.0236) -- (6.2216,-1.0219) -- (6.2204,-1.0201) -- (6.2192,-1.0184) -- (6.2180,-1.0167) -- (6.2169,-1.0150) -- (6.2157,-1.0133) -- (6.2145,-1.0116) -- (6.2133,-1.0098) -- (6.2121,-1.0081) -- (6.2109,-1.0064) -- (6.2098,-1.0047) -- (6.2086,-1.0030) -- (6.2074,-1.0013) -- (6.2062,-0.9995) -- (6.2050,-0.9978) -- (6.2038,-0.9961) -- (6.2026,-0.9944) -- (6.2014,-0.9927) -- (6.2003,-0.9910) -- (6.1991,-0.9892) -- (6.1979,-0.9875) -- (6.1967,-0.9858) -- (6.1955,-0.9841) -- (6.1943,-0.9824) -- (6.1931,-0.9807) -- (6.1920,-0.9790) -- (6.1908,-0.9772) -- (6.1896,-0.9755) -- (6.1884,-0.9738) -- (6.1872,-0.9721) -- (6.1860,-0.9704) -- (6.1848,-0.9687) -- (6.1836,-0.9670) -- (6.1824,-0.9652) -- (6.1813,-0.9635) -- (6.1801,-0.9618) -- (6.1789,-0.9601) -- (6.1777,-0.9584) -- (6.1765,-0.9567) -- (6.1753,-0.9550) -- (6.1741,-0.9533) -- (6.1729,-0.9515) -- (6.1717,-0.9498) -- (6.1705,-0.9481) -- (6.1694,-0.9464) -- (6.1682,-0.9447) -- (6.1670,-0.9430) -- (6.1658,-0.9413) -- (6.1646,-0.9396) -- (6.1634,-0.9379) -- (6.1622,-0.9362) -- (6.1610,-0.9344) -- (6.1598,-0.9327) -- (6.1586,-0.9310) -- (6.1574,-0.9293) -- (6.1562,-0.9276) -- (6.1551,-0.9259) -- (6.1539,-0.9242) -- (6.1527,-0.9225) -- (6.1515,-0.9208) -- (6.1503,-0.9191) -- (6.1491,-0.9174) -- (6.1479,-0.9156) -- (6.1467,-0.9139) -- (6.1455,-0.9122) -- (6.1443,-0.9105) -- (6.1431,-0.9088) -- (6.1419,-0.9071) -- (6.1407,-0.9054) -- (6.1395,-0.9037) -- (6.1383,-0.9020) -- (6.1371,-0.9003) -- (6.1359,-0.8986) -- (6.1348,-0.8969) -- (6.1336,-0.8952) -- (6.1324,-0.8935) -- (6.1312,-0.8918) -- (6.1300,-0.8900) -- (6.1288,-0.8883) -- (6.1276,-0.8866) -- (6.1264,-0.8849) -- (6.1252,-0.8832) -- (6.1240,-0.8815) -- (6.1228,-0.8798) -- (6.1216,-0.8781) -- (6.1204,-0.8764) -- (6.1192,-0.8747) -- (6.1180,-0.8730) -- (6.1168,-0.8713) -- (6.1156,-0.8696) -- (6.1144,-0.8679) -- (6.1132,-0.8662) -- (6.1120,-0.8645) -- (6.1108,-0.8628) -- (6.1096,-0.8611) -- (6.1084,-0.8594) -- (6.1072,-0.8577) -- (6.1060,-0.8560) -- (6.1048,-0.8543) -- (6.1036,-0.8526) -- (6.1024,-0.8509) -- (6.1012,-0.8492) -- (6.1000,-0.8475) -- (6.0988,-0.8458) -- (6.0976,-0.8441) -- (6.0964,-0.8424) -- (6.0952,-0.8407) -- (6.0940,-0.8390) -- (6.0928,-0.8373) -- (6.0916,-0.8356) -- (6.0904,-0.8339) -- (6.0892,-0.8322) -- (6.0880,-0.8305) -- (6.0868,-0.8288) -- (6.0856,-0.8271) -- (6.0844,-0.8254) -- (6.0832,-0.8237) -- (6.0820,-0.8220) -- (6.0808,-0.8203) -- (6.0796,-0.8186) -- (6.0784,-0.8169) -- (6.0772,-0.8152) -- (6.0760,-0.8135) -- (6.0748,-0.8118) -- (6.0736,-0.8101) -- (6.0724,-0.8084) -- (6.0712,-0.8067) -- (6.0700,-0.8050) -- (6.0688,-0.8033) -- (6.0676,-0.8016) -- (6.0664,-0.7999) -- (6.0652,-0.7982) -- (6.0639,-0.7965) -- (6.0627,-0.7948) -- (6.0615,-0.7931) -- (6.0603,-0.7914) -- (6.0591,-0.7897) -- (6.0579,-0.7880) -- (6.0567,-0.7863) -- (6.0555,-0.7846) -- (6.0543,-0.7829) -- (6.0531,-0.7812) -- (6.0519,-0.7795) -- (6.0507,-0.7778) -- (6.0495,-0.7761) -- (6.0483,-0.7744) -- (6.0471,-0.7728) -- (6.0459,-0.7711) -- (6.0446,-0.7694) -- (6.0434,-0.7677) -- (6.0422,-0.7660) -- (6.0410,-0.7643) -- (6.0398,-0.7626) -- (6.0386,-0.7609) -- (6.0374,-0.7592) -- (6.0362,-0.7575) -- (6.0350,-0.7558) -- (6.0338,-0.7541) -- (6.0326,-0.7524) -- (6.0314,-0.7507) -- (6.0301,-0.7490) -- (6.0289,-0.7474) -- (6.0277,-0.7457) -- (6.0265,-0.7440) -- (6.0253,-0.7423) -- (6.0241,-0.7406) -- (6.0229,-0.7389) -- (6.0217,-0.7372) -- (6.0205,-0.7355) -- (6.0193,-0.7338) -- (6.0180,-0.7321) -- (6.0168,-0.7304) -- (6.0156,-0.7287) -- (6.0144,-0.7271) -- (6.0132,-0.7254) -- (6.0120,-0.7237) -- (6.0108,-0.7220) -- (6.0096,-0.7203) -- (6.0084,-0.7186) -- (6.0071,-0.7169) -- (6.0059,-0.7152) -- (6.0047,-0.7135) -- (6.0035,-0.7118) -- (6.0023,-0.7102) -- (6.0011,-0.7085) -- (5.9999,-0.7068) -- (5.9986,-0.7051) -- (5.9974,-0.7034) -- (5.9962,-0.7017) -- (5.9950,-0.7000) -- (5.9938,-0.6983) -- (5.9926,-0.6967) -- (5.9914,-0.6950) -- (5.9901,-0.6933) -- (5.9889,-0.6916) -- (5.9877,-0.6899) -- (5.9865,-0.6882) -- (5.9853,-0.6865) -- (5.9841,-0.6848) -- (5.9829,-0.6832) -- (5.9816,-0.6815) -- (5.9804,-0.6798) -- (5.9792,-0.6781) -- (5.9780,-0.6764) -- (5.9768,-0.6747) -- (5.9756,-0.6730) -- (5.9743,-0.6714) -- (5.9731,-0.6697) -- (5.9719,-0.6680) -- (5.9707,-0.6663) -- (5.9695,-0.6646) -- (5.9683,-0.6629) -- (5.9670,-0.6612) -- (5.9658,-0.6596) -- (5.9646,-0.6579) -- (5.9634,-0.6562) -- (5.9622,-0.6545) -- (5.9609,-0.6528) -- (5.9597,-0.6511) -- (5.9585,-0.6495) -- (5.9573,-0.6478) -- (5.9561,-0.6461) -- (5.9548,-0.6444) -- (5.9536,-0.6427) -- (5.9524,-0.6410) -- (5.9512,-0.6394) -- (5.9500,-0.6377) -- (5.9487,-0.6360) -- (5.9475,-0.6343) -- (5.9463,-0.6326) -- (5.9451,-0.6310) -- (5.9439,-0.6293) -- (5.9426,-0.6276) -- (5.9414,-0.6259) -- (5.9402,-0.6242) -- (5.9390,-0.6225) -- (5.9378,-0.6209) -- (5.9365,-0.6192) -- (5.9353,-0.6175) -- (5.9341,-0.6158) -- (5.9329,-0.6141) -- (5.9316,-0.6125) -- (5.9304,-0.6108) -- (5.9292,-0.6091) -- (5.9280,-0.6074) -- (5.9268,-0.6057) -- (5.9255,-0.6041) -- (5.9243,-0.6024) -- (5.9231,-0.6007) -- (5.9219,-0.5990) -- (5.9206,-0.5973) -- (5.9194,-0.5957) -- (5.9182,-0.5940) -- (5.9170,-0.5923) -- (5.9157,-0.5906) -- (5.9145,-0.5890) -- (5.9133,-0.5873) -- (5.9121,-0.5856) -- (5.9108,-0.5839) -- (5.9096,-0.5822) -- (5.9084,-0.5806) -- (5.9072,-0.5789) -- (5.9059,-0.5772) -- (5.9047,-0.5755) -- (5.9035,-0.5739) -- (5.9022,-0.5722) -- (5.9010,-0.5705) -- (5.8998,-0.5688) -- (5.8986,-0.5672) -- (5.8973,-0.5655) -- (5.8961,-0.5638) -- (5.8949,-0.5621) -- (5.8937,-0.5604) -- (5.8924,-0.5588) -- (5.8912,-0.5571) -- (5.8900,-0.5554) -- (5.8887,-0.5537) -- (5.8875,-0.5521) -- (5.8863,-0.5504) -- (5.8851,-0.5487) -- (5.8838,-0.5470) -- (5.8826,-0.5454) -- (5.8814,-0.5437) -- (5.8801,-0.5420) -- (5.8789,-0.5404) -- (5.8777,-0.5387) -- (5.8764,-0.5370) -- (5.8752,-0.5353) -- (5.8740,-0.5337) -- (5.8728,-0.5320) -- (5.8715,-0.5303) -- (5.8703,-0.5286) -- (5.8691,-0.5270) -- (5.8678,-0.5253) -- (5.8666,-0.5236) -- (5.8654,-0.5219) -- (5.8641,-0.5203) -- (5.8629,-0.5186) -- (5.8617,-0.5169) -- (5.8604,-0.5153) -- (5.8592,-0.5136) -- (5.8580,-0.5119) -- (5.8567,-0.5102) -- (5.8555,-0.5086) -- (5.8543,-0.5069) -- (5.8530,-0.5052) -- (5.8518,-0.5036) -- (5.8506,-0.5019) -- (5.8493,-0.5002) -- (5.8481,-0.4985) -- (5.8469,-0.4969) -- (5.8456,-0.4952) -- (5.8444,-0.4935) -- (5.8432,-0.4919) -- (5.8419,-0.4902) -- (5.8407,-0.4885) -- (5.8395,-0.4869) -- (5.8382,-0.4852) -- (5.8370,-0.4835) -- (5.8357,-0.4819) -- (5.8345,-0.4802) -- (5.8333,-0.4785) -- (5.8320,-0.4768) -- (5.8308,-0.4752) -- (5.8296,-0.4735) -- (5.8283,-0.4718) -- (5.8271,-0.4702) -- (5.8259,-0.4685) -- (5.8246,-0.4668) -- (5.8234,-0.4652) -- (5.8221,-0.4635) -- (5.8209,-0.4618) -- (5.8197,-0.4602) -- (5.8184,-0.4585) -- (5.8172,-0.4568) -- (5.8159,-0.4552) -- (5.8147,-0.4535) -- (5.8135,-0.4518) -- (5.8122,-0.4502) -- (5.8110,-0.4485) -- (5.8098,-0.4468) -- (5.8085,-0.4452) -- (5.8073,-0.4435) -- (5.8060,-0.4418) -- (5.8048,-0.4402) -- (5.8035,-0.4385) -- (5.8023,-0.4368) -- (5.8011,-0.4352) -- (5.7998,-0.4335) -- (5.7986,-0.4318) -- (5.7973,-0.4302) -- (5.7961,-0.4285) -- (5.7949,-0.4268) -- (5.7936,-0.4252) -- (5.7924,-0.4235) -- (5.7911,-0.4219) -- (5.7899,-0.4202) -- (5.7886,-0.4185) -- (5.7874,-0.4169) -- (5.7862,-0.4152) -- (5.7849,-0.4135) -- (5.7837,-0.4119) -- (5.7824,-0.4102) -- (5.7812,-0.4085) -- (5.7799,-0.4069) -- (5.7787,-0.4052) -- (5.7775,-0.4036) -- (5.7762,-0.4019) -- (5.7750,-0.4002) -- (5.7737,-0.3986) -- (5.7725,-0.3969) -- (5.7712,-0.3952) -- (5.7700,-0.3936) -- (5.7687,-0.3919) -- (5.7675,-0.3903) -- (5.7663,-0.3886) -- (5.7650,-0.3869) -- (5.7638,-0.3853) -- (5.7625,-0.3836) -- (5.7613,-0.3820) -- (5.7600,-0.3803) -- (5.7588,-0.3786) -- (5.7575,-0.3770) -- (5.7563,-0.3753) -- (5.7550,-0.3737) -- (5.7538,-0.3720) -- (5.7525,-0.3703) -- (5.7513,-0.3687) -- (5.7500,-0.3670) -- (5.7488,-0.3654) -- (5.7475,-0.3637) -- (5.7463,-0.3620) -- (5.7450,-0.3604) -- (5.7438,-0.3587) -- (5.7425,-0.3571) -- (5.7413,-0.3554) -- (5.7400,-0.3537) -- (5.7388,-0.3521) -- (5.7375,-0.3504) -- (5.7363,-0.3488) -- (5.7350,-0.3471) -- (5.7338,-0.3454) -- (5.7325,-0.3438) -- (5.7313,-0.3421) -- (5.7300,-0.3405) -- (5.7288,-0.3388) -- (5.7275,-0.3372) -- (5.7263,-0.3355) -- (5.7250,-0.3338) -- (5.7238,-0.3322) -- (5.7225,-0.3305) -- (5.7213,-0.3289) -- (5.7200,-0.3272) -- (5.7188,-0.3256) -- (5.7175,-0.3239) -- (5.7163,-0.3222) -- (5.7150,-0.3206) -- (5.7138,-0.3189) -- (5.7125,-0.3173) -- (5.7113,-0.3156) -- (5.7100,-0.3140) -- (5.7087,-0.3123) -- (5.7075,-0.3107) -- (5.7062,-0.3090) -- (5.7050,-0.3073) -- (5.7037,-0.3057) -- (5.7025,-0.3040) -- (5.7012,-0.3024) -- (5.7000,-0.3007) -- (5.6987,-0.2991) -- (5.6975,-0.2974) -- (5.6962,-0.2958) -- (5.6949,-0.2941) -- (5.6937,-0.2925) -- (5.6924,-0.2908) -- (5.6912,-0.2892) -- (5.6899,-0.2875) -- (5.6887,-0.2858) -- (5.6874,-0.2842) -- (5.6861,-0.2825) -- (5.6849,-0.2809) -- (5.6836,-0.2792) -- (5.6824,-0.2776) -- (5.6811,-0.2759) -- (5.6799,-0.2743) -- (5.6786,-0.2726) -- (5.6773,-0.2710) -- (5.6761,-0.2693) -- (5.6748,-0.2677) -- (5.6736,-0.2660) -- (5.6723,-0.2644) -- (5.6710,-0.2627) -- (5.6698,-0.2611) -- (5.6685,-0.2594) -- (5.6673,-0.2578) -- (5.6660,-0.2561) -- (5.6647,-0.2545) -- (5.6635,-0.2528) -- (5.6622,-0.2512) -- (5.6610,-0.2495) -- (5.6597,-0.2479) -- (5.6584,-0.2462) -- (5.6572,-0.2446) -- (5.6559,-0.2429) -- (5.6546,-0.2413) -- (5.6534,-0.2396) -- (5.6521,-0.2380) -- (5.6509,-0.2363) -- (5.6496,-0.2347) -- (5.6483,-0.2330) -- (5.6471,-0.2314) -- (5.6458,-0.2297) -- (5.6445,-0.2281) -- (5.6433,-0.2264) -- (5.6420,-0.2248) -- (5.6408,-0.2231) -- (5.6395,-0.2215) -- (5.6382,-0.2198) -- (5.6370,-0.2182) -- (5.6357,-0.2165) -- (5.6344,-0.2149) -- (5.6332,-0.2132) -- (5.6319,-0.2116) -- (5.6306,-0.2099) -- (5.6294,-0.2083) -- (5.6281,-0.2066) -- (5.6268,-0.2050) -- (5.6256,-0.2033) -- (5.6243,-0.2017) -- (5.6230,-0.2000) -- (5.6218,-0.1984) -- (5.6205,-0.1967) -- (5.6192,-0.1951) -- (5.6180,-0.1935) -- (5.6167,-0.1918) -- (5.6154,-0.1902) -- (5.6142,-0.1885) -- (5.6129,-0.1869) -- (5.6116,-0.1852) -- (5.6104,-0.1836) -- (5.6091,-0.1819) -- (5.6078,-0.1803) -- (5.6066,-0.1786) -- (5.6053,-0.1770) -- (5.6040,-0.1754) -- (5.6027,-0.1737) -- (5.6015,-0.1721) -- (5.6002,-0.1704) -- (5.5989,-0.1688) -- (5.5977,-0.1671) -- (5.5964,-0.1655) -- (5.5951,-0.1638) -- (5.5939,-0.1622) -- (5.5926,-0.1606) -- (5.5913,-0.1589) -- (5.5900,-0.1573) -- (5.5888,-0.1556) -- (5.5875,-0.1540) -- (5.5862,-0.1523) -- (5.5850,-0.1507) -- (5.5837,-0.1490) -- (5.5824,-0.1474) -- (5.5811,-0.1458) -- (5.5799,-0.1441) -- (5.5786,-0.1425) -- (5.5773,-0.1408) -- (5.5760,-0.1392) -- (5.5748,-0.1375) -- (5.5735,-0.1359) -- (5.5722,-0.1343) -- (5.5709,-0.1326) -- (5.5697,-0.1310) -- (5.5684,-0.1293) -- (5.5671,-0.1277) -- (5.5658,-0.1261) -- (5.5646,-0.1244) -- (5.5633,-0.1228) -- (5.5620,-0.1211) -- (5.5607,-0.1195) -- (5.5595,-0.1179) -- (5.5582,-0.1162) -- (5.5569,-0.1146) -- (5.5556,-0.1129) -- (5.5544,-0.1113) -- (5.5531,-0.1096) -- (5.5518,-0.1080) -- (5.5505,-0.1064) -- (5.5493,-0.1047) -- (5.5480,-0.1031) -- (5.5467,-0.1014) -- (5.5454,-0.0998) -- (5.5441,-0.0982) -- (5.5429,-0.0965) -- (5.5416,-0.0949) -- (5.5403,-0.0933) -- (5.5390,-0.0916) -- (5.5377,-0.0900) -- (5.5365,-0.0883) -- (5.5352,-0.0867) -- (5.5339,-0.0851) -- (5.5326,-0.0834) -- (5.5313,-0.0818) -- (5.5301,-0.0801) -- (5.5288,-0.0785) -- (5.5275,-0.0769) -- (5.5262,-0.0752) -- (5.5249,-0.0736);
\draw[fmRed, line width=0.8pt, dashed] (5.5249,-0.0736) -- (5.5207,-0.0681) -- (5.5164,-0.0627) -- (5.5121,-0.0572) -- (5.5079,-0.0518) -- (5.5036,-0.0463) -- (5.4993,-0.0409) -- (5.4950,-0.0354) -- (5.4907,-0.0300) -- (5.4864,-0.0245) -- (5.4821,-0.0191) -- (5.4778,-0.0136) -- (5.4735,-0.0081) -- (5.4692,-0.0027) -- (5.4649,0.0028) -- (5.4606,0.0082) -- (5.4563,0.0137) -- (5.4519,0.0192) -- (5.4476,0.0246) -- (5.4433,0.0301) -- (5.4389,0.0355) -- (5.4346,0.0410) -- (5.4303,0.0465) -- (5.4259,0.0519) -- (5.4216,0.0574) -- (5.4172,0.0628) -- (5.4129,0.0683) -- (5.4085,0.0738) -- (5.4041,0.0792) -- (5.3998,0.0847);
\draw[fmGreen!70!black, line width=1.2pt, -{Stealth[length=1.8mm]}] (5.5249,-0.0736) -- (5.5336,-0.0650) -- (5.5423,-0.0564) -- (5.5509,-0.0478) -- (5.5596,-0.0391) -- (5.5683,-0.0305) -- (5.5769,-0.0219) -- (5.5856,-0.0132) -- (5.5943,-0.0046) -- (5.6030,0.0041) -- (5.6117,0.0127) -- (5.6204,0.0214) -- (5.6291,0.0300) -- (5.6379,0.0387) -- (5.6466,0.0474) -- (5.6553,0.0561) -- (5.6640,0.0648) -- (5.6728,0.0735) -- (5.6815,0.0822) -- (5.6903,0.0909) -- (5.6991,0.0996) -- (5.7078,0.1083) -- (5.7166,0.1170) -- (5.7254,0.1258) -- (5.7342,0.1345) -- (5.7430,0.1432) -- (5.7518,0.1520) -- (5.7606,0.1608) -- (5.7694,0.1695) -- (5.7782,0.1783);
\coordinate (F) at (6.2982,-1.1339);
\coordinate (FA) at (0.6596,0.2199);
\coordinate (Fcross) at (5.5249,-0.0736);
\coordinate (nPrim) at (0.8785,-0.8785);
\coordinate (dPrim) at (0.0000,0.9941);
\coordinate (nConj) at (-0.8785,0.8785);
\coordinate (dConj) at (0.9941,0.0000);
\coordinate (nCrit) at (0.8785,0.8785);
\coordinate (dCrit) at (0.0000,-0.9941);
\coordinate (nCross) at (-0.8785,-0.8785);
\coordinate (A100) at (0.0000,0.0000);
\coordinate (A010) at (2.4000,0.0000);
\coordinate (A001) at (0.0000,2.4000);
\coordinate (A110) at (0.9941,0.0000);

  % (а) вся проекция
  \fill[fmRed] (FA) circle (0.035);
  \node[font=\scriptsize, anchor=north east, inner sep=1pt] at (A100) {$[100]$};
  \node[font=\scriptsize, anchor=west] at (A010) {$[010]$};
  \node[font=\scriptsize, anchor=south] at (A001) {$[001]$};
  \node[ring=fmRed] at (dPrim) {};
  \node[font=\scriptsize, fmRed, anchor=south east, inner sep=1pt] at (dPrim) {$[101]$};
  \node[ring=fmRed] at (nPrim) {};
  \node[font=\scriptsize, fmRed, anchor=north west, inner sep=2pt] at (nPrim) {$(11\br1)$};
  \node[ring=fmGreen] at (nConj) {};
  \node[font=\scriptsize, fmGreen, anchor=south east, inner sep=2pt] at (nConj) {$(1\br11)$};
  \node[ring=fmPurple] at (nCrit) {};
  \node[font=\scriptsize, fmPurple, anchor=south west, inner sep=2pt] at (nCrit) {$(111)$};
  \node[ring=fmTeal] at (nCross) {};
  \node[font=\scriptsize, fmTeal, anchor=north east, inner sep=2pt] at (nCross) {$(1\br1\br1)$};
  \node[font=\small] at (0,-2.75) {\textit{а}};
  % (б) окрестность стандартного треугольника
  \node[font=\scriptsize, anchor=north east, inner sep=1pt] at (B100) {$[100]$};
  \node[font=\scriptsize, anchor=north, inner sep=2pt] at (B110) {$[110]$};
  \node[font=\scriptsize, anchor=south west, inner sep=1pt] at (B111) {$[111]$};
  \node[font=\scriptsize, anchor=west, inner sep=2pt] at (B101) {$[101]$};
  \node[font=\scriptsize, fmGreen!70!black, anchor=north west, inner sep=1pt] at (B211) {$[211]$};
  \fill[fmRed] (F) circle (0.05) node[below right=-1pt, font=\small, fmRed]{$F$};
  \node[font=\small] at (5.6,-2.75) {\textit{б}};
\end{tikzpicture}
\caption{Системы скольжения ГЦК-кристалла на стереографической проекции
вдоль $[100]$: \textit{а}~— вся проекция; кружками отмечены полюса
плоскостей скольжения первичной $(11\br1)$ (направление $[101]$),
сопряжённой $(1\br11)$ (направление $[110]$), критической $(111)$
(направление $[10\br1]$) и поперечной $(1\br1\br1)$ (направление $[101]$)
систем; квадраты~— $\nfam{100}$, овалы~— $\nfam{110}$, треугольники~—
$\nfam{111}$. \textit{б}~— стандартный треугольник $[100]$--$[110]$--$[111]$
крупно: при одиночном скольжении ось $F\parallel[531]$ поворачивается к
$[101]$ (красная линия), немного «перескакивает» сторону $[100]$--$[111]$
(пунктир), после чего при двойном скольжении движется к $[211]$
(зелёная стрелка)}
\label{fig:slip-stereo}
\end{figure}

Системам скольжения дают названия по их роли в деформации
(рис.~\ref{fig:slip-stereo}):
\begin{description}
\item[Первичная] $(11\br1)[101]$~— с наибольшим фактором Шмида для любой
оси внутри треугольника; включается первой.
\item[Сопряжённая] $(1\br11)[110]$~— зеркальное отражение первичной в
стороне $[100]$--$[111]$. При растяжении ось $F$ поворачивается к
направлению скольжения $[101]$ и достигает этой стороны; там факторы
Шмида первичной и сопряжённой систем равны, и начинается \emph{двойное
скольжение}: ось движется вдоль стороны к $[211]=[101]+[110]$. Обычно
ось немного «перескакивает» границу в соседний треугольник
$[100]$--$[101]$--$[111]$ (пунктир), потому что сопряжённая система
упрочнена дислокациями первичной.
\item[Критическая] $(111)[10\br1]$~— отражение первичной в стороне
$[100]$--$[110]$. При \emph{сжатии} ось движется к нормали плоскости
скольжения $(11\br1)$, попадает в треугольник $[100]$--$[110]$--$[11\br1]$
и включает критическую систему.
\item[Поперечная] $(1\br1\br1)[101]$~— имеет то же направление
скольжения (тот же вектор Бюргерса), что и первичная; в эту плоскость
переходят винтовые дислокации при поперечном скольжении (гл.~2,
п.~2.2.2).
\end{description}

\subsubsection{Скольжение в ОЦК-металлах}

В ОЦК-металлах ($\alpha$-Fe, Mo, W, Nb, Ta, $\beta$-Ti) направление
скольжения~— плотноупакованное $\nfam{111}$, а плоскостью скольжения в
зависимости от металла и температуры служат $\plfam{110}$, $\plfam{112}$
или $\plfam{123}$~— все они содержат направление $\nfam{111}$.
Систем: $\plfam{110}\nfam{111}$~— 12, $\plfam{112}\nfam{111}$~— 12,
$\plfam{123}\nfam{111}$~— 24. Часто наблюдается и
\term{некристаллографическое} («карандашное») скольжение: направление
$\nfam{111}$, а плоскость~— любая, проходящая через него, с наибольшим
касательным напряжением ($\chi=90^\circ-\lambda$). Следы такого скольжения
волнистые.

Как правило, скольжение \emph{неполярно}: сдвиг вдоль $+\vb m$ и $-\vb m$
кристаллографически эквивалентен. Но если плоскость скольжения~— не
плоскость симметрии и в ней нет оси 2, переводящей $\vb m$ в $-\vb m$,
то направление скольжения~— \emph{полярный} вектор в этой плоскости, и
сопротивление сдвигу «туда» и «обратно» различно (несимметричный
потенциальный рельеф решётки). Типичный пример~— системы
$\plfam{112}\nfam{111}$ и $\plfam{123}\nfam{111}$ в ОЦК-металлах: в
одном направлении сдвиг по $\plfam{112}$ идёт «в двойниковом», в обратном~—
«в антидвойниковом» смысле. В системе $\plfam{110}\nfam{111}$ полярности
нет: ось 2 вдоль нормали $\nfam{110}$ переводит $[111]$ в $[\br1\br1\br1]$.

\subsubsection{Скольжение в ГПУ-металлах}

В ГПУ-металлах направление скольжения~— обычно плотноупакованное
$\nfam{11\br20}$ (то же семейство, что $\nfam{2\br1\br10}$), а плоскости
скольжения разнообразны (рис.~\ref{fig:hcp-slip}): \term{базисная}
$(0001)$, \term{призматические} $\plfam{10\br10}$ и \term{пирамидальные}
$\plfam{10\br11}$, $\plfam{10\br12}$, $\plfam{11\br21}$,
$\plfam{11\br22}$\ispr{В оригинале пирамидальные плоскости названы
плоскостями I--IV рода в порядке перечисления. В литературе чаще
различают пирамиды I рода $\plfam{h0\br hl}$ (сюда относятся
$\plfam{10\br11}$ и $\plfam{10\br12}$) и II рода $\plfam{hh\,\overline{2h}\,l}$
($\plfam{11\br21}$, $\plfam{11\br22}$).}.

\begin{figure}[ht]
\centering
\begin{tikzpicture}[cubeproj, scale=0.95]
  % создано tools/gen_plastic.py: плоскости скольжения ГПУ (cubeproj; c/a = 1.633)
\begin{scope}[shift={(0.0,0,0.0)}]
\draw[atomLine!45, dashed] (0.5000,0.8660,0.0000) -- (-0.5000,0.8660,0.0000);
\draw[atomLine!45, dashed] (-0.5000,0.8660,0.0000) -- (-1.0000,0.0000,0.0000);
\draw[atomLine!45, dashed] (-0.5000,0.8660,0.0000) -- (-0.5000,0.8660,1.6330);
\fill[fmRed!35, opacity=0.85] (-1.0000,0.0000,0.8165) -- (-0.5000,-0.8660,0.8165) -- (0.5000,-0.8660,0.8165) -- (1.0000,0.0000,0.8165) -- (0.5000,0.8660,0.8165) -- (-0.5000,0.8660,0.8165) -- cycle;
\draw[fmRed!80!black, line width=0.7pt] (-1.0000,0.0000,0.8165) -- (-0.5000,-0.8660,0.8165) -- (0.5000,-0.8660,0.8165) -- (1.0000,0.0000,0.8165) -- (0.5000,0.8660,0.8165) -- (-0.5000,0.8660,0.8165) -- cycle;
\draw[atomLine, line width=0.5pt] (1.0000,0.0000,0.0000) -- (0.5000,0.8660,0.0000);
\draw[atomLine, line width=0.5pt] (-1.0000,0.0000,0.0000) -- (-0.5000,-0.8660,0.0000);
\draw[atomLine, line width=0.5pt] (-0.5000,-0.8660,0.0000) -- (0.5000,-0.8660,0.0000);
\draw[atomLine, line width=0.5pt] (0.5000,-0.8660,0.0000) -- (1.0000,0.0000,0.0000);
\draw[atomLine, line width=0.5pt] (1.0000,0.0000,1.6330) -- (0.5000,0.8660,1.6330);
\draw[atomLine, line width=0.5pt] (0.5000,0.8660,1.6330) -- (-0.5000,0.8660,1.6330);
\draw[atomLine, line width=0.5pt] (-0.5000,0.8660,1.6330) -- (-1.0000,0.0000,1.6330);
\draw[atomLine, line width=0.5pt] (-1.0000,0.0000,1.6330) -- (-0.5000,-0.8660,1.6330);
\draw[atomLine, line width=0.5pt] (-0.5000,-0.8660,1.6330) -- (0.5000,-0.8660,1.6330);
\draw[atomLine, line width=0.5pt] (0.5000,-0.8660,1.6330) -- (1.0000,0.0000,1.6330);
\draw[atomLine, line width=0.5pt] (1.0000,0.0000,0.0000) -- (1.0000,0.0000,1.6330);
\draw[atomLine, line width=0.5pt] (0.5000,0.8660,0.0000) -- (0.5000,0.8660,1.6330);
\draw[atomLine, line width=0.5pt] (-1.0000,0.0000,0.0000) -- (-1.0000,0.0000,1.6330);
\draw[atomLine, line width=0.5pt] (-0.5000,-0.8660,0.0000) -- (-0.5000,-0.8660,1.6330);
\draw[atomLine, line width=0.5pt] (0.5000,-0.8660,0.0000) -- (0.5000,-0.8660,1.6330);
\node[font=\small, anchor=north] at (0,-1.05,0) {$1$: $(0001)$};
\end{scope}
\begin{scope}[shift={(3.5,0,0.0)}]
\draw[atomLine!45, dashed] (0.5000,0.8660,0.0000) -- (-0.5000,0.8660,0.0000);
\draw[atomLine!45, dashed] (-0.5000,0.8660,0.0000) -- (-1.0000,0.0000,0.0000);
\draw[atomLine!45, dashed] (-0.5000,0.8660,0.0000) -- (-0.5000,0.8660,1.6330);
\fill[fmRed!35, opacity=0.85] (-1.0000,0.0000,1.6330) -- (1.0000,0.0000,1.6330) -- (1.0000,0.0000,0.0000) -- (-1.0000,0.0000,0.0000) -- cycle;
\draw[fmRed!80!black, line width=0.7pt] (-1.0000,0.0000,1.6330) -- (1.0000,0.0000,1.6330) -- (1.0000,0.0000,0.0000) -- (-1.0000,0.0000,0.0000) -- cycle;
\draw[atomLine, line width=0.5pt] (1.0000,0.0000,0.0000) -- (0.5000,0.8660,0.0000);
\draw[atomLine, line width=0.5pt] (-1.0000,0.0000,0.0000) -- (-0.5000,-0.8660,0.0000);
\draw[atomLine, line width=0.5pt] (-0.5000,-0.8660,0.0000) -- (0.5000,-0.8660,0.0000);
\draw[atomLine, line width=0.5pt] (0.5000,-0.8660,0.0000) -- (1.0000,0.0000,0.0000);
\draw[atomLine, line width=0.5pt] (1.0000,0.0000,1.6330) -- (0.5000,0.8660,1.6330);
\draw[atomLine, line width=0.5pt] (0.5000,0.8660,1.6330) -- (-0.5000,0.8660,1.6330);
\draw[atomLine, line width=0.5pt] (-0.5000,0.8660,1.6330) -- (-1.0000,0.0000,1.6330);
\draw[atomLine, line width=0.5pt] (-1.0000,0.0000,1.6330) -- (-0.5000,-0.8660,1.6330);
\draw[atomLine, line width=0.5pt] (-0.5000,-0.8660,1.6330) -- (0.5000,-0.8660,1.6330);
\draw[atomLine, line width=0.5pt] (0.5000,-0.8660,1.6330) -- (1.0000,0.0000,1.6330);
\draw[atomLine, line width=0.5pt] (1.0000,0.0000,0.0000) -- (1.0000,0.0000,1.6330);
\draw[atomLine, line width=0.5pt] (0.5000,0.8660,0.0000) -- (0.5000,0.8660,1.6330);
\draw[atomLine, line width=0.5pt] (-1.0000,0.0000,0.0000) -- (-1.0000,0.0000,1.6330);
\draw[atomLine, line width=0.5pt] (-0.5000,-0.8660,0.0000) -- (-0.5000,-0.8660,1.6330);
\draw[atomLine, line width=0.5pt] (0.5000,-0.8660,0.0000) -- (0.5000,-0.8660,1.6330);
\node[font=\small, anchor=north] at (0,-1.05,0) {$2$: $(01\br10)$};
\end{scope}
\begin{scope}[shift={(7.0,0,0.0)}]
\draw[atomLine!45, dashed] (0.5000,0.8660,0.0000) -- (-0.5000,0.8660,0.0000);
\draw[atomLine!45, dashed] (-0.5000,0.8660,0.0000) -- (-1.0000,0.0000,0.0000);
\draw[atomLine!45, dashed] (-0.5000,0.8660,0.0000) -- (-0.5000,0.8660,1.6330);
\fill[fmRed!35, opacity=0.85] (-1.0000,0.0000,1.6330) -- (-0.5000,-0.8660,0.0000) -- (0.5000,-0.8660,0.0000) -- (1.0000,0.0000,1.6330) -- cycle;
\draw[fmRed!80!black, line width=0.7pt] (-1.0000,0.0000,1.6330) -- (-0.5000,-0.8660,0.0000) -- (0.5000,-0.8660,0.0000) -- (1.0000,0.0000,1.6330) -- cycle;
\draw[atomLine, line width=0.5pt] (1.0000,0.0000,0.0000) -- (0.5000,0.8660,0.0000);
\draw[atomLine, line width=0.5pt] (-1.0000,0.0000,0.0000) -- (-0.5000,-0.8660,0.0000);
\draw[atomLine, line width=0.5pt] (-0.5000,-0.8660,0.0000) -- (0.5000,-0.8660,0.0000);
\draw[atomLine, line width=0.5pt] (0.5000,-0.8660,0.0000) -- (1.0000,0.0000,0.0000);
\draw[atomLine, line width=0.5pt] (1.0000,0.0000,1.6330) -- (0.5000,0.8660,1.6330);
\draw[atomLine, line width=0.5pt] (0.5000,0.8660,1.6330) -- (-0.5000,0.8660,1.6330);
\draw[atomLine, line width=0.5pt] (-0.5000,0.8660,1.6330) -- (-1.0000,0.0000,1.6330);
\draw[atomLine, line width=0.5pt] (-1.0000,0.0000,1.6330) -- (-0.5000,-0.8660,1.6330);
\draw[atomLine, line width=0.5pt] (-0.5000,-0.8660,1.6330) -- (0.5000,-0.8660,1.6330);
\draw[atomLine, line width=0.5pt] (0.5000,-0.8660,1.6330) -- (1.0000,0.0000,1.6330);
\draw[atomLine, line width=0.5pt] (1.0000,0.0000,0.0000) -- (1.0000,0.0000,1.6330);
\draw[atomLine, line width=0.5pt] (0.5000,0.8660,0.0000) -- (0.5000,0.8660,1.6330);
\draw[atomLine, line width=0.5pt] (-1.0000,0.0000,0.0000) -- (-1.0000,0.0000,1.6330);
\draw[atomLine, line width=0.5pt] (-0.5000,-0.8660,0.0000) -- (-0.5000,-0.8660,1.6330);
\draw[atomLine, line width=0.5pt] (0.5000,-0.8660,0.0000) -- (0.5000,-0.8660,1.6330);
\node[font=\small, anchor=north] at (0,-1.05,0) {$3$: $(0\br111)$};
\end{scope}
\begin{scope}[shift={(0.0,0,-3.0)}]
\draw[atomLine!45, dashed] (0.5000,0.8660,0.0000) -- (-0.5000,0.8660,0.0000);
\draw[atomLine!45, dashed] (-0.5000,0.8660,0.0000) -- (-1.0000,0.0000,0.0000);
\draw[atomLine!45, dashed] (-0.5000,0.8660,0.0000) -- (-0.5000,0.8660,1.6330);
\fill[fmRed!35, opacity=0.85] (0.5000,0.8660,1.6330) -- (-0.5000,0.8660,1.6330) -- (-1.0000,0.0000,0.8165) -- (-0.5000,-0.8660,0.0000) -- (0.5000,-0.8660,0.0000) -- (1.0000,0.0000,0.8165) -- cycle;
\draw[fmRed!80!black, line width=0.7pt] (0.5000,0.8660,1.6330) -- (-0.5000,0.8660,1.6330) -- (-1.0000,0.0000,0.8165) -- (-0.5000,-0.8660,0.0000) -- (0.5000,-0.8660,0.0000) -- (1.0000,0.0000,0.8165) -- cycle;
\draw[atomLine, line width=0.5pt] (1.0000,0.0000,0.0000) -- (0.5000,0.8660,0.0000);
\draw[atomLine, line width=0.5pt] (-1.0000,0.0000,0.0000) -- (-0.5000,-0.8660,0.0000);
\draw[atomLine, line width=0.5pt] (-0.5000,-0.8660,0.0000) -- (0.5000,-0.8660,0.0000);
\draw[atomLine, line width=0.5pt] (0.5000,-0.8660,0.0000) -- (1.0000,0.0000,0.0000);
\draw[atomLine, line width=0.5pt] (1.0000,0.0000,1.6330) -- (0.5000,0.8660,1.6330);
\draw[atomLine, line width=0.5pt] (0.5000,0.8660,1.6330) -- (-0.5000,0.8660,1.6330);
\draw[atomLine, line width=0.5pt] (-0.5000,0.8660,1.6330) -- (-1.0000,0.0000,1.6330);
\draw[atomLine, line width=0.5pt] (-1.0000,0.0000,1.6330) -- (-0.5000,-0.8660,1.6330);
\draw[atomLine, line width=0.5pt] (-0.5000,-0.8660,1.6330) -- (0.5000,-0.8660,1.6330);
\draw[atomLine, line width=0.5pt] (0.5000,-0.8660,1.6330) -- (1.0000,0.0000,1.6330);
\draw[atomLine, line width=0.5pt] (1.0000,0.0000,0.0000) -- (1.0000,0.0000,1.6330);
\draw[atomLine, line width=0.5pt] (0.5000,0.8660,0.0000) -- (0.5000,0.8660,1.6330);
\draw[atomLine, line width=0.5pt] (-1.0000,0.0000,0.0000) -- (-1.0000,0.0000,1.6330);
\draw[atomLine, line width=0.5pt] (-0.5000,-0.8660,0.0000) -- (-0.5000,-0.8660,1.6330);
\draw[atomLine, line width=0.5pt] (0.5000,-0.8660,0.0000) -- (0.5000,-0.8660,1.6330);
\node[font=\small, anchor=north] at (0,-1.05,0) {$4$: $(0\br112)$};
\end{scope}
\begin{scope}[shift={(3.5,0,-3.0)}]
\draw[atomLine!45, dashed] (0.5000,0.8660,0.0000) -- (-0.5000,0.8660,0.0000);
\draw[atomLine!45, dashed] (-0.5000,0.8660,0.0000) -- (-1.0000,0.0000,0.0000);
\draw[atomLine!45, dashed] (-0.5000,0.8660,0.0000) -- (-0.5000,0.8660,1.6330);
\fill[fmRed!35, opacity=0.85] (-0.7500,0.4330,1.6330) -- (0.7500,-0.4330,1.6330) -- (1.0000,0.0000,0.0000) -- (-0.5000,0.8660,0.0000) -- cycle;
\draw[fmRed!80!black, line width=0.7pt] (-0.7500,0.4330,1.6330) -- (0.7500,-0.4330,1.6330) -- (1.0000,0.0000,0.0000) -- (-0.5000,0.8660,0.0000) -- cycle;
\draw[atomLine, line width=0.5pt] (1.0000,0.0000,0.0000) -- (0.5000,0.8660,0.0000);
\draw[atomLine, line width=0.5pt] (-1.0000,0.0000,0.0000) -- (-0.5000,-0.8660,0.0000);
\draw[atomLine, line width=0.5pt] (-0.5000,-0.8660,0.0000) -- (0.5000,-0.8660,0.0000);
\draw[atomLine, line width=0.5pt] (0.5000,-0.8660,0.0000) -- (1.0000,0.0000,0.0000);
\draw[atomLine, line width=0.5pt] (1.0000,0.0000,1.6330) -- (0.5000,0.8660,1.6330);
\draw[atomLine, line width=0.5pt] (0.5000,0.8660,1.6330) -- (-0.5000,0.8660,1.6330);
\draw[atomLine, line width=0.5pt] (-0.5000,0.8660,1.6330) -- (-1.0000,0.0000,1.6330);
\draw[atomLine, line width=0.5pt] (-1.0000,0.0000,1.6330) -- (-0.5000,-0.8660,1.6330);
\draw[atomLine, line width=0.5pt] (-0.5000,-0.8660,1.6330) -- (0.5000,-0.8660,1.6330);
\draw[atomLine, line width=0.5pt] (0.5000,-0.8660,1.6330) -- (1.0000,0.0000,1.6330);
\draw[atomLine, line width=0.5pt] (1.0000,0.0000,0.0000) -- (1.0000,0.0000,1.6330);
\draw[atomLine, line width=0.5pt] (0.5000,0.8660,0.0000) -- (0.5000,0.8660,1.6330);
\draw[atomLine, line width=0.5pt] (-1.0000,0.0000,0.0000) -- (-1.0000,0.0000,1.6330);
\draw[atomLine, line width=0.5pt] (-0.5000,-0.8660,0.0000) -- (-0.5000,-0.8660,1.6330);
\draw[atomLine, line width=0.5pt] (0.5000,-0.8660,0.0000) -- (0.5000,-0.8660,1.6330);
\node[font=\small, anchor=north] at (0,-1.05,0) {$5$: $(11\br21)$};
\end{scope}
\begin{scope}[shift={(7.0,0,-3.0)}]
\draw[atomLine!45, dashed] (0.5000,0.8660,0.0000) -- (-0.5000,0.8660,0.0000);
\draw[atomLine!45, dashed] (-0.5000,0.8660,0.0000) -- (-1.0000,0.0000,0.0000);
\draw[atomLine!45, dashed] (-0.5000,0.8660,0.0000) -- (-0.5000,0.8660,1.6330);
\fill[fmRed!35, opacity=0.85] (0.5000,-0.8660,1.6330) -- (1.0000,0.0000,0.0000) -- (-0.5000,0.8660,0.0000) -- (-1.0000,0.0000,1.6330) -- cycle;
\draw[fmRed!80!black, line width=0.7pt] (0.5000,-0.8660,1.6330) -- (1.0000,0.0000,0.0000) -- (-0.5000,0.8660,0.0000) -- (-1.0000,0.0000,1.6330) -- cycle;
\draw[atomLine, line width=0.5pt] (1.0000,0.0000,0.0000) -- (0.5000,0.8660,0.0000);
\draw[atomLine, line width=0.5pt] (-1.0000,0.0000,0.0000) -- (-0.5000,-0.8660,0.0000);
\draw[atomLine, line width=0.5pt] (-0.5000,-0.8660,0.0000) -- (0.5000,-0.8660,0.0000);
\draw[atomLine, line width=0.5pt] (0.5000,-0.8660,0.0000) -- (1.0000,0.0000,0.0000);
\draw[atomLine, line width=0.5pt] (1.0000,0.0000,1.6330) -- (0.5000,0.8660,1.6330);
\draw[atomLine, line width=0.5pt] (0.5000,0.8660,1.6330) -- (-0.5000,0.8660,1.6330);
\draw[atomLine, line width=0.5pt] (-0.5000,0.8660,1.6330) -- (-1.0000,0.0000,1.6330);
\draw[atomLine, line width=0.5pt] (-1.0000,0.0000,1.6330) -- (-0.5000,-0.8660,1.6330);
\draw[atomLine, line width=0.5pt] (-0.5000,-0.8660,1.6330) -- (0.5000,-0.8660,1.6330);
\draw[atomLine, line width=0.5pt] (0.5000,-0.8660,1.6330) -- (1.0000,0.0000,1.6330);
\draw[atomLine, line width=0.5pt] (1.0000,0.0000,0.0000) -- (1.0000,0.0000,1.6330);
\draw[atomLine, line width=0.5pt] (0.5000,0.8660,0.0000) -- (0.5000,0.8660,1.6330);
\draw[atomLine, line width=0.5pt] (-1.0000,0.0000,0.0000) -- (-1.0000,0.0000,1.6330);
\draw[atomLine, line width=0.5pt] (-0.5000,-0.8660,0.0000) -- (-0.5000,-0.8660,1.6330);
\draw[atomLine, line width=0.5pt] (0.5000,-0.8660,0.0000) -- (0.5000,-0.8660,1.6330);
\node[font=\small, anchor=north] at (0,-1.05,0) {$6$: $(11\br22)$};
\end{scope}

\end{tikzpicture}
\caption{Плоскости скольжения в ГПУ-кристалле (гексагональная призма,
$c/a=1{,}633$): 1~— базисная $(0001)$; 2~— призматическая $(01\br10)$
из семейства $\plfam{10\br10}$; 3,\,4~— пирамидальные I рода $(0\br111)$ и
$(0\br112)$ из семейств $\plfam{10\br11}$ и $\plfam{10\br12}$; 5,\,6~—
пирамидальные II рода $(11\br21)$ и $(11\br22)$}
\label{fig:hcp-slip}
\end{figure}

Какая система работает, грубо коррелирует с отношением $c/a$
(табл.~\ref{tab:slip-systems}):
\begin{itemize}
\item при $c/a$ больше идеального $1{,}633$ базисная плоскость~— самая
плотноупакованная и далеко отстоящая от соседних; действует
\emph{базисное} скольжение (Zn, $c/a=1{,}86$; Cd, $1{,}89$); у Mg
($c/a=1{,}62\approx1{,}633$)~— тоже базисное;
\item при $c/a<1{,}633$ ($\alpha$-Ti, $\alpha$-Zr, $c/a\approx1{,}59$)
преобладает \emph{призматическое} скольжение;
\item но бериллий ($c/a=1{,}57$) скользит по базису~— правило $c/a$
лишь эмпирическое: решают детали межатомной связи.
\end{itemize}

\begin{table}[ht]
\centering
\caption{Основные системы скольжения металлов}
\label{tab:slip-systems}
\small
\begin{tabular}{@{}llccl@{}}
\toprule
Решётка & Система & Число систем & Независимых & Примеры\\
\midrule
ГЦК & $\plfam{111}\nfam{1\br10}$ & 12 & 5 & Al, Cu, Ni, Au, $\gamma$-Fe\\
ОЦК & $\plfam{110}\nfam{\br111}$ & 12 & 5 & $\alpha$-Fe, Mo, W, Nb\\
    & $\plfam{112}\nfam{\br1\br11}$ & 12 & 5 & $\alpha$-Fe, W\\
    & $\plfam{123}\nfam{\br1\br11}$ & 24 & 5 & $\alpha$-Fe\\
ГПУ & $(0001)\nfam{11\br20}$ & 3 & 2 & Zn, Cd, Mg, Be\\
    & $\plfam{10\br10}\nfam{11\br20}$ & 3 & 2 & Ti, Zr\\
    & $\plfam{10\br11}\nfam{11\br20}$ & 6 & 4 & Ti, Mg\\
    & $\plfam{11\br22}\nfam{11\br2\br3}$ & 6 & 5 & Zn, Cd, Ti ($\vb c+\vb a$)\\
\bottomrule
\end{tabular}
\end{table}

Сдвиг вдоль $\nfam{11\br20}$ не меняет размеров кристалла вдоль оси $c$,
сколько бы таких систем ни работало. Поэтому для произвольного
формоизменения ГПУ-металлу нужны пирамидальное скольжение с вектором
$\vb c+\vb a$ (направление $\nfam{11\br2\br3}$) или двойникование~— иначе
поликристалл хрупок.

\subsubsection{Механическое двойникование}

Двойникование играет в деформации сравнительно небольшую роль и
проявляется, когда скольжение затруднено. Ему способствуют низкая
температура и ударное нагружение: с понижением температуры и ростом
скорости деформации критическое напряжение скольжения растёт быстрее, чем
двойникования. Легче всего двойникуются гексагональные металлы, труднее~—
ОЦК, труднее всего~— ГЦК.

В отличие от скольжения, при двойниковании атомы смещаются на \emph{доли}
вектора трансляции, а смещение каждого атома пропорционально его
расстоянию от плоскости двойникования (рис.~\ref{fig:slip-twin},\,\textit{в}).
Такой \term{простой сдвиг}~— однородная деформация: мысленно выделенный в
кристалле шар превращается в эллипсоид. Опишем её
(рис.~\ref{fig:twin-shear}). Ось $y$ направим вдоль направления сдвига
$\eta_1$, ось $z$~— по нормали к плоскости сдвига $K_1$; верхняя
половина шара единичного радиуса сдвигается: $(y,z)\to(y+sz,\,z)$.

\begin{figure}[ht]
\centering
\begin{tikzpicture}[scale=1.9]
  \def\s{0.707}
  \fill[fmBlueBg] (-1,0) arc[start angle=180, end angle=360, radius=1] -- cycle;
  \fill[fmAmber!15] plot[domain=0:180, samples=73, smooth] ({cos(\x)+\s*sin(\x)},{sin(\x)}) -- cycle;
  \draw[atomLine, dashed] (1,0) arc[start angle=0, end angle=180, radius=1];
  \draw[fmAmber!80!black, line width=0.9pt] plot[domain=0:180, samples=73, smooth] ({cos(\x)+\s*sin(\x)},{sin(\x)});
  \draw[atomLine, line width=0.6pt] (-1,0) arc[start angle=180, end angle=360, radius=1];
  \draw[fmGray, line width=1.1pt] (-1.6,0) -- (1.9,0) node[right, font=\small]{$K_1$};
  \draw[bvec, fmBlue] (1.05,-0.18) -- (1.65,-0.18) node[below, font=\small]{$\eta_1$};
  % неискажённое направление: до (A) и после (A') сдвига
  \pgfmathsetmacro{\nn}{sqrt(1+\s*\s/4)}
  \coordinate (A) at ({-\s/2/\nn},{1/\nn});
  \coordinate (A2) at ({\s/2/\nn},{1/\nn});
  \draw[fmRed, line width=0.8pt] ($-0.5*(A)$) -- ($1.45*(A)$) node[above, font=\small]{$K_2$};
  \draw[fmRed!70!black, line width=0.8pt, dashed] (0,0) -- ($1.45*(A2)$) node[above, font=\small]{$K_2'$};
  \fill (A) circle (0.03) node[left, font=\small]{$A$};
  \fill (A2) circle (0.03) node[right, font=\small]{$A'$};
  \draw[arr, fmRed] (A) -- (A2);
  \node[font=\footnotesize, fmRed] at ($(A)!0.5!(A2)+(0,0.16)$) {$\eta_2$};
  \fill (0,0) circle (0.025) node[below, font=\small]{$O$};
  \draw[arr, fmAmber!80!black] (0,1) -- (\s,1);
  \node[font=\footnotesize, above] at ({\s/2+0.05},1.02) {$s$};
  \draw[fmGray, densely dotted] (0,0) -- (0,1) (\s,1) -- (0,0);
\end{tikzpicture}
\caption{Однородный сдвиг при двойниковании: верхняя половина единичной
сферы (пунктир) переходит в половину эллипсоида. Неискажёнными остаются
плоскость сдвига $K_1$ и плоскость $K_2$: направление $OA$ ($\eta_2$)
переходит в $OA'$ той же длины; $s$~— кристаллографический сдвиг
(на рисунке $s=0{,}707$, как у ГЦК и ОЦК)}
\label{fig:twin-shear}
\end{figure}

У эллипсоида два круговых \emph{неискажённых} сечения: плоскость сдвига
$K_1$ (она не меняется вовсе) и плоскость $K_2$, перпендикулярная
плоскости рисунка: направление $OA$ ($\eta_2$) переходит в $OA'$ без
изменения длины. Действительно, вектор $(y,z)$ сохраняет длину, если
$(y+sz)^2+z^2=y^2+z^2$, т.\,е. $y=-sz/2$. Отсюда угол $2\phi$ между $K_1$
и $K_2$ связан со сдвигом:
\begin{equation}\label{eq:twin-shear}
s=2\operatorname{ctg}2\phi .
\end{equation}

\begin{opredelenie}
\term{Элементы двойникования}: плоскости $K_1$ (плоскость двойникования)
и $K_2$ (вторая неискажённая плоскость), направления $\eta_1$ (направление
двойникования, лежит в $K_1$) и $\eta_2$ (лежит в $K_2$, перпендикулярно
линии пересечения $K_1$ и $K_2$) и величина $s$~— \term{кристаллографический
сдвиг}.
\end{opredelenie}

\begin{table}[ht]
\centering
\caption{Элементы двойникования металлов}
\label{tab:twinning}
\small
\begin{tabular}{@{}lccccl@{}}
\toprule
Решётка & $K_1$ & $\eta_1$ & $K_2$ & $\eta_2$ & $s$\\
\midrule
ГЦК & $(111)$ & $[11\br2]$ & $(11\br1)$ & $[112]$ & $1/\sqrt2=0{,}707$\\
ОЦК & $(112)$ & $[11\br1]$ & $(11\br2)$ & $[111]$ & $1/\sqrt2=0{,}707$\\
ГПУ & $(10\br12)$ & $[\br1011]$ & $(\br1012)$ & $[10\br11]$ & $\dfrac{|3-(c/a)^2|}{\sqrt3\,c/a}$\\
\bottomrule
\end{tabular}
\end{table}

В ГЦК- и ОЦК-кристаллах обе плоскости $K_1$, $K_2$ принадлежат одному
семейству ($\plfam{111}$ и $\plfam{112}$), направления $\eta_1$, $\eta_2$~—
семействам $\nfam{112}$ и $\nfam{111}$ соответственно, а $s=0{,}707$
(табл.~\ref{tab:twinning}; угол между плоскостями $\plfam{111}$
$2\phi=70{,}5^\circ$ и по~\eqref{eq:twin-shear} $s=2\operatorname{ctg}70{,}5^\circ=0{,}707$).
В ГПУ-металлах ($\alpha$-Zr, Ti, Mg, Zn) главная система~—
$\plfam{10\br12}\nfam{\br1011}$, а сдвиг зависит от $c/a$: у Zr
($c/a=1{,}593$) $s=0{,}167$, у Mg $0{,}129$, у Zn $0{,}138$. При
$c/a<\sqrt3$ такой двойник удлиняет кристалл вдоль оси $c$, при $c/a>\sqrt3$~—
укорачивает.

\begin{vnimanie}
Скольжение не меняет ориентацию решётки внутри сдвинутой части (лишь
поворачивает кристалл как целое), а двойникование \emph{меняет}: двойник~—
это область с другой, симметрично связанной ориентацией. Поэтому
двойники видны на шлифе после травления как полосы иного оттенка, а
следы скольжения~— только как ступеньки на поверхности.
\end{vnimanie}

\subsection{Кристаллографическая текстура}\label{ssec:texture}

При пластической деформации поликристалла каждое зерно (если в нём
работают пять независимых систем деформации) меняет форму вместе с
соседями без образования пустот. При этом решётка зерна \emph{поворачивается}
(антисимметричная часть $\omega_{ij}$ тензора дисторсии), и ориентации
разных зёрен постепенно сближаются.

\begin{opredelenie}
\term{Текстура}~— преимущественная ориентация кристаллитов в
поликристалле.
\end{opredelenie}

Текстуры возникают при деформации (волочении, прокатке), рекристаллизации,
литье, осаждении плёнок и сильно влияют на анизотропию свойств изделия:
упругих модулей, магнитной проницаемости трансформаторной стали,
радиационного роста циркониевых оболочек твэлов.

\subsubsection{Аксиальная и ограниченная текстуры}
В простейшем случае определённые кристаллографические направления
$\nfam{mnp}$ всех кристаллитов параллельны одному внешнему направлению
$F$~— \term{оси текстуры}, а вокруг этой оси кристаллиты повёрнуты
произвольно. Такую текстуру называют \term{аксиальной} (осевой,
волокнистой, неограниченной). Она типична для проволоки: в ОЦК-металлах
ось текстуры волочения~— $\nfam{110}$, в ГЦК~— $\nfam{111}$ и
$\nfam{100}$ (доля $\nfam{100}$ растёт с уменьшением энергии дефекта
упаковки).

В прокатанных листах, лентах и фольге возникает \term{ограниченная}
текстура (текстура прокатки): определённая плоскость $(hkl)$
кристаллитов параллельна плоскости прокатки (поверхности листа), а
направление $[uvw]$~— направлению прокатки. Её записывают как
$\plfam{hkl}\nfam{uvw}$. Например, в ГЦК-металлах распространены
компоненты «медная» $\plfam{112}\nfam{111}$, «латунная»
$\plfam{110}\nfam{112}$, Госса $\plfam{110}\nfam{001}$ и после
рекристаллизации~— кубическая $\plfam{001}\nfam{100}$.

Текстуру описывают тремя способами: прямыми полюсными фигурами (ППФ),
обратными полюсными фигурами (ОПФ) и функцией распределения ориентаций
(ФРО). Все три сравнивают между собой две системы координат; начнём с
них, а затем разберём каждый способ.

\subsubsection{Две системы координат}

\begin{itemize}
\item \term{Система координат образца} привязана к изделию. Для листа это
направление прокатки (НП), поперечное направление (ПН) и нормаль к
листу (НН); для проволоки~— её ось. Она одна на весь лист.
\item \term{Система координат кристалла}~— оси $[100]$, $[010]$, $[001]$
решётки. В каждом зерне она своя.
\end{itemize}
\term{Ориентация} зерна~— поворот, совмещающий одну систему с другой
(матрица поворота или углы Эйлера, п.~\ref{ssec:orth}). Удобно представить
каждое зерно кубиком, рёбра которого идут вдоль $\nfam{100}$
(рис.~\ref{fig:texture-frames}). В металле без текстуры кубики повёрнуты
как попало. При \emph{кубической} текстуре $\plfam{001}\nfam{100}$ почти все
они «стоят» одинаково: грани параллельны поверхности и краям листа,
синие грани $\perp[001]$ смотрят вверх (рис.~\ref{fig:texture-frames},\,\textit{б}).

\begin{figure}[p]
\centering
\begin{tikzpicture}
  % создано tools/gen_plastic.py: системы координат образца и кристалла (текстура)
\fill[atomA!25] (4.600,0.900) -- (10.300,0.900) -- (10.300,1.565) -- (4.600,1.565) -- cycle;
\fill[atomA!40] (10.300,0.900) -- (11.896,2.040) -- (11.896,2.705) -- (10.300,1.565) -- cycle;
\fill[atomA!12] (4.600,1.565) -- (10.300,1.565) -- (11.896,2.705) -- (6.196,2.705) -- cycle;
\draw[atomLine!70, line width=0.5pt] (4.600,1.565) -- (5.650,1.565) -- (6.211,1.869) -- (5.235,2.019) -- cycle;
\draw[atomLine!70, line width=0.5pt] (5.650,1.565) -- (7.052,1.565) -- (7.097,1.885) -- (6.686,1.933) -- (6.211,1.869) -- cycle;
\draw[atomLine!70, line width=0.5pt] (7.052,1.565) -- (7.672,1.565) -- (8.495,1.918) -- (7.969,1.981) -- (7.097,1.885) -- cycle;
\draw[atomLine!70, line width=0.5pt] (7.672,1.565) -- (9.209,1.565) -- (9.497,1.854) -- (8.917,1.964) -- (8.495,1.918) -- cycle;
\draw[atomLine!70, line width=0.5pt] (9.209,1.565) -- (10.300,1.565) -- (11.043,2.095) -- (9.497,1.854) -- cycle;
\draw[atomLine!70, line width=0.5pt] (5.580,2.265) -- (5.235,2.019) -- (6.211,1.869) -- (6.686,1.933) -- (7.225,2.262) -- (7.020,2.313) -- cycle;
\draw[atomLine!70, line width=0.5pt] (7.225,2.262) -- (6.686,1.933) -- (7.097,1.885) -- (7.969,1.981) -- (8.179,2.302) -- (7.973,2.331) -- cycle;
\draw[atomLine!70, line width=0.5pt] (8.179,2.302) -- (7.969,1.981) -- (8.495,1.918) -- (8.917,1.964) -- (9.448,2.286) -- (9.086,2.368) -- cycle;
\draw[atomLine!70, line width=0.5pt] (11.043,2.095) -- (11.081,2.123) -- (10.222,2.374) -- (9.448,2.286) -- (8.917,1.964) -- (9.497,1.854) -- cycle;
\draw[atomLine!70, line width=0.5pt] (7.417,2.705) -- (6.196,2.705) -- (5.580,2.265) -- (7.020,2.313) -- cycle;
\draw[atomLine!70, line width=0.5pt] (8.447,2.705) -- (7.417,2.705) -- (7.020,2.313) -- (7.225,2.262) -- (7.973,2.331) -- cycle;
\draw[atomLine!70, line width=0.5pt] (9.474,2.705) -- (8.447,2.705) -- (7.973,2.331) -- (8.179,2.302) -- (9.086,2.368) -- cycle;
\draw[atomLine!70, line width=0.5pt] (10.883,2.705) -- (9.474,2.705) -- (9.086,2.368) -- (9.448,2.286) -- (10.222,2.374) -- cycle;
\draw[atomLine!70, line width=0.5pt] (11.081,2.123) -- (11.896,2.705) -- (10.883,2.705) -- (10.222,2.374) -- cycle;
\draw[fmAmber!80!black, line width=1.3pt] (8.179,2.302) -- (7.969,1.981) -- (8.495,1.918) -- (8.917,1.964) -- (9.448,2.286) -- (9.086,2.368) -- cycle;
\draw[atomLine, line width=0.6pt] (4.600,0.900) -- (10.300,0.900) -- (11.896,2.040) -- (11.896,2.705) -- (6.196,2.705) -- (4.600,1.565) -- cycle (4.600,1.565) -- (10.300,1.565) -- (11.896,2.705) (10.300,1.565) -- (10.300,0.900);
\filldraw[fmRed!35, draw=fmRed!60!black, line width=0.3pt] (9.913,2.552) -- (10.187,2.729) -- (10.162,3.192) -- (9.887,3.015) -- cycle;
\filldraw[fmGreen!35, draw=fmGreen!60!black, line width=0.3pt] (9.478,2.572) -- (9.452,3.035) -- (9.887,3.015) -- (9.913,2.552) -- cycle;
\filldraw[fmBlue!35, draw=fmBlue!60!black, line width=0.3pt] (9.452,3.035) -- (9.887,3.015) -- (10.162,3.192) -- (9.727,3.212) -- cycle;
\filldraw[fmRed!35, draw=fmRed!60!black, line width=0.3pt] (9.069,3.149) -- (8.607,3.252) -- (8.379,2.927) -- (8.841,2.824) -- cycle;
\filldraw[fmGreen!35, draw=fmGreen!60!black, line width=0.3pt] (9.063,2.788) -- (8.835,2.464) -- (8.841,2.824) -- (9.069,3.149) -- cycle;
\filldraw[fmBlue!35, draw=fmBlue!60!black, line width=0.3pt] (8.835,2.464) -- (8.841,2.824) -- (8.379,2.927) -- (8.373,2.567) -- cycle;
\filldraw[fmRed!35, draw=fmRed!60!black, line width=0.3pt] (7.809,2.513) -- (8.015,2.693) -- (8.020,3.155) -- (7.814,2.975) -- cycle;
\filldraw[fmGreen!35, draw=fmGreen!60!black, line width=0.3pt] (7.337,2.536) -- (7.341,2.998) -- (7.814,2.975) -- (7.809,2.513) -- cycle;
\filldraw[fmBlue!35, draw=fmBlue!60!black, line width=0.3pt] (7.341,2.998) -- (7.814,2.975) -- (8.020,3.155) -- (7.547,3.178) -- cycle;
\filldraw[fmRed!35, draw=fmRed!60!black, line width=0.3pt] (6.905,2.501) -- (7.052,2.701) -- (7.019,3.154) -- (6.872,2.955) -- cycle;
\filldraw[fmGreen!35, draw=fmGreen!60!black, line width=0.3pt] (6.412,2.473) -- (6.379,2.926) -- (6.872,2.955) -- (6.905,2.501) -- cycle;
\filldraw[fmBlue!35, draw=fmBlue!60!black, line width=0.3pt] (6.379,2.926) -- (6.872,2.955) -- (7.019,3.154) -- (6.526,3.126) -- cycle;
\filldraw[fmRed!35, draw=fmRed!60!black, line width=0.3pt] (11.287,2.510) -- (11.432,2.603) -- (11.473,3.090) -- (11.327,2.997) -- cycle;
\filldraw[fmGreen!35, draw=fmGreen!60!black, line width=0.3pt] (10.794,2.520) -- (10.835,3.007) -- (11.327,2.997) -- (11.287,2.510) -- cycle;
\filldraw[fmBlue!35, draw=fmBlue!60!black, line width=0.3pt] (10.835,3.007) -- (11.327,2.997) -- (11.473,3.090) -- (10.980,3.100) -- cycle;
\filldraw[fmRed!35, draw=fmRed!60!black, line width=0.3pt] (8.293,2.244) -- (8.732,2.137) -- (8.947,2.598) -- (8.509,2.705) -- cycle;
\filldraw[fmGreen!35, draw=fmGreen!60!black, line width=0.3pt] (8.131,2.393) -- (8.346,2.854) -- (8.509,2.705) -- (8.293,2.244) -- cycle;
\filldraw[fmBlue!35, draw=fmBlue!60!black, line width=0.3pt] (8.346,2.854) -- (8.509,2.705) -- (8.947,2.598) -- (8.785,2.747) -- cycle;
\filldraw[fmRed!35, draw=fmRed!60!black, line width=0.3pt] (10.015,2.167) -- (10.197,2.286) -- (10.205,2.767) -- (10.023,2.648) -- cycle;
\filldraw[fmGreen!35, draw=fmGreen!60!black, line width=0.3pt] (9.533,2.161) -- (9.541,2.642) -- (10.023,2.648) -- (10.015,2.167) -- cycle;
\filldraw[fmBlue!35, draw=fmBlue!60!black, line width=0.3pt] (9.541,2.642) -- (10.023,2.648) -- (10.205,2.767) -- (9.723,2.762) -- cycle;
\filldraw[fmRed!35, draw=fmRed!60!black, line width=0.3pt] (6.505,2.143) -- (6.695,2.307) -- (6.699,2.775) -- (6.509,2.611) -- cycle;
\filldraw[fmGreen!35, draw=fmGreen!60!black, line width=0.3pt] (6.026,2.153) -- (6.030,2.621) -- (6.509,2.611) -- (6.505,2.143) -- cycle;
\filldraw[fmBlue!35, draw=fmBlue!60!black, line width=0.3pt] (6.030,2.621) -- (6.509,2.611) -- (6.699,2.775) -- (6.220,2.785) -- cycle;
\filldraw[fmRed!35, draw=fmRed!60!black, line width=0.3pt] (7.757,2.141) -- (7.939,2.308) -- (7.894,2.773) -- (7.712,2.606) -- cycle;
\filldraw[fmGreen!35, draw=fmGreen!60!black, line width=0.3pt] (7.277,2.101) -- (7.232,2.567) -- (7.712,2.606) -- (7.757,2.141) -- cycle;
\filldraw[fmBlue!35, draw=fmBlue!60!black, line width=0.3pt] (7.232,2.567) -- (7.712,2.606) -- (7.894,2.773) -- (7.414,2.733) -- cycle;
\filldraw[fmRed!35, draw=fmRed!60!black, line width=0.3pt] (8.017,2.107) -- (7.643,2.290) -- (7.313,1.900) -- (7.687,1.718) -- cycle;
\filldraw[fmGreen!35, draw=fmGreen!60!black, line width=0.3pt] (7.513,2.536) -- (7.183,2.146) -- (7.313,1.900) -- (7.643,2.290) -- cycle;
\filldraw[fmBlue!35, draw=fmBlue!60!black, line width=0.3pt] (7.887,2.353) -- (8.017,2.107) -- (7.643,2.290) -- (7.513,2.536) -- cycle;
\filldraw[fmRed!35, draw=fmRed!60!black, line width=0.3pt] (10.311,1.818) -- (10.480,1.872) -- (10.520,2.365) -- (10.351,2.311) -- cycle;
\filldraw[fmGreen!35, draw=fmGreen!60!black, line width=0.3pt] (9.826,1.819) -- (9.866,2.312) -- (10.351,2.311) -- (10.311,1.818) -- cycle;
\filldraw[fmBlue!35, draw=fmBlue!60!black, line width=0.3pt] (9.866,2.312) -- (10.351,2.311) -- (10.520,2.365) -- (10.035,2.366) -- cycle;
\filldraw[fmRed!35, draw=fmRed!60!black, line width=0.3pt] (5.404,1.748) -- (5.652,1.953) -- (5.606,2.405) -- (5.358,2.200) -- cycle;
\filldraw[fmGreen!35, draw=fmGreen!60!black, line width=0.3pt] (4.955,1.751) -- (4.909,2.203) -- (5.358,2.200) -- (5.404,1.748) -- cycle;
\filldraw[fmBlue!35, draw=fmBlue!60!black, line width=0.3pt] (4.909,2.203) -- (5.358,2.200) -- (5.606,2.405) -- (5.157,2.408) -- cycle;
\filldraw[fmRed!35, draw=fmRed!60!black, line width=0.3pt] (8.710,1.734) -- (8.932,1.920) -- (8.863,2.378) -- (8.641,2.192) -- cycle;
\filldraw[fmGreen!35, draw=fmGreen!60!black, line width=0.3pt] (8.250,1.694) -- (8.182,2.152) -- (8.641,2.192) -- (8.710,1.734) -- cycle;
\filldraw[fmBlue!35, draw=fmBlue!60!black, line width=0.3pt] (8.182,2.152) -- (8.641,2.192) -- (8.863,2.378) -- (8.403,2.337) -- cycle;
\filldraw[fmRed!35, draw=fmRed!60!black, line width=0.3pt] (6.716,1.716) -- (6.893,1.887) -- (6.869,2.352) -- (6.692,2.181) -- cycle;
\filldraw[fmGreen!35, draw=fmGreen!60!black, line width=0.3pt] (6.233,1.697) -- (6.209,2.162) -- (6.692,2.181) -- (6.716,1.716) -- cycle;
\filldraw[fmBlue!35, draw=fmBlue!60!black, line width=0.3pt] (6.209,2.162) -- (6.692,2.181) -- (6.869,2.352) -- (6.386,2.333) -- cycle;
\fill[atomA!40] (0.150,1.250) -- (3.550,1.250) -- (3.550,1.410) -- (0.150,1.410) -- cycle;
\fill[atomA!55] (3.550,1.250) -- (4.474,1.910) -- (4.474,2.070) -- (3.550,1.410) -- cycle;
\fill[atomA!18] (0.150,1.410) -- (3.550,1.410) -- (4.474,2.070) -- (1.074,2.070) -- cycle;
\draw[atomLine, line width=0.6pt] (0.150,1.250) -- (3.550,1.250) -- (4.474,1.910) -- (4.474,2.070) -- (1.074,2.070) -- (0.150,1.410) -- cycle (0.150,1.410) -- (3.550,1.410) -- (4.474,2.070) (3.550,1.410) -- (3.550,1.250);
\draw[fmAmber!80!black, line width=0.8pt] (2.744,1.620) -- (3.194,1.620) -- (3.320,1.710) -- (2.870,1.710) -- cycle;
\draw[atomLine!50, densely dotted] (2.870,1.710) -- (6.196,2.705) (2.744,1.620) -- (4.600,1.565);
\draw[-{Stealth[length=1.8mm]}, atomLine, line width=0.8pt] (-0.452,1.070) -- (0.548,1.070) node[font=\scriptsize, anchor=north] {НП};
\draw[-{Stealth[length=1.8mm]}, atomLine, line width=0.8pt] (-0.452,1.070) -- (0.010,1.400) node[font=\scriptsize, anchor=east] {ПН};
\draw[-{Stealth[length=1.8mm]}, atomLine, line width=0.8pt] (-0.452,1.070) -- (-0.452,2.070) node[font=\scriptsize, anchor=south] {НН};
\coordinate (Alab) at (1.550,0.400);
\coordinate (Bsel) at (8.539,2.154);
\coordinate (Bfront) at (7.450,0.900);
\coordinate (Bright) at (11.098,1.802);
\coordinate (Btl) at (6.196,2.705);
\coordinate (Bbl) at (4.600,0.900);
\filldraw[fmRed!35, draw=fmRed!60!black, line width=0.3pt, opacity=0.75] (0.952,-3.561) -- (2.107,-3.842) -- (2.675,-2.629) -- (1.520,-2.348) -- cycle;
\filldraw[fmGreen!35, draw=fmGreen!60!black, line width=0.3pt, opacity=0.75] (0.525,-3.171) -- (1.093,-1.958) -- (1.520,-2.348) -- (0.952,-3.561) -- cycle;
\filldraw[fmBlue!35, draw=fmBlue!60!black, line width=0.3pt, opacity=0.75] (1.093,-1.958) -- (1.520,-2.348) -- (2.675,-2.629) -- (2.248,-2.239) -- cycle;
\draw[-{Stealth[length=2mm]}, fmRed!80!black, line width=1.1pt] (1.6,-2.9) -- (2.129,-3.384) node[font=\scriptsize, fmRed!80!black, anchor=west] {$[100]$};
\draw[-{Stealth[length=2mm]}, fmGreen!80!black, line width=1.1pt] (1.6,-2.9) -- (0.168,-2.551) node[font=\scriptsize, fmGreen!80!black, anchor=east] {$[0\br10]$};
\draw[-{Stealth[length=2mm]}, fmBlue!80!black, line width=1.1pt] (1.6,-2.9) -- (2.304,-1.396) node[font=\scriptsize, fmBlue!80!black, anchor=west] {$[001]$};
\draw[-{Stealth[length=2mm]}, atomLine, dashed, line width=0.8pt] (1.6,-2.9) -- (3.500,-2.900) node[font=\scriptsize, anchor=west] {НП};
\draw[-{Stealth[length=2mm]}, atomLine, dashed, line width=0.8pt] (1.6,-2.9) -- (2.398,-2.330) node[font=\scriptsize, anchor=south west] {ПН};
\draw[-{Stealth[length=2mm]}, atomLine, dashed, line width=0.8pt] (1.6,-2.9) -- (1.600,-1.000) node[font=\scriptsize, anchor=south] {НН};
\coordinate (Csel) at (1.6,-1.5499999999999998);
\draw[atomLine, line width=0.6pt] (6.3,-2.75) circle (1.45);
\draw[atomLine!40] (4.85,-2.75) -- (7.75,-2.75) (6.3,-4.2) -- (6.3,-1.3);
\filldraw[fmRed] (6.141,-1.348) circle (0.055);
\filldraw[fmGreen] (7.558,-2.612) circle (0.055);
\filldraw[fmBlue] (6.204,-2.781) circle (0.055);
\filldraw[fmRed] (6.364,-1.340) circle (0.055);
\filldraw[fmGreen] (7.664,-2.814) circle (0.055);
\filldraw[fmBlue] (6.255,-2.767) circle (0.055);
\filldraw[fmRed] (7.146,-3.411) circle (0.055);
\filldraw[fmGreen] (5.719,-3.231) circle (0.055);
\filldraw[fmBlue] (6.452,-2.244) circle (0.055);
\filldraw[fmRed] (6.222,-1.445) circle (0.055);
\filldraw[fmGreen] (7.619,-2.684) circle (0.055);
\filldraw[fmBlue] (6.236,-2.829) circle (0.055);
\filldraw[fmRed] (6.223,-4.172) circle (0.055);
\filldraw[fmGreen] (5.095,-2.681) circle (0.055);
\filldraw[fmBlue] (6.433,-2.744) circle (0.055);
\filldraw[fmRed] (6.270,-4.162) circle (0.055);
\filldraw[fmGreen] (7.684,-2.778) circle (0.055);
\filldraw[fmBlue] (6.267,-2.730) circle (0.055);
\filldraw[fmRed] (6.343,-1.405) circle (0.055);
\filldraw[fmGreen] (7.676,-2.799) circle (0.055);
\filldraw[fmBlue] (6.261,-2.803) circle (0.055);
\filldraw[fmRed, draw=fmAmber!80!black, line width=0.9pt] (7.299,-3.625) circle (0.075);
\filldraw[fmGreen, draw=fmAmber!80!black, line width=0.9pt] (5.657,-3.428) circle (0.075);
\filldraw[fmBlue, draw=fmAmber!80!black, line width=0.9pt] (6.462,-2.473) circle (0.075);
\filldraw[fmRed] (6.354,-1.303) circle (0.055);
\filldraw[fmGreen] (4.919,-2.698) circle (0.055);
\filldraw[fmBlue] (6.335,-2.752) circle (0.055);
\filldraw[fmRed] (6.444,-1.348) circle (0.055);
\filldraw[fmGreen] (7.572,-2.885) circle (0.055);
\filldraw[fmBlue] (6.208,-2.761) circle (0.055);
\filldraw[fmRed] (6.316,-4.136) circle (0.055);
\filldraw[fmGreen] (7.638,-2.730) circle (0.055);
\filldraw[fmBlue] (6.241,-2.718) circle (0.055);
\filldraw[fmRed] (5.974,-2.603) circle (0.055);
\filldraw[fmGreen] (6.311,-3.915) circle (0.055);
\filldraw[fmBlue] (7.230,-2.648) circle (0.055);
\filldraw[fmRed] (6.062,-1.333) circle (0.055);
\filldraw[fmGreen] (7.623,-2.528) circle (0.055);
\filldraw[fmBlue] (6.245,-2.766) circle (0.055);
\filldraw[fmRed] (6.156,-4.119) circle (0.055);
\filldraw[fmGreen] (4.998,-2.606) circle (0.055);
\filldraw[fmBlue] (6.377,-2.720) circle (0.055);
\coordinate (PFc) at (6.3,-2.75); \coordinate (PFtop) at (6.3,-1.3); \coordinate (PFright) at (7.75,-2.75);
\filldraw[fmAmber!12, draw=atomLine, line width=0.6pt] (9.050,-3.600) -- (9.139,-3.600) -- (9.229,-3.600) -- (9.318,-3.600) -- (9.408,-3.600) -- (9.498,-3.600) -- (9.587,-3.600) -- (9.677,-3.600) -- (9.768,-3.600) -- (9.858,-3.600) -- (9.949,-3.600) -- (10.040,-3.600) -- (10.132,-3.600) -- (10.224,-3.600) -- (10.316,-3.600) -- (10.409,-3.600) -- (10.503,-3.600) -- (10.597,-3.600) -- (10.691,-3.600) -- (10.787,-3.600) -- (10.882,-3.600) -- (10.979,-3.600) -- (11.077,-3.600) -- (11.175,-3.600) -- (11.274,-3.600) -- (11.374,-3.600) -- (11.475,-3.600) -- (11.577,-3.600) -- (11.680,-3.600) -- (11.784,-3.600) -- (11.784,-3.600) -- (11.783,-3.518) -- (11.782,-3.436) -- (11.781,-3.354) -- (11.778,-3.272) -- (11.775,-3.190) -- (11.771,-3.107) -- (11.766,-3.025) -- (11.761,-2.943) -- (11.754,-2.860) -- (11.748,-2.778) -- (11.740,-2.695) -- (11.731,-2.613) -- (11.722,-2.530) -- (11.712,-2.447) -- (11.702,-2.364) -- (11.690,-2.281) -- (11.678,-2.198) -- (11.665,-2.114) -- (11.651,-2.031) -- (11.636,-1.947) -- (11.621,-1.863) -- (11.604,-1.779) -- (11.587,-1.695) -- (11.569,-1.610) -- (11.550,-1.525) -- (11.531,-1.440) -- (11.510,-1.355) -- (11.488,-1.270) -- (11.466,-1.184) -- (11.466,-1.184) -- (11.369,-1.281) -- (11.274,-1.376) -- (11.180,-1.470) -- (11.088,-1.562) -- (10.997,-1.653) -- (10.908,-1.742) -- (10.819,-1.831) -- (10.732,-1.918) -- (10.646,-2.004) -- (10.560,-2.090) -- (10.476,-2.174) -- (10.392,-2.258) -- (10.309,-2.341) -- (10.227,-2.423) -- (10.146,-2.504) -- (10.065,-2.585) -- (9.985,-2.665) -- (9.905,-2.745) -- (9.826,-2.824) -- (9.747,-2.903) -- (9.669,-2.981) -- (9.590,-3.060) -- (9.513,-3.137) -- (9.435,-3.215) -- (9.358,-3.292) -- (9.281,-3.369) -- (9.204,-3.446) -- (9.127,-3.523) -- (9.050,-3.600) -- cycle;
\filldraw[atomLine] (9.498,-3.510) circle (0.055);
\filldraw[atomLine] (9.249,-3.511) circle (0.055);
\filldraw[atomLine] (11.195,-2.509) circle (0.055);
\filldraw[atomLine] (9.392,-3.291) circle (0.055);
\filldraw[atomLine] (9.653,-3.539) circle (0.055);
\filldraw[atomLine] (9.202,-3.512) circle (0.055);
\filldraw[atomLine] (9.296,-3.430) circle (0.055);
\filldraw[atomLine, draw=fmAmber!80!black, line width=0.9pt] (10.480,-3.297) circle (0.075);
\filldraw[atomLine] (9.209,-3.596) circle (0.055);
\filldraw[atomLine] (9.462,-3.505) circle (0.055);
\filldraw[atomLine] (9.315,-3.450) circle (0.055);
\filldraw[atomLine] (10.492,-2.846) circle (0.055);
\filldraw[atomLine] (9.307,-3.570) circle (0.055);
\filldraw[atomLine] (9.385,-3.427) circle (0.055);
\coordinate (I100) at (9.050,-3.600);
\coordinate (I110) at (11.784,-3.600);
\coordinate (I111) at (11.466,-1.184);

  \draw[-{Stealth[length=2mm]}, fmAmber!80!black, line width=0.9pt] (Bsel) to[bend right=25] (Csel);
  \node[font=\small] at (Alab) {\textit{а}\ лист (см)};
  \node[font=\small] at (8.0,0.25) {\textit{б}\ зёрна (10--100 мкм)};
  \node[font=\small] at (1.6,-4.75) {\textit{в}\ одно зерно};
  \node[font=\small] at (6.3,-4.75) {\textit{г}\ ППФ $(100)$};
  \node[font=\small] at (10.4,-4.75) {\textit{д}\ ОПФ для НН};
  \node[font=\scriptsize, anchor=south] at (PFtop) {НП};
  \node[font=\scriptsize, anchor=west] at (PFright) {ПН};
  \node[font=\scriptsize, fill=white, inner sep=0.5pt] at ($(PFc)+(0.42,-0.3)$) {НН};
  \node[font=\scriptsize, anchor=north east] at (I100) {$[100]$};
  \node[font=\scriptsize, anchor=north west] at (I110) {$[110]$};
  \node[font=\scriptsize, anchor=south] at (I111) {$[111]$};
\end{tikzpicture}
\caption{Две системы координат при описании текстуры. \textit{а}~—
лист с осями образца НП, ПН, НН (сантиметры); рамкой отмечен
увеличиваемый участок. \textit{б}~— участок поверхности: каждое зерно
изображено кубиком с рёбрами вдоль $\nfam{100}$ (грани $\perp[100]$~—
красные, $\perp[010]$~— зелёные, $\perp[001]$~— синие); большинство зёрен
имеют кубическую текстуру $\plfam{001}\nfam{100}$ с разбросом
$\sim10^\circ$, три зерна ориентированы случайно. \textit{в}~— выделенное
зерно (оранжевая граница) крупно: оси кристалла (цветные стрелки) и оси
образца (пунктир) из одной точки. \textit{г}~— прямая полюсная фигура
$(100)$ в осях образца (НН в центре): каждая точка~— направление одного
ребра одного кубика, её цвет~— какое это ребро; у кубической текстуры
точки собираются у НН, НП и ПН. \textit{д}~— обратная полюсная фигура
для НН в стандартном треугольнике осей кристалла: каждая точка~— одно
зерно; у кубической текстуры НН $\parallel\nfam{100}$, и точки собираются у
вершины $[100]$. Выделенное зерно обведено оранжевым}
\label{fig:texture-frames}
\end{figure}

Все способы описания отличаются тем, \emph{какую} систему держат
неподвижной:
\begin{itemize}
\item \vazhno{прямая полюсная фигура}: неподвижны оси образца; для каждого
зерна отмечают, куда смотрят его направления $\nfam{100}$ (или нормали к
другим плоскостям $\plfam{hkl}$)~— рис.~\ref{fig:texture-frames},\,\textit{г};
\item \vazhno{обратная полюсная фигура}: неподвижны оси кристалла; для
каждого зерна отмечают, куда смотрит выбранная ось образца, например
НН,~— рис.~\ref{fig:texture-frames},\,\textit{д};
\item \vazhno{функция распределения ориентаций}: отмечают сам поворот
(три угла Эйлера) каждого зерна~— точку в трёхмерном пространстве
ориентаций.
\end{itemize}

\begin{masshtab}[текстура]{-9.7/-9.2/{оси кристалла задаёт решётка}, -5.2/-3.8/{зерно}, -3.3/-1.3/{образец и пятно рентгена}}
Оси кристалла задаёт решётка (ячейка~— доли нанометра), но ориентация
одинакова во всём зерне размером 10--100~мкм. Оси образца привязаны к
изделию (лист, проволока~— сантиметры и больше). Полюсные фигуры
измеряют рентгеном или нейтронами по пятну размером в миллиметры: в него
попадают $10^4$--$10^6$ зёрен, поэтому полюсная фигура~— статистика
ориентаций множества зёрен, а не картинка одного кристалла.
\end{masshtab}

\subsubsection{Прямые полюсные фигуры}
\begin{opredelenie}
\term{Прямая полюсная фигура} (ППФ) $(hkl)$~— стереографическая проекция
в \emph{системе координат образца} распределения нормалей к плоскостям
$\plfam{hkl}$ всех кристаллитов.
\end{opredelenie}

Пусть ось текстуры $F\parallel[123]$. Нормали к плоскостям $\plfam{100}$
образуют с $F$ углы $\alpha_1$, $\alpha_2$, $\alpha_3$:
\[
\cos\alpha_1=\tfrac1{\sqrt{14}},\quad \cos\alpha_2=\tfrac2{\sqrt{14}},\quad
\cos\alpha_3=\tfrac3{\sqrt{14}}\qquad(\alpha_1=74{,}5^\circ,\ \alpha_2=57{,}7^\circ,\ \alpha_3=36{,}7^\circ).
\]
При повороте кристаллита вокруг $F$ каждая нормаль описывает конус.
Поэтому если ось $F$ находится в центре проекции, ППФ $(100)$~— три
концентрические окружности, а если на основном круге~— три пары
«параллелей» сетки Вульфа (рис.~\ref{fig:pf123}).

\begin{figure}[ht]
\centering
\begin{tikzpicture}[scale=1.55]
  % создано tools/gen_plastic.py: ППФ (100) аксиальной текстуры F || [123] (r = 1)
\begin{scope}
\draw[atomLine, line width=0.7pt] (0,0) circle (1);
\draw[fmBlue, line width=1.1pt] (0,0) circle (0.7604);
\draw[fmRed, line width=1.1pt] (0,0) circle (0.5508);
\draw[fmGreen!70!black, line width=1.1pt] (0,0) circle (0.3317);
\fill[fmBlue] (0.5201,0.5547) circle (0.045);
\fill[fmRed] (-0.4616,0.3004) circle (0.045);
\fill[fmGreen!70!black] (0.1402,-0.3006) circle (0.045);
\fill (0,0) circle (0.03);
\end{scope}
\begin{scope}[shift={(2.7,0)}]
\draw[atomLine, line width=0.7pt] (0,0) circle (1);
\draw[fmBlue, line width=1.1pt] (0.9636,0.2673) -- (0.9317,0.2586) -- (0.9007,0.2504) -- (0.8706,0.2428) -- (0.8414,0.2357) -- (0.8130,0.2290) -- (0.7852,0.2227) -- (0.7582,0.2167) -- (0.7319,0.2112) -- (0.7062,0.2059) -- (0.6811,0.2010) -- (0.6565,0.1964) -- (0.6324,0.1920) -- (0.6089,0.1879) -- (0.5858,0.1840) -- (0.5632,0.1804) -- (0.5410,0.1769) -- (0.5191,0.1737) -- (0.4977,0.1706) -- (0.4766,0.1677) -- (0.4558,0.1650) -- (0.4354,0.1625) -- (0.4152,0.1601) -- (0.3953,0.1578) -- (0.3757,0.1557) -- (0.3564,0.1538) -- (0.3372,0.1519) -- (0.3183,0.1502) -- (0.2995,0.1486) -- (0.2810,0.1471) -- (0.2626,0.1457) -- (0.2444,0.1444) -- (0.2264,0.1432) -- (0.2084,0.1421) -- (0.1906,0.1411) -- (0.1730,0.1403) -- (0.1554,0.1395) -- (0.1379,0.1387) -- (0.1205,0.1381) -- (0.1031,0.1376) -- (0.0859,0.1371) -- (0.0686,0.1368) -- (0.0514,0.1365) -- (0.0343,0.1363) -- (0.0171,0.1361) -- (0.0000,0.1361) -- (-0.0171,0.1361) -- (-0.0343,0.1363) -- (-0.0514,0.1365) -- (-0.0686,0.1368) -- (-0.0859,0.1371) -- (-0.1031,0.1376) -- (-0.1205,0.1381) -- (-0.1379,0.1387) -- (-0.1554,0.1395) -- (-0.1730,0.1403) -- (-0.1906,0.1411) -- (-0.2084,0.1421) -- (-0.2264,0.1432) -- (-0.2444,0.1444) -- (-0.2626,0.1457) -- (-0.2810,0.1471) -- (-0.2995,0.1486) -- (-0.3183,0.1502) -- (-0.3372,0.1519) -- (-0.3564,0.1538) -- (-0.3757,0.1557) -- (-0.3953,0.1578) -- (-0.4152,0.1601) -- (-0.4354,0.1625) -- (-0.4558,0.1650) -- (-0.4766,0.1677) -- (-0.4977,0.1706) -- (-0.5191,0.1737) -- (-0.5410,0.1769) -- (-0.5632,0.1804) -- (-0.5858,0.1840) -- (-0.6089,0.1879) -- (-0.6324,0.1920) -- (-0.6565,0.1964) -- (-0.6811,0.2010) -- (-0.7062,0.2059) -- (-0.7319,0.2112) -- (-0.7582,0.2167) -- (-0.7852,0.2227) -- (-0.8130,0.2290) -- (-0.8414,0.2357) -- (-0.8706,0.2428) -- (-0.9007,0.2504) -- (-0.9317,0.2586) -- (-0.9636,0.2673);
\draw[fmBlue, line width=1.1pt] (0.9636,-0.2673) -- (0.9317,-0.2586) -- (0.9007,-0.2504) -- (0.8706,-0.2428) -- (0.8414,-0.2357) -- (0.8130,-0.2290) -- (0.7852,-0.2227) -- (0.7582,-0.2167) -- (0.7319,-0.2112) -- (0.7062,-0.2059) -- (0.6811,-0.2010) -- (0.6565,-0.1964) -- (0.6324,-0.1920) -- (0.6089,-0.1879) -- (0.5858,-0.1840) -- (0.5632,-0.1804) -- (0.5410,-0.1769) -- (0.5191,-0.1737) -- (0.4977,-0.1706) -- (0.4766,-0.1677) -- (0.4558,-0.1650) -- (0.4354,-0.1625) -- (0.4152,-0.1601) -- (0.3953,-0.1578) -- (0.3757,-0.1557) -- (0.3564,-0.1538) -- (0.3372,-0.1519) -- (0.3183,-0.1502) -- (0.2995,-0.1486) -- (0.2810,-0.1471) -- (0.2626,-0.1457) -- (0.2444,-0.1444) -- (0.2264,-0.1432) -- (0.2084,-0.1421) -- (0.1906,-0.1411) -- (0.1730,-0.1403) -- (0.1554,-0.1395) -- (0.1379,-0.1387) -- (0.1205,-0.1381) -- (0.1031,-0.1376) -- (0.0859,-0.1371) -- (0.0686,-0.1368) -- (0.0514,-0.1365) -- (0.0343,-0.1363) -- (0.0171,-0.1361) -- (0.0000,-0.1361) -- (-0.0171,-0.1361) -- (-0.0343,-0.1363) -- (-0.0514,-0.1365) -- (-0.0686,-0.1368) -- (-0.0859,-0.1371) -- (-0.1031,-0.1376) -- (-0.1205,-0.1381) -- (-0.1379,-0.1387) -- (-0.1554,-0.1395) -- (-0.1730,-0.1403) -- (-0.1906,-0.1411) -- (-0.2084,-0.1421) -- (-0.2264,-0.1432) -- (-0.2444,-0.1444) -- (-0.2626,-0.1457) -- (-0.2810,-0.1471) -- (-0.2995,-0.1486) -- (-0.3183,-0.1502) -- (-0.3372,-0.1519) -- (-0.3564,-0.1538) -- (-0.3757,-0.1557) -- (-0.3953,-0.1578) -- (-0.4152,-0.1601) -- (-0.4354,-0.1625) -- (-0.4558,-0.1650) -- (-0.4766,-0.1677) -- (-0.4977,-0.1706) -- (-0.5191,-0.1737) -- (-0.5410,-0.1769) -- (-0.5632,-0.1804) -- (-0.5858,-0.1840) -- (-0.6089,-0.1879) -- (-0.6324,-0.1920) -- (-0.6565,-0.1964) -- (-0.6811,-0.2010) -- (-0.7062,-0.2059) -- (-0.7319,-0.2112) -- (-0.7582,-0.2167) -- (-0.7852,-0.2227) -- (-0.8130,-0.2290) -- (-0.8414,-0.2357) -- (-0.8706,-0.2428) -- (-0.9007,-0.2504) -- (-0.9317,-0.2586) -- (-0.9636,-0.2673);
\draw[fmRed, line width=1.1pt] (0.8452,0.5345) -- (0.8204,0.5192) -- (0.7962,0.5048) -- (0.7723,0.4911) -- (0.7488,0.4783) -- (0.7258,0.4661) -- (0.7031,0.4546) -- (0.6808,0.4438) -- (0.6589,0.4335) -- (0.6373,0.4238) -- (0.6161,0.4147) -- (0.5952,0.4060) -- (0.5746,0.3978) -- (0.5543,0.3900) -- (0.5342,0.3827) -- (0.5145,0.3757) -- (0.4950,0.3692) -- (0.4758,0.3630) -- (0.4568,0.3571) -- (0.4381,0.3516) -- (0.4195,0.3464) -- (0.4012,0.3414) -- (0.3831,0.3368) -- (0.3651,0.3324) -- (0.3474,0.3283) -- (0.3298,0.3245) -- (0.3123,0.3208) -- (0.2950,0.3175) -- (0.2779,0.3143) -- (0.2609,0.3114) -- (0.2440,0.3086) -- (0.2272,0.3061) -- (0.2106,0.3038) -- (0.1940,0.3016) -- (0.1775,0.2997) -- (0.1611,0.2979) -- (0.1448,0.2963) -- (0.1285,0.2949) -- (0.1123,0.2937) -- (0.0962,0.2926) -- (0.0801,0.2917) -- (0.0640,0.2910) -- (0.0480,0.2904) -- (0.0320,0.2900) -- (0.0160,0.2898) -- (0.0000,0.2897) -- (-0.0160,0.2898) -- (-0.0320,0.2900) -- (-0.0480,0.2904) -- (-0.0640,0.2910) -- (-0.0801,0.2917) -- (-0.0962,0.2926) -- (-0.1123,0.2937) -- (-0.1285,0.2949) -- (-0.1448,0.2963) -- (-0.1611,0.2979) -- (-0.1775,0.2997) -- (-0.1940,0.3016) -- (-0.2106,0.3038) -- (-0.2272,0.3061) -- (-0.2440,0.3086) -- (-0.2609,0.3114) -- (-0.2779,0.3143) -- (-0.2950,0.3175) -- (-0.3123,0.3208) -- (-0.3298,0.3245) -- (-0.3474,0.3283) -- (-0.3651,0.3324) -- (-0.3831,0.3368) -- (-0.4012,0.3414) -- (-0.4195,0.3464) -- (-0.4381,0.3516) -- (-0.4568,0.3571) -- (-0.4758,0.3630) -- (-0.4950,0.3692) -- (-0.5145,0.3757) -- (-0.5342,0.3827) -- (-0.5543,0.3900) -- (-0.5746,0.3978) -- (-0.5952,0.4060) -- (-0.6161,0.4147) -- (-0.6373,0.4238) -- (-0.6589,0.4335) -- (-0.6808,0.4438) -- (-0.7031,0.4546) -- (-0.7258,0.4661) -- (-0.7488,0.4783) -- (-0.7723,0.4911) -- (-0.7962,0.5048) -- (-0.8204,0.5192) -- (-0.8452,0.5345);
\draw[fmRed, line width=1.1pt] (0.8452,-0.5345) -- (0.8204,-0.5192) -- (0.7962,-0.5048) -- (0.7723,-0.4911) -- (0.7488,-0.4783) -- (0.7258,-0.4661) -- (0.7031,-0.4546) -- (0.6808,-0.4438) -- (0.6589,-0.4335) -- (0.6373,-0.4238) -- (0.6161,-0.4147) -- (0.5952,-0.4060) -- (0.5746,-0.3978) -- (0.5543,-0.3900) -- (0.5342,-0.3827) -- (0.5145,-0.3757) -- (0.4950,-0.3692) -- (0.4758,-0.3630) -- (0.4568,-0.3571) -- (0.4381,-0.3516) -- (0.4195,-0.3464) -- (0.4012,-0.3414) -- (0.3831,-0.3368) -- (0.3651,-0.3324) -- (0.3474,-0.3283) -- (0.3298,-0.3245) -- (0.3123,-0.3208) -- (0.2950,-0.3175) -- (0.2779,-0.3143) -- (0.2609,-0.3114) -- (0.2440,-0.3086) -- (0.2272,-0.3061) -- (0.2106,-0.3038) -- (0.1940,-0.3016) -- (0.1775,-0.2997) -- (0.1611,-0.2979) -- (0.1448,-0.2963) -- (0.1285,-0.2949) -- (0.1123,-0.2937) -- (0.0962,-0.2926) -- (0.0801,-0.2917) -- (0.0640,-0.2910) -- (0.0480,-0.2904) -- (0.0320,-0.2900) -- (0.0160,-0.2898) -- (0.0000,-0.2897) -- (-0.0160,-0.2898) -- (-0.0320,-0.2900) -- (-0.0480,-0.2904) -- (-0.0640,-0.2910) -- (-0.0801,-0.2917) -- (-0.0962,-0.2926) -- (-0.1123,-0.2937) -- (-0.1285,-0.2949) -- (-0.1448,-0.2963) -- (-0.1611,-0.2979) -- (-0.1775,-0.2997) -- (-0.1940,-0.3016) -- (-0.2106,-0.3038) -- (-0.2272,-0.3061) -- (-0.2440,-0.3086) -- (-0.2609,-0.3114) -- (-0.2779,-0.3143) -- (-0.2950,-0.3175) -- (-0.3123,-0.3208) -- (-0.3298,-0.3245) -- (-0.3474,-0.3283) -- (-0.3651,-0.3324) -- (-0.3831,-0.3368) -- (-0.4012,-0.3414) -- (-0.4195,-0.3464) -- (-0.4381,-0.3516) -- (-0.4568,-0.3571) -- (-0.4758,-0.3630) -- (-0.4950,-0.3692) -- (-0.5145,-0.3757) -- (-0.5342,-0.3827) -- (-0.5543,-0.3900) -- (-0.5746,-0.3978) -- (-0.5952,-0.4060) -- (-0.6161,-0.4147) -- (-0.6373,-0.4238) -- (-0.6589,-0.4335) -- (-0.6808,-0.4438) -- (-0.7031,-0.4546) -- (-0.7258,-0.4661) -- (-0.7488,-0.4783) -- (-0.7723,-0.4911) -- (-0.7962,-0.5048) -- (-0.8204,-0.5192) -- (-0.8452,-0.5345);
\draw[fmGreen!70!black, line width=1.1pt] (0.5976,0.8018) -- (0.5850,0.7854) -- (0.5723,0.7697) -- (0.5594,0.7546) -- (0.5464,0.7402) -- (0.5332,0.7264) -- (0.5200,0.7132) -- (0.5066,0.7005) -- (0.4932,0.6884) -- (0.4798,0.6768) -- (0.4663,0.6657) -- (0.4527,0.6551) -- (0.4392,0.6450) -- (0.4256,0.6353) -- (0.4121,0.6261) -- (0.3985,0.6173) -- (0.3849,0.6089) -- (0.3713,0.6010) -- (0.3578,0.5934) -- (0.3443,0.5861) -- (0.3307,0.5793) -- (0.3173,0.5728) -- (0.3038,0.5666) -- (0.2903,0.5607) -- (0.2769,0.5552) -- (0.2635,0.5500) -- (0.2501,0.5451) -- (0.2368,0.5405) -- (0.2235,0.5362) -- (0.2102,0.5321) -- (0.1969,0.5283) -- (0.1837,0.5248) -- (0.1704,0.5216) -- (0.1572,0.5186) -- (0.1441,0.5159) -- (0.1309,0.5134) -- (0.1177,0.5112) -- (0.1046,0.5092) -- (0.0915,0.5075) -- (0.0784,0.5060) -- (0.0653,0.5047) -- (0.0523,0.5037) -- (0.0392,0.5029) -- (0.0261,0.5023) -- (0.0131,0.5020) -- (0.0000,0.5019) -- (-0.0131,0.5020) -- (-0.0261,0.5023) -- (-0.0392,0.5029) -- (-0.0523,0.5037) -- (-0.0653,0.5047) -- (-0.0784,0.5060) -- (-0.0915,0.5075) -- (-0.1046,0.5092) -- (-0.1177,0.5112) -- (-0.1309,0.5134) -- (-0.1441,0.5159) -- (-0.1572,0.5186) -- (-0.1704,0.5216) -- (-0.1837,0.5248) -- (-0.1969,0.5283) -- (-0.2102,0.5321) -- (-0.2235,0.5362) -- (-0.2368,0.5405) -- (-0.2501,0.5451) -- (-0.2635,0.5500) -- (-0.2769,0.5552) -- (-0.2903,0.5607) -- (-0.3038,0.5666) -- (-0.3173,0.5728) -- (-0.3307,0.5793) -- (-0.3443,0.5861) -- (-0.3578,0.5934) -- (-0.3713,0.6010) -- (-0.3849,0.6089) -- (-0.3985,0.6173) -- (-0.4121,0.6261) -- (-0.4256,0.6353) -- (-0.4392,0.6450) -- (-0.4527,0.6551) -- (-0.4663,0.6657) -- (-0.4798,0.6768) -- (-0.4932,0.6884) -- (-0.5066,0.7005) -- (-0.5200,0.7132) -- (-0.5332,0.7264) -- (-0.5464,0.7402) -- (-0.5594,0.7546) -- (-0.5723,0.7697) -- (-0.5850,0.7854) -- (-0.5976,0.8018);
\draw[fmGreen!70!black, line width=1.1pt] (0.5976,-0.8018) -- (0.5850,-0.7854) -- (0.5723,-0.7697) -- (0.5594,-0.7546) -- (0.5464,-0.7402) -- (0.5332,-0.7264) -- (0.5200,-0.7132) -- (0.5066,-0.7005) -- (0.4932,-0.6884) -- (0.4798,-0.6768) -- (0.4663,-0.6657) -- (0.4527,-0.6551) -- (0.4392,-0.6450) -- (0.4256,-0.6353) -- (0.4121,-0.6261) -- (0.3985,-0.6173) -- (0.3849,-0.6089) -- (0.3713,-0.6010) -- (0.3578,-0.5934) -- (0.3443,-0.5861) -- (0.3307,-0.5793) -- (0.3173,-0.5728) -- (0.3038,-0.5666) -- (0.2903,-0.5607) -- (0.2769,-0.5552) -- (0.2635,-0.5500) -- (0.2501,-0.5451) -- (0.2368,-0.5405) -- (0.2235,-0.5362) -- (0.2102,-0.5321) -- (0.1969,-0.5283) -- (0.1837,-0.5248) -- (0.1704,-0.5216) -- (0.1572,-0.5186) -- (0.1441,-0.5159) -- (0.1309,-0.5134) -- (0.1177,-0.5112) -- (0.1046,-0.5092) -- (0.0915,-0.5075) -- (0.0784,-0.5060) -- (0.0653,-0.5047) -- (0.0523,-0.5037) -- (0.0392,-0.5029) -- (0.0261,-0.5023) -- (0.0131,-0.5020) -- (0.0000,-0.5019) -- (-0.0131,-0.5020) -- (-0.0261,-0.5023) -- (-0.0392,-0.5029) -- (-0.0523,-0.5037) -- (-0.0653,-0.5047) -- (-0.0784,-0.5060) -- (-0.0915,-0.5075) -- (-0.1046,-0.5092) -- (-0.1177,-0.5112) -- (-0.1309,-0.5134) -- (-0.1441,-0.5159) -- (-0.1572,-0.5186) -- (-0.1704,-0.5216) -- (-0.1837,-0.5248) -- (-0.1969,-0.5283) -- (-0.2102,-0.5321) -- (-0.2235,-0.5362) -- (-0.2368,-0.5405) -- (-0.2501,-0.5451) -- (-0.2635,-0.5500) -- (-0.2769,-0.5552) -- (-0.2903,-0.5607) -- (-0.3038,-0.5666) -- (-0.3173,-0.5728) -- (-0.3307,-0.5793) -- (-0.3443,-0.5861) -- (-0.3578,-0.5934) -- (-0.3713,-0.6010) -- (-0.3849,-0.6089) -- (-0.3985,-0.6173) -- (-0.4121,-0.6261) -- (-0.4256,-0.6353) -- (-0.4392,-0.6450) -- (-0.4527,-0.6551) -- (-0.4663,-0.6657) -- (-0.4798,-0.6768) -- (-0.4932,-0.6884) -- (-0.5066,-0.7005) -- (-0.5200,-0.7132) -- (-0.5332,-0.7264) -- (-0.5464,-0.7402) -- (-0.5594,-0.7546) -- (-0.5723,-0.7697) -- (-0.5850,-0.7854) -- (-0.5976,-0.8018);
\end{scope}

  \node[font=\small, anchor=north west] at (0.05,-0.02) {$F$};
  \fill (2.7,1) circle (0.035);
  \node[font=\small, above] at (2.7,1.02) {$F$};
  \fill (2.7,-1) circle (0.035);
  \node[font=\small] at (0,-1.25) {\textit{а}};
  \node[font=\small] at (2.7,-1.25) {\textit{б}};
  \begin{scope}[shift={(4.25,0.45)}, font=\footnotesize]
    \draw[fmGreen!70!black, line width=1.1pt] (0,0.4) -- (0.3,0.4) node[right, black]{$\alpha_3=36{,}7^\circ$};
    \draw[fmRed, line width=1.1pt] (0,0) -- (0.3,0) node[right, black]{$\alpha_2=57{,}7^\circ$};
    \draw[fmBlue, line width=1.1pt] (0,-0.4) -- (0.3,-0.4) node[right, black]{$\alpha_1=74{,}5^\circ$};
  \end{scope}
\end{tikzpicture}
\caption{Прямая полюсная фигура $(100)$ аксиальной текстуры
$F\parallel[123]$: \textit{а}~— ось $F$ в центре проекции (точки~—
полюса $\plfam{100}$ одного кристаллита); \textit{б}~— ось $F$ на основном
круге проекции}
\label{fig:pf123}
\end{figure}

\subsubsection{Обратные полюсные фигуры}
\begin{opredelenie}
\term{Обратная полюсная фигура} (ОПФ)~— стереографическая проекция в
\emph{кристаллографической системе координат} распределения выбранного
внешнего направления (оси образца).
\end{opredelenie}

Благодаря симметрии кристалла ОПФ кубического материала достаточно
строить в одном стандартном треугольнике $[100]$--$[110]$--$[111]$.
Для аксиальной текстуры $F\parallel\nfam{123}$ ОПФ направления $F$~—
\emph{одна точка} $[321]$ (тот из 24 эквивалентных векторов
$\nfam{123}$, который попадает в треугольник). Для текстуры
$F\parallel\nfam{110}$ ОПФ направления $G\perp F$~— множество направлений,
перпендикулярных $[110]$ (и эквивалентных им); в стандартном треугольнике
оно занимает две его стороны: $[110]$--$[111]$ и $[111]$--$[100]$
(рис.~\ref{fig:ipf}).

\begin{figure}[ht]
\centering
\begin{tikzpicture}[scale=9]
  % создано tools/gen_plastic.py: стандартный треугольник [100]-[110]-[111] (r = 1)
\fill[fmAmber!15] (0.0000,0.0000) -- (0.0101,0.0000) -- (0.0201,0.0000) -- (0.0302,0.0000) -- (0.0403,0.0000) -- (0.0504,0.0000) -- (0.0605,0.0000) -- (0.0706,0.0000) -- (0.0807,0.0000) -- (0.0909,0.0000) -- (0.1010,0.0000) -- (0.1112,0.0000) -- (0.1214,0.0000) -- (0.1317,0.0000) -- (0.1419,0.0000) -- (0.1522,0.0000) -- (0.1625,0.0000) -- (0.1729,0.0000) -- (0.1833,0.0000) -- (0.1937,0.0000) -- (0.2042,0.0000) -- (0.2147,0.0000) -- (0.2252,0.0000) -- (0.2358,0.0000) -- (0.2465,0.0000) -- (0.2572,0.0000) -- (0.2679,0.0000) -- (0.2788,0.0000) -- (0.2897,0.0000) -- (0.3006,0.0000) -- (0.3116,0.0000) -- (0.3227,0.0000) -- (0.3338,0.0000) -- (0.3451,0.0000) -- (0.3564,0.0000) -- (0.3678,0.0000) -- (0.3792,0.0000) -- (0.3908,0.0000) -- (0.4025,0.0000) -- (0.4142,0.0000) -- (0.4142,0.0000) -- (0.4142,0.0092) -- (0.4141,0.0185) -- (0.4139,0.0277) -- (0.4137,0.0370) -- (0.4135,0.0462) -- (0.4131,0.0555) -- (0.4127,0.0647) -- (0.4123,0.0740) -- (0.4118,0.0833) -- (0.4112,0.0925) -- (0.4105,0.1018) -- (0.4098,0.1111) -- (0.4091,0.1204) -- (0.4083,0.1297) -- (0.4074,0.1390) -- (0.4064,0.1483) -- (0.4054,0.1576) -- (0.4043,0.1669) -- (0.4032,0.1763) -- (0.4020,0.1856) -- (0.4007,0.1950) -- (0.3994,0.2044) -- (0.3980,0.2138) -- (0.3965,0.2232) -- (0.3950,0.2326) -- (0.3934,0.2420) -- (0.3917,0.2514) -- (0.3899,0.2609) -- (0.3881,0.2704) -- (0.3862,0.2798) -- (0.3843,0.2893) -- (0.3823,0.2989) -- (0.3802,0.3084) -- (0.3780,0.3180) -- (0.3758,0.3275) -- (0.3734,0.3371) -- (0.3710,0.3467) -- (0.3686,0.3564) -- (0.3660,0.3660) -- (0.3660,0.3660) -- (0.3551,0.3551) -- (0.3443,0.3443) -- (0.3337,0.3337) -- (0.3232,0.3232) -- (0.3127,0.3127) -- (0.3024,0.3024) -- (0.2923,0.2923) -- (0.2822,0.2822) -- (0.2722,0.2722) -- (0.2623,0.2623) -- (0.2525,0.2525) -- (0.2427,0.2427) -- (0.2331,0.2331) -- (0.2235,0.2235) -- (0.2141,0.2141) -- (0.2046,0.2046) -- (0.1953,0.1953) -- (0.1860,0.1860) -- (0.1768,0.1768) -- (0.1676,0.1676) -- (0.1585,0.1585) -- (0.1494,0.1494) -- (0.1404,0.1404) -- (0.1314,0.1314) -- (0.1224,0.1224) -- (0.1135,0.1135) -- (0.1047,0.1047) -- (0.0958,0.0958) -- (0.0870,0.0870) -- (0.0783,0.0783) -- (0.0695,0.0695) -- (0.0608,0.0608) -- (0.0521,0.0521) -- (0.0434,0.0434) -- (0.0347,0.0347) -- (0.0260,0.0260) -- (0.0173,0.0173) -- (0.0087,0.0087) -- (0.0000,0.0000) -- cycle;
\draw[atomLine, line width=0.6pt] (0.0000,0.0000) -- (0.0101,0.0000) -- (0.0201,0.0000) -- (0.0302,0.0000) -- (0.0403,0.0000) -- (0.0504,0.0000) -- (0.0605,0.0000) -- (0.0706,0.0000) -- (0.0807,0.0000) -- (0.0909,0.0000) -- (0.1010,0.0000) -- (0.1112,0.0000) -- (0.1214,0.0000) -- (0.1317,0.0000) -- (0.1419,0.0000) -- (0.1522,0.0000) -- (0.1625,0.0000) -- (0.1729,0.0000) -- (0.1833,0.0000) -- (0.1937,0.0000) -- (0.2042,0.0000) -- (0.2147,0.0000) -- (0.2252,0.0000) -- (0.2358,0.0000) -- (0.2465,0.0000) -- (0.2572,0.0000) -- (0.2679,0.0000) -- (0.2788,0.0000) -- (0.2897,0.0000) -- (0.3006,0.0000) -- (0.3116,0.0000) -- (0.3227,0.0000) -- (0.3338,0.0000) -- (0.3451,0.0000) -- (0.3564,0.0000) -- (0.3678,0.0000) -- (0.3792,0.0000) -- (0.3908,0.0000) -- (0.4025,0.0000) -- (0.4142,0.0000) -- (0.4142,0.0000) -- (0.4142,0.0092) -- (0.4141,0.0185) -- (0.4139,0.0277) -- (0.4137,0.0370) -- (0.4135,0.0462) -- (0.4131,0.0555) -- (0.4127,0.0647) -- (0.4123,0.0740) -- (0.4118,0.0833) -- (0.4112,0.0925) -- (0.4105,0.1018) -- (0.4098,0.1111) -- (0.4091,0.1204) -- (0.4083,0.1297) -- (0.4074,0.1390) -- (0.4064,0.1483) -- (0.4054,0.1576) -- (0.4043,0.1669) -- (0.4032,0.1763) -- (0.4020,0.1856) -- (0.4007,0.1950) -- (0.3994,0.2044) -- (0.3980,0.2138) -- (0.3965,0.2232) -- (0.3950,0.2326) -- (0.3934,0.2420) -- (0.3917,0.2514) -- (0.3899,0.2609) -- (0.3881,0.2704) -- (0.3862,0.2798) -- (0.3843,0.2893) -- (0.3823,0.2989) -- (0.3802,0.3084) -- (0.3780,0.3180) -- (0.3758,0.3275) -- (0.3734,0.3371) -- (0.3710,0.3467) -- (0.3686,0.3564) -- (0.3660,0.3660) -- (0.3660,0.3660) -- (0.3551,0.3551) -- (0.3443,0.3443) -- (0.3337,0.3337) -- (0.3232,0.3232) -- (0.3127,0.3127) -- (0.3024,0.3024) -- (0.2923,0.2923) -- (0.2822,0.2822) -- (0.2722,0.2722) -- (0.2623,0.2623) -- (0.2525,0.2525) -- (0.2427,0.2427) -- (0.2331,0.2331) -- (0.2235,0.2235) -- (0.2141,0.2141) -- (0.2046,0.2046) -- (0.1953,0.1953) -- (0.1860,0.1860) -- (0.1768,0.1768) -- (0.1676,0.1676) -- (0.1585,0.1585) -- (0.1494,0.1494) -- (0.1404,0.1404) -- (0.1314,0.1314) -- (0.1224,0.1224) -- (0.1135,0.1135) -- (0.1047,0.1047) -- (0.0958,0.0958) -- (0.0870,0.0870) -- (0.0783,0.0783) -- (0.0695,0.0695) -- (0.0608,0.0608) -- (0.0521,0.0521) -- (0.0434,0.0434) -- (0.0347,0.0347) -- (0.0260,0.0260) -- (0.0173,0.0173) -- (0.0087,0.0087) -- (0.0000,0.0000) -- cycle;
\draw[fmBlue, line width=2.2pt, line cap=round] (0.4142,0.0000) -- (0.4142,0.0092) -- (0.4141,0.0185) -- (0.4139,0.0277) -- (0.4137,0.0370) -- (0.4135,0.0462) -- (0.4131,0.0555) -- (0.4127,0.0647) -- (0.4123,0.0740) -- (0.4118,0.0833) -- (0.4112,0.0925) -- (0.4105,0.1018) -- (0.4098,0.1111) -- (0.4091,0.1204) -- (0.4083,0.1297) -- (0.4074,0.1390) -- (0.4064,0.1483) -- (0.4054,0.1576) -- (0.4043,0.1669) -- (0.4032,0.1763) -- (0.4020,0.1856) -- (0.4007,0.1950) -- (0.3994,0.2044) -- (0.3980,0.2138) -- (0.3965,0.2232) -- (0.3950,0.2326) -- (0.3934,0.2420) -- (0.3917,0.2514) -- (0.3899,0.2609) -- (0.3881,0.2704) -- (0.3862,0.2798) -- (0.3843,0.2893) -- (0.3823,0.2989) -- (0.3802,0.3084) -- (0.3780,0.3180) -- (0.3758,0.3275) -- (0.3734,0.3371) -- (0.3710,0.3467) -- (0.3686,0.3564) -- (0.3660,0.3660) -- (0.3551,0.3551) -- (0.3443,0.3443) -- (0.3337,0.3337) -- (0.3232,0.3232) -- (0.3127,0.3127) -- (0.3024,0.3024) -- (0.2923,0.2923) -- (0.2822,0.2822) -- (0.2722,0.2722) -- (0.2623,0.2623) -- (0.2525,0.2525) -- (0.2427,0.2427) -- (0.2331,0.2331) -- (0.2235,0.2235) -- (0.2141,0.2141) -- (0.2046,0.2046) -- (0.1953,0.1953) -- (0.1860,0.1860) -- (0.1768,0.1768) -- (0.1676,0.1676) -- (0.1585,0.1585) -- (0.1494,0.1494) -- (0.1404,0.1404) -- (0.1314,0.1314) -- (0.1224,0.1224) -- (0.1135,0.1135) -- (0.1047,0.1047) -- (0.0958,0.0958) -- (0.0870,0.0870) -- (0.0783,0.0783) -- (0.0695,0.0695) -- (0.0608,0.0608) -- (0.0521,0.0521) -- (0.0434,0.0434) -- (0.0347,0.0347) -- (0.0260,0.0260) -- (0.0173,0.0173) -- (0.0087,0.0087) -- (0.0000,0.0000);
\coordinate (T100) at (0.0000,0.0000);
\coordinate (T110) at (0.4142,0.0000);
\coordinate (T111) at (0.3660,0.3660);
\coordinate (T321) at (0.2967,0.1483);
\coordinate (T210) at (0.2361,0.0000);
\coordinate (T211) at (0.2247,0.2247);
\coordinate (T221) at (0.4000,0.2000);
\coordinate (T310) at (0.1623,0.0000);
\coordinate (T311) at (0.1583,0.1583);

  \foreach \p/\a/\t in {T100/north east/{[100]}, T110/north west/{[110]}, T111/south/{[111]}}{
    \fill[atomLine] (\p) circle (0.006);
    \node[font=\small, anchor=\a] at (\p) {$\t$};}
  \fill[fmRed] (T321) circle (0.009) node[below=1pt, font=\small, fmRed]{$[321]$};
  \foreach \p in {T211, T210}{\fill[atomLine!60] (\p) circle (0.004);}
  \node[font=\scriptsize, atomLine!80, anchor=north] at (T210) {$[210]$};
  \node[font=\scriptsize, atomLine!80, anchor=south east] at (T211) {$[211]$};
  \node[font=\footnotesize, fmBlue, anchor=west, align=left] at (0.43,0.22)
    {$G\perp F$ для\\текстуры $\nfam{110}$};
\end{tikzpicture}
\caption{Обратные полюсные фигуры в стандартном треугольнике: ось
текстуры $F\parallel\nfam{123}$ изображается точкой $[321]$ (красная);
направления $G\perp F$ при текстуре $F\parallel\nfam{110}$~— сторонами
$[110]$--$[111]$--$[100]$ (синие)}
\label{fig:ipf}
\end{figure}

\subsubsection{Функция распределения ориентаций}
Полюсные фигуры~— двумерные «проекции» трёхмерного распределения
ориентаций, и по одной ППФ или ОПФ ориентацию кристаллитов однозначно не
восстановить. Полное описание даёт ФРО.

\begin{opredelenie}
\term{Функция распределения ориентаций} $f(g)=f(\varphi,\theta,\psi)$~—
плотность вероятности найти кристаллит с ориентацией $g$, заданной
углами Эйлера~\eqref{eq:euler}: объёмная доля кристаллитов с ориентациями
между $g$ и $g+\mathrm dg$ равна $\Delta V/V=f(g)\,\mathrm dg$, где
$\mathrm dg=\sin\theta\,\mathrm d\varphi\,\mathrm d\theta\,\mathrm d\psi/8\pi^2$.
\end{opredelenie}

Углы Эйлера меняются в пределах $0\le\varphi<2\pi$, $0\le\theta\le\pi$,
$0\le\psi<2\pi$, т.\,е. ФРО задана в \term{ориентационном пространстве}~—
«параллелепипеде» $2\pi\times\pi\times2\pi$. Его геометрический объём
$4\pi^3\approx12\pi^2$, но объём с учётом веса $\sin\theta$ равен
\[
\int_0^{2\pi}\!\!\mathrm d\varphi\int_0^\pi\!\sin\theta\,\mathrm d\theta\int_0^{2\pi}\!\!\mathrm d\psi=8\pi^2 .
\]
Вес $\sin\theta$ появляется потому, что ФРО задана на непрерывной группе
вращений, а равные интервалы углов Эйлера соответствуют неравным
«объёмам» ориентаций: при $\theta\to0$ углы $\varphi$ и $\psi$ описывают
почти один и тот же поворот.

\begin{figure}[ht]
\centering
\begin{tikzpicture}[cubeproj, scale=1.0]
  \def\Lx{4} \def\Ly{4} \def\Lz{2}
  \draw[atomLine!50, dashed] (0,\Ly,0) -- (0,0,0) (0,\Ly,0) -- (\Lx,\Ly,0) (0,\Ly,0) -- (0,\Ly,\Lz);
  \draw[atomLine] (0,0,0) -- (\Lx,0,0) -- (\Lx,\Ly,0) -- (\Lx,\Ly,\Lz) -- (0,\Ly,\Lz) -- (0,0,\Lz) -- cycle
    (0,0,\Lz) -- (\Lx,0,\Lz) -- (\Lx,\Ly,\Lz) (\Lx,0,0) -- (\Lx,0,\Lz);
  \draw[arr] (0,0,0) -- (4.6,0,0) node[right, font=\small]{$\varphi$};
  \draw[arr] (0,0,0) -- (0,0,2.6) node[above, font=\small]{$\theta$};
  \draw[arr, atomLine!60] (0,\Ly,0) -- (0,4.7,0) node[above right=-2pt, font=\small, black]{$\psi$};
  \node[font=\scriptsize, below] at (\Lx,0,0) {$2\pi$};
  \node[font=\scriptsize, left] at (0,0,\Lz) {$\pi$};
  \node[font=\scriptsize, right] at (\Lx,\Ly,0) {$2\pi$};
  % стержень: θ = 36,7°, ψ = 26,6°
  \def\th{0.408} \def\ps{0.295}
  \draw[fmGray, densely dotted] (0,\ps,0) -- (0,\ps,\th) (0,0,\th) -- (0,\ps,\th);
  \draw[fmRed, line width=2.2pt] (0,\ps,\th) -- (\Lx,\ps,\th);
  \node[font=\footnotesize, fmRed, anchor=south] at (2.4,\ps,\th) {$\theta,\psi=\text{const}$, $\varphi$~— любое};
\end{tikzpicture}
\caption{ФРО аксиальной текстуры $F\parallel[123]$ в пространстве углов
Эйлера~— стержень, параллельный оси $\varphi$ ($\cos\theta=3/\sqrt{14}$,
$\operatorname{tg}\psi=1/2$; ещё 23 симметрично эквивалентных стержня не показаны)}
\label{fig:odf}
\end{figure}

Для аксиальной текстуры удобно совместить ось $F$ с осью $Z$ образца.
Первые два поворота Эйлера ($\varphi$ вокруг $Z$, затем $\theta$) задают,
где окажется ось $[001]$ кристалла, а последний~— поворот кристалла
вокруг неё. По третьей строке матрицы~\eqref{eq:euler} ось $Z$ образца
имеет в осях кристалла направляющие косинусы
$(\sin\theta\sin\psi,\ \sin\theta\cos\psi,\ \cos\theta)$. Если
$Z\parallel[mnp]$, то
\begin{equation}\label{eq:fiber-euler}
\cos\theta=\frac{p}{\sqrt{m^2+n^2+p^2}},\qquad \operatorname{tg}\psi=\frac mn ,
\end{equation}
а угол $\varphi$ (поворот вокруг оси текстуры) любой. Значит, \vazhno{ФРО
аксиальной текстуры~— стержень}: $\varphi$ меняется, $\theta$ и $\psi$
постоянны (рис.~\ref{fig:odf}). Поскольку у кубического кристалла
24 поворотных элемента симметрии, систему координат кристалла можно
выбрать 24 способами, и в полном ориентационном пространстве таких
стержней 24 (симметрично эквивалентных).

\begin{itogi}
\begin{itemize}
\item Пластическая деформация идёт скольжением (сдвиг на целое число
периодов, решётка не меняется) или двойникованием (сдвиг на доли периода,
решётка зеркально отражается).
\item Закон Шмида: скольжение начинается при
$\tau=\sigma\cos\chi\cos\lambda=\tau_\text{кр}$; фактор Шмида $\le\frac12$.
\item Тензор пластической дисторсии $u^0_{ij}=\alpha m_ih_j$: след нулевой
(объём сохраняется), антисимметричная часть~— поворот решётки. Для
произвольного формоизменения нужны 5 независимых систем.
\item ГЦК: 12 систем $\plfam{111}\nfam{110}$; ОЦК: направление $\nfam{111}$,
плоскости $\plfam{110}$, $\plfam{112}$, $\plfam{123}$; ГПУ: направление
$\nfam{11\br20}$, плоскость зависит от $c/a$.
\item Элементы двойникования $K_1,K_2,\eta_1,\eta_2,s$; в ГЦК и ОЦК
$s=1/\sqrt2$.
\item Текстура: аксиальная или ограниченная; описание~— ППФ (нормали в
осях образца), ОПФ (оси образца в осях кристалла), ФРО в пространстве
углов Эйлера объёмом $8\pi^2$.
\end{itemize}
\end{itogi}

\begin{proverka}
\item Почему при скольжении решётка не меняется, а при двойниковании
меняется?
\item При какой ориентации оси растяжения фактор Шмида равен
$\frac12$? Может ли он быть больше?
\item Найдите фактор Шмида системы $(\br111)[101]$ ГЦК-кристалла при
растяжении вдоль $[001]$. Сколько систем имеют такой же фактор?
\item Почему ГПУ-металлы с одним базисным скольжением хрупки в
поликристаллическом состоянии?
\item Вычислите угол $2\phi$ между плоскостями $K_1$ и $K_2$ двойникования
ОЦК-кристалла и проверьте формулу~\eqref{eq:twin-shear}.
\item Какой будет ППФ $(100)$ для аксиальной текстуры $F\parallel[001]$?
\end{proverka}
